% Examen blanc — Algèbre linéaire et déterminants (Le Collé, PT)
% Compiler tel quel : le sujet. Pour le corrigé : pdflatex -jobname=eb-al-det-corrige "\def\CORRIGE{}% Examen blanc — Algèbre linéaire et déterminants (Le Collé, PT)
% Compiler tel quel : le sujet. Pour le corrigé : pdflatex -jobname=eb-al-det-corrige "\def\CORRIGE{}% Examen blanc — Algèbre linéaire et déterminants (Le Collé, PT)
% Compiler tel quel : le sujet. Pour le corrigé : pdflatex -jobname=eb-al-det-corrige "\def\CORRIGE{}% Examen blanc — Algèbre linéaire et déterminants (Le Collé, PT)
% Compiler tel quel : le sujet. Pour le corrigé : pdflatex -jobname=eb-al-det-corrige "\def\CORRIGE{}\input{examen-al-determinants.tex}"
\documentclass[a4paper,11pt]{article}

% ============================================================
%  Le Collé — Algèbre linéaire (compléments de deuxième année)
%  © 2026 Ahmed Amine Ouqlani
% ============================================================

% ---------- Polices ----------
\usepackage[utf8]{inputenc}
\usepackage[T1]{fontenc}
\usepackage[sc]{mathpazo}          % texte et mathématiques : Palatino
\usepackage[scaled=0.88]{helvet}   % titres et repères : linéale
\linespread{1.07}
\usepackage{amsmath,amssymb,amsthm,mathrsfs,stmaryrd}
\usepackage{microtype}

% ---------- Mise en page ----------
\usepackage[a4paper,top=2.3cm,bottom=2.3cm,left=2.4cm,right=2.4cm,
            headheight=14pt,headsep=0.8cm,footskip=1.1cm]{geometry}
\setlength{\parindent}{0pt}
\setlength{\parskip}{0.45em plus 0.15em minus 0.1em}
\raggedbottom
\emergencystretch=1.5em
\clubpenalty=10000 \widowpenalty=10000 \displaywidowpenalty=10000 \brokenpenalty=10000

% ---------- Couleurs : un bleu de nuit, un ocre, des gris ----------
\usepackage{xcolor}
\definecolor{Ink}{HTML}{1C2533}
\definecolor{Bleu}{HTML}{1F3A5F}
\definecolor{BleuMoyen}{HTML}{4A6A8F}
\definecolor{BleuClair}{HTML}{9DB2C9}
\definecolor{Fond}{HTML}{F2F4F7}
\definecolor{Ocre}{HTML}{B07A26}
\definecolor{Brique}{HTML}{A23B2A}
\definecolor{Gris}{HTML}{5E6672}
\definecolor{Filet}{HTML}{C5CCD6}

% ---------- Packages ----------
\usepackage{enumitem}
\setlist{itemsep=0.25em,topsep=0.35em,parsep=0pt}
\setlist[itemize]{label=\textcolor{BleuMoyen}{\small\textbullet},leftmargin=1.4em}
\setlist[enumerate]{leftmargin=1.8em,label=\textsf{\small\color{Bleu}\arabic*.}}
\usepackage{booktabs,array,multicol}
\usepackage{needspace,float}
\usepackage[font=small,labelfont={sf,color=Bleu},labelsep=period,skip=6pt]{caption}
\usepackage{tikz}
\usetikzlibrary{arrows.meta,calc,decorations.markings}
\usepackage{pgfplots}
\pgfplotsset{compat=1.17}
\usepackage{tkz-tab}
\usepackage[most]{tcolorbox}
\usepackage{titlesec,titletoc}
\usepackage{fancyhdr,lastpage,etoolbox}
\usepackage[colorlinks=true,linkcolor=Bleu,urlcolor=Bleu,citecolor=Bleu,
            pdftitle={Examen blanc — Algèbre linéaire et déterminants},pdfauthor={Ahmed Amine Ouqlani}]{hyperref}

\numberwithin{equation}{section}
\numberwithin{figure}{section}
\numberwithin{table}{section}

% ---------- Titres ----------
% numéro du chapitre dans le parcours PT (affiché), distinct de la numérotation interne
\newcommand{\grandnum}{\thesection}\newcommand{\ongletcouleur}{Bleu}
% Chapitres (sections) : grand numéro, titre en linéale, filet ocre
\titleformat{\section}[display]
  {\normalfont\sffamily\color{Ink}}
  {{\small\color{Gris}\textls[200]{\MakeUppercase{Chapitre}}}\\[-0.2em]%
   {\fontsize{64}{64}\selectfont\color{BleuClair}\grandnum}}
  {0.4em}
  {\fontsize{25}{30}\selectfont\bfseries}
  [\vspace{0.5em}{\color{Ocre}\rule{3.2cm}{1.4pt}}]
\titlespacing*{\section}{0pt}{0pt}{1.6em}
\titleformat{name=\section,numberless}[display]
  {\normalfont\sffamily\color{Ink}}{}{0pt}{\fontsize{25}{30}\selectfont\bfseries}
  [\vspace{0.5em}{\color{Ocre}\rule{3.2cm}{1.4pt}}]
\titlespacing*{name=\section,numberless}{0pt}{0pt}{1.6em}

\titleformat{\subsection}
  {\needspace{5\baselineskip}\normalfont\sffamily\Large\bfseries\color{Ink}}
  {\color{Bleu}\thesubsection}{0.8em}{}
\titlespacing*{\subsection}{0pt}{1.8em}{0.5em}
\titleformat{\subsubsection}
  {\needspace{5\baselineskip}\normalfont\sffamily\large\bfseries\color{Ink}}
  {\color{Bleu}\thesubsubsection}{0.7em}{}
\titlespacing*{\subsubsection}{0pt}{1.6em}{0.4em}
\titleformat{\paragraph}[runin]{\normalfont\sffamily\bfseries\color{Ink}}{}{0pt}{}[.]
\titlespacing*{\paragraph}{0pt}{0.9em}{0.6em}
\setcounter{secnumdepth}{3}
\setcounter{tocdepth}{2}

% ---------- Table des matières ----------
\titlecontents{section}[0pt]{\addvspace{0.7em}\sffamily\bfseries\color{Ink}}
  {}{}{\hfill\contentspage}[\addvspace{0.2em}]
\titlecontents{subsection}[2.2em]{\small\linespread{1}\selectfont}
  {\makebox[2.6em][l]{\color{Gris}\thecontentslabel}}{}{\titlerule*[0.7pc]{.}\contentspage}

% ---------- En-têtes et pieds de page ----------
\pagestyle{fancy}
\fancyhf{}
\renewcommand{\sectionmark}[1]{\markboth{#1}{}}
\fancyhead[L]{\footnotesize\sffamily\color{Gris}\leftmark}
\fancyhead[R]{\footnotesize\sffamily\scshape\color{Gris}Le Collé}
\fancyfoot[C]{\small\sffamily\color{Ink}\thepage}
\fancyfoot[L]{\scriptsize\sffamily\hyperlink{toc}{\color{Filet}sommaire}}
\renewcommand{\headrulewidth}{0.4pt}
\renewcommand{\headrule}{{\color{Filet}\hrule height\headrulewidth width\headwidth}}
\renewcommand{\footrulewidth}{0pt}
\fancypagestyle{plain}{\fancyhf{}\fancyfoot[C]{\small\sffamily\color{Ink}\thepage}\renewcommand{\headrulewidth}{0pt}}

% ---------- Énoncés ----------
\newtheoremstyle{enonce}{0pt}{0pt}{}{}{\sffamily\bfseries\color{Bleu}}{.}{0.5em}
  {\thmname{#1}\thmnumber{~#2}\thmnote{\ \normalfont\sffamily\color{Ink}(#3)}}
\newtheoremstyle{exemple}{0.9em}{0.9em}{}{}{\sffamily\bfseries\color{BleuMoyen}}{.}{0.5em}
  {\thmname{#1}\thmnumber{~#2}\thmnote{\ \normalfont\sffamily\color{Ink}(#3)}}
\newtheoremstyle{remarque}{0.7em}{0.7em}{}{}{\sffamily\itshape}{.}{0.5em}{\thmname{#1}\thmnote{\ (#3)}}
\theoremstyle{enonce}
\newtheorem{definition}{Définition}[section]
\newtheorem{theoreme}[definition]{Théorème}
\newtheorem{proposition}[definition]{Proposition}
\newtheorem{corollaire}[definition]{Corollaire}
\theoremstyle{exemple}
\newtheorem{exemple}[definition]{Exemple}
\theoremstyle{remarque}
\newtheorem*{remarque}{Remarque}

\tcbset{
  defi/.style={enhanced,breakable,colback=white,colframe=white,boxrule=0pt,arc=0pt,
     borderline west={1.2pt}{0pt}{Bleu},left=11pt,right=4pt,top=3pt,bottom=3pt,
     before skip=1.1em,after skip=1.1em},
  thm/.style={enhanced,breakable,colback=Fond,colframe=Fond,boxrule=0pt,arc=0pt,
     borderline west={2.4pt}{0pt}{Bleu},left=11pt,right=8pt,top=6pt,bottom=6pt,
     before skip=1.2em,after skip=1.2em}}
\tcolorboxenvironment{definition}{defi}
\tcolorboxenvironment{theoreme}{thm}
\tcolorboxenvironment{proposition}{thm}
\tcolorboxenvironment{corollaire}{thm}

\renewcommand{\proofname}{Démonstration}
\renewcommand{\qedsymbol}{\textcolor{BleuClair}{$\blacksquare$}}
\makeatletter
\renewenvironment{proof}[1][\proofname]{\par\pushQED{\qed}\normalfont
  \topsep0.5em\relax\trivlist\item[\hskip\labelsep\sffamily\itshape\color{Gris}#1.]\ignorespaces}
  {\popQED\endtrivlist\@endpefalse}
\makeatother

% ---------- Méthode : filets haut et bas, titre en bandeau ----------
\makeatletter\newcommand{\methlabelfix}{\protected@edef\@currentlabel{\thetcbcounter}}\makeatother
\newtcolorbox[auto counter,number within=section]{methode}[2][]{%
  enhanced,unbreakable,colback=white,colframe=white,boxrule=0pt,arc=0pt,
  borderline north={0.9pt}{0pt}{Bleu},borderline south={0.9pt}{0pt}{Bleu},
  left=4pt,right=4pt,top=8pt,bottom=8pt,before skip=1.4em,after skip=1.4em,
  before upper={\methlabelfix\phantomsection\addcontentsline{met}{metentry}{\protect\makebox[2.6em][l]{\thetcbcounter}#2}%
     {\sffamily\bfseries\small\color{white}\colorbox{Bleu}{\strut\ MÉTHODE~\thetcbcounter\ }}%
     \ \ {\sffamily\bfseries\color{Bleu}#2}\par\smallskip},#1}

% ---------- Repères pédagogiques : un filet de couleur, un mot-clé ----------
\newcommand{\definirbloc}[3]{% nom, couleur, libellé
  \newtcolorbox{#1}[1][]{enhanced,breakable,blanker,
    borderline west={1.4pt}{0pt}{#2},left=10pt,right=2pt,top=1pt,bottom=1pt,
    before skip=1em,after skip=1em,
    before upper={{\sffamily\bfseries\footnotesize\color{#2}\textls[120]{\MakeUppercase{#3}}}\quad},##1}}
\definirbloc{reflexe}{Bleu}{Réflexe}
\definirbloc{astuce}{Ocre}{Astuce}
\definirbloc{attention}{Brique}{Attention}
\definirbloc{enpratique}{BleuMoyen}{En pratique}
\definirbloc{plusloin}{Gris}{Pour aller plus loin}
\definirbloc{strategie}{Bleu}{Stratégie}
\definirbloc{lecture}{BleuMoyen}{Lire le schéma}
% compatibilité avec les parties reprises
\newenvironment{garde}{\begin{attention}}{\end{attention}}
\newenvironment{complementHP}{\begin{plusloin}}{\end{plusloin}}
\newenvironment{interpretation}{\begin{enpratique}}{\end{enpratique}}

% ---------- Exercices, corrigés ----------
\newcounter{exo}[section]
\renewcommand{\theexo}{\thesection.\arabic{exo}}
\newcommand{\niveau}[1]{{\color{Ocre}\footnotesize
  \ifcase#1\or$\bullet\circ\circ\circ\circ$\or$\bullet\bullet\circ\circ\circ$\or$\bullet\bullet\bullet\circ\circ$%
  \or$\bullet\bullet\bullet\bullet\circ$\or$\bullet\bullet\bullet\bullet\bullet$\fi}}
\newcommand{\serie}[1]{\par\needspace{8\baselineskip}\vspace{1.5em}%
  \noindent{\sffamily\bfseries\color{Bleu}#1}\par\nobreak\vspace{-0.4em}%
  \noindent{\color{Filet}\rule{\linewidth}{0.4pt}}\par\nobreak\vspace{0.1em}}
\newenvironment{exo}[3]{%
  \par\needspace{5\baselineskip}\vspace{0.75em}\refstepcounter{exo}\label{e:#1}%
  \def\exolab{#1}%
  \noindent{\sffamily\bfseries\color{Bleu}Exercice~\theexo}\quad\niveau{#2}%
  \hfill{\footnotesize\sffamily\color{Gris}#3}\par\nobreak\vspace{0.15em}}%
  {\par\nopagebreak\noindent\hfill{\scriptsize\sffamily\hyperref[c:\exolab]{\color{Gris}corrigé p.~\pageref{c:\exolab}}}\par}
\newenvironment{corrige}[1]{%
  \par\needspace{6\baselineskip}\vspace{1em}\phantomsection\label{c:#1}%
  \noindent{\sffamily\bfseries\color{Bleu}Exercice~\ref{e:#1}}\quad
  {\scriptsize\sffamily\hyperref[e:#1]{\color{Gris}énoncé p.~\pageref{e:#1}}}\par\nobreak\vspace{0.2em}}%
  {\hfill\textcolor{BleuClair}{$\square$}\par}
\newcommand{\etape}[1]{\par\smallskip\noindent{\sffamily\bfseries\small\color{BleuMoyen}#1.}\ }
\newenvironment{solution}{\par\noindent{\sffamily\bfseries\small\color{BleuMoyen}Solution.}\ }{\hfill\textcolor{BleuClair}{$\square$}\par}
\newcommand{\chercher}{\par\smallskip\noindent{\color{Filet}\rule[0.5ex]{1.2cm}{0.4pt}}\
  {\footnotesize\sffamily\itshape\color{Gris}laisser chercher l'étudiant avant de poursuivre}\par\smallskip}

% ---------- Guide de séance ----------

\newcommand{\marge}[1]{\leavevmode\marginpar{\raggedleft\scriptsize\sffamily\color{Gris}\textls[100]{\MakeUppercase{#1}}}\ignorespaces}
\newcommand{\Etape}[4]{%
  \subsection[Étape #1 --- #3 {\color{Gris}(#2)}]{#3}\label{etape:#1}%
  \noindent{\small\sffamily\color{Gris}{\color{Ocre}\bfseries Étape #1}\quad·\quad #2\quad·\quad #4}\par\medskip}
\newcommand{\question}[1]{\par\smallskip\noindent{\sffamily\bfseries\small\color{BleuMoyen}Question.}\ \textit{#1}\par\smallskip}
\newcommand{\rappel}[1]{\par\smallskip\noindent\phantomsection\addcontentsline{rap}{rapentry}{#1}%
  {\sffamily\bfseries\scriptsize\colorbox{Fond}{\color{Gris}\strut RAPPEL}}\ {\sffamily\small\color{Gris}#1.}\ }
\newcommand{\pourquoi}{\par\noindent{\sffamily\bfseries\small\color{Ocre}Pourquoi ?}\ }

% Ouverture de chapitre : \ouverture{objectifs}{prérequis}{durée}
\newcommand{\ouverture}[3]{%
  \begin{tcolorbox}[enhanced,colback=Fond,colframe=Fond,arc=0pt,boxrule=0pt,left=12pt,right=12pt,top=10pt,bottom=10pt,
     before skip=1.2em,after skip=1.6em]
  \begin{minipage}[t]{0.56\linewidth}
  {\sffamily\bfseries\small\color{Bleu}OBJECTIFS}\par\vspace{0.3em}\small #1
  \end{minipage}\hfill
  \begin{minipage}[t]{0.38\linewidth}
  {\sffamily\bfseries\small\color{Bleu}PRÉREQUIS}\par\vspace{0.3em}\small #2
  \ifx\relax#3\relax\else\par\vspace{0.6em}{\sffamily\bfseries\small\color{Bleu}DURÉE}\par\vspace{0.2em}\small #3\fi
  \end{minipage}
  \end{tcolorbox}}
% Fin de chapitre
\newenvironment{essentiel}{\par\needspace{12\baselineskip}\vspace{1.5em}%
  \begin{tcolorbox}[enhanced,unbreakable,colback=white,colframe=white,boxrule=0pt,arc=0pt,
   borderline north={1.4pt}{0pt}{Ocre},left=2pt,right=2pt,top=8pt,bottom=4pt]
  {\sffamily\bfseries\color{Ocre}L'ESSENTIEL DU CHAPITRE}\par\medskip}{\end{tcolorbox}}

% ---------- Notations ----------
\newcommand{\K}{\mathbb{K}}
\newcommand{\R}{\mathbb{R}}
\newcommand{\C}{\mathbb{C}}
\newcommand{\N}{\mathbb{N}}
\newcommand{\Z}{\mathbb{Z}}
\newcommand{\Mn}[1]{\mathcal{M}_{n}(#1)}
\newcommand{\GLn}[1]{\mathrm{GL}_n(#1)}
\newcommand{\LE}{\mathcal{L}(E)}
\newcommand{\rg}{\operatorname{rg}}
\newcommand{\vect}{\operatorname{Vect}}
\newcommand{\Mat}{\operatorname{Mat}}
\newcommand{\Id}{\operatorname{Id}}
\newcommand{\sh}{\operatorname{sh}}
\newcommand{\ch}{\operatorname{ch}}
\renewcommand{\th}{\operatorname{th}}
\DeclareMathOperator{\argsh}{argsh}
\newcommand{\vv}[1]{\overrightarrow{#1}}
\newcommand{\ve}[1]{\vec{#1}}
\newcommand{\ient}[2]{\llbracket #1,#2\rrbracket}
\newcommand{\dt}[1]{\left|\begin{matrix}#1\end{matrix}\right|}
\newcommand{\dis}{\displaystyle}

% ---------- Figures ----------
\pgfplotsset{trig format plots=rad}
\pgfplotsset{courbe/.style={axis lines=middle,axis equal image,
  axis line style={Gris,-{Stealth[length=5pt]}},
  xlabel={$x$},ylabel={$y$},ticklabel style={font=\scriptsize\sffamily,color=Gris},
  every axis x label/.style={at={(ticklabel* cs:1)},anchor=west,font=\small},
  every axis y label/.style={at={(ticklabel* cs:1)},anchor=south,font=\small},
  samples=300,smooth,width=0.6\linewidth}}
\tikzset{arcetude/.style={Bleu,line width=1.25pt},
         reste/.style={BleuClair,line width=1pt},
         tgte/.style={Ocre,line width=0.9pt},
         asy/.style={Gris,dashed,line width=0.6pt},
         pt/.style={circle,fill=Ink,inner sep=1.2pt}}

% ---------- Navigation ----------
\makeatletter
\newcommand*\l@rapentry{\@dottedtocline{1}{0em}{0em}}
\newcommand*\l@metentry{\@dottedtocline{1}{0em}{0em}}
\newcommand{\listemethodes}{\@starttoc{met}}
\newcommand{\listerappels}{\@starttoc{rap}}
\makeatother
% onglets de chapitre sur le bord droit
\usepackage{eso-pic}
\newif\ifonglets \ongletsfalse
\newcommand{\dessineronglet}{%
  \edef\ongcol{\ongletcouleur}%
  \edef\ongnum{\grandnum}%
  \pgfmathsetmacro{\ongy}{3.2+2.3*(\ongnum-1)}%
  \begin{tikzpicture}[remember picture,overlay]
  \fill[\ongcol] ([xshift=-0.75cm,yshift=-\ongy cm]current page.north east)
     rectangle ([yshift=-\ongy cm-1.9cm]current page.north east);
  \node[text=white,font=\sffamily\bfseries\large] at ([xshift=-0.37cm,yshift=-\ongy cm-0.95cm]current page.north east)
     {\ongnum};
  \end{tikzpicture}}
\AddToShipoutPictureBG{\ifonglets\ifnum\value{section}>0 \dessineronglet\fi\fi}
% blocs de navigation et d'enseignement
\definirbloc{commencer}{Gris}{Si vous commencez ici}
\definirbloc{horsprogramme}{Brique}{Hors programme}
\newcommand{\HP}{\,{\sffamily\bfseries\scriptsize\textcolor{Brique}{[HP]}}}
\newcommand{\source}[1]{{\footnotesize\sffamily\color{Gris}\,(#1)}}
\newcommand{\coupdepouce}[1]{\par\vspace{-0.2em}\noindent{\footnotesize\sffamily\color{Gris}%
  \textbf{Coup de pouce} (à donner si l'étudiant bloque) : #1}\par}
\newcommand{\pointarret}[1]{\par\vspace{1em}\noindent{\color{Filet}\rule[0.5ex]{\linewidth}{0.4pt}}\par\nobreak
  \noindent{\footnotesize\sffamily\color{Ocre}\bfseries POINT D'ARRÊT POSSIBLE}\quad{\footnotesize\sffamily\color{Gris}#1}\par\vspace{0.6em}}
\newcommand{\demarche}[1]{\noindent{\sffamily\bfseries\small\color{BleuMoyen}Démarche.}\ \textit{#1}\par\smallskip}

% ---------- Corrigés : idées, astuces, pièges, à retenir ----------
\newcommand{\choix}[1]{\par\noindent{\sffamily\bfseries\small\color{BleuMoyen}Choix de la méthode.}\ #1\par\smallskip}
\newcommand{\gain}[1]{\par\noindent{\sffamily\bfseries\small\color{Ocre}Astuce.}\ #1\par}
\newcommand{\piege}[1]{\par\noindent{\sffamily\bfseries\small\color{Brique}Piège.}\ #1\par}
\newtcolorbox{retenirbox}{enhanced,breakable,colback=Fond,colframe=Fond,arc=0pt,boxrule=0pt,
  borderline west={2pt}{0pt}{Ocre},left=8pt,right=6pt,top=4pt,bottom=4pt,before skip=0.7em,after skip=0.5em}
\newcommand{\retenir}[1]{\begin{retenirbox}{\sffamily\bfseries\footnotesize\color{Ocre}À RETENIR}\quad #1\end{retenirbox}}
\newcommand{\TD}[1]{{\color{Ocre}\bfseries TD\,#1}\ ·\ }
\newcommand{\q}[1]{\par\smallskip\noindent{\sffamily\bfseries\color{Bleu}#1}\ }
\providecommand{\iddots}{\mathinner{\mkern1mu\raise1pt\hbox{.}\mkern2mu\raise4pt\hbox{.}\mkern2mu\raise7pt\hbox{.}\mkern1mu}}


% =====================================================================
%  Examen blanc — Algèbre linéaire et déterminants (Le Collé, PT)
%  Une source, deux PDF : \def\CORRIGE{} avant \input donne le corrigé.
% =====================================================================
\providecommand{\Ker}{\operatorname{Ker}}\renewcommand{\Im}{\operatorname{Im}}
\fancyfoot[L]{}\fancyhead[L]{\footnotesize\sffamily\color{Gris}Examen blanc · Algèbre linéaire et déterminants\ifdefined\CORRIGE\ · corrigé\fi}
\newcommand{\pts}[1]{\hfill{\footnotesize\sffamily\color{Gris}(#1)}}
\newcommand{\Ex}[3]{\par\needspace{6\baselineskip}\vspace{1.3em}\noindent
  {\sffamily\bfseries\large\color{Bleu}Exercice #1 \,·\, #2}\hfill{\sffamily\small\color{Ocre}#3}\par
  \vspace{2pt}{\color{Filet}\rule{\linewidth}{0.4pt}}\par\smallskip}
\newcommand{\Pb}[2]{\par\needspace{6\baselineskip}\vspace{1.3em}\noindent
  {\sffamily\bfseries\large\color{Bleu}Problème \,·\, #1}\hfill{\sffamily\small\color{Ocre}#2}\par
  \vspace{2pt}{\color{Filet}\rule{\linewidth}{0.4pt}}\par\smallskip}
\newcommand{\Partie}[1]{\par\needspace{4\baselineskip}\medskip\noindent{\sffamily\bfseries\color{BleuMoyen}#1}\par\nobreak\smallskip}
\newcommand{\Q}[1]{\par\needspace{3\baselineskip}\medskip\noindent{\sffamily\bfseries\color{Bleu}#1}\enspace}
\newcommand{\autre}[1]{\par\noindent{\sffamily\bfseries\small\color{BleuMoyen}Autre méthode.}\ #1\par}
\newcommand{\lienDL}[1]{\par\noindent{\sffamily\bfseries\small\color{Ocre}Lien avec le DL.}\ #1\par}
\setlist[enumerate,1]{label=\arabic*.,leftmargin=1.6em,itemsep=4pt,topsep=4pt}
\setlist[enumerate,2]{label=(\alph*),leftmargin=1.6em,itemsep=2pt,topsep=2pt}

\begin{document}
\pagestyle{fancy}
\thispagestyle{plain}

% ------------------------------------------------------------------ en-tête
\begin{tikzpicture}[remember picture,overlay]
\fill[Bleu] (current page.north west) rectangle ([yshift=-4.2cm]current page.north east);
\node[anchor=north west,text=white,font=\sffamily\footnotesize,inner sep=0] at ([xshift=2.2cm,yshift=-1.1cm]current page.north west)
  {\textls[120]{MATHÉMATIQUES · CPGE PT · EXAMEN BLANC}};
\node[anchor=north west,text=white,font=\sffamily\bfseries\LARGE,inner sep=0] at ([xshift=2.2cm,yshift=-1.75cm]current page.north west)
  {Algèbre linéaire et déterminants};
\node[anchor=north west,text=white,font=\sffamily\normalsize,inner sep=0] at ([xshift=2.2cm,yshift=-2.75cm]current page.north west)
  {\ifdefined\CORRIGE Corrigé détaillé et diagnostic\else Durée : 3 heures \quad·\quad Calculatrice interdite\fi};
\node[anchor=north east,text=white,font=\sffamily\scshape\large,inner sep=0] at ([xshift=-2.2cm,yshift=-1.05cm]current page.north east)
  {Le Collé};
\end{tikzpicture}
\vspace*{2.9cm}

\ifdefined\CORRIGE
% =====================================================================
%                              CORRIGÉ
% =====================================================================
\markboth{Corrigé · Examen blanc}{}
\paragraph{Comment utiliser ce corrigé} Ne l'ouvre qu'après avoir composé les 3 heures. Corrige d'abord ta copie avec un
stylo d'une autre couleur, question par question, en comptant les points. Pour chaque question ratée, lis le \emph{choix de la
méthode} : c'est lui qu'il faut retenir, plus que le calcul. La grille de la dernière page dit quoi revoir.

\subsection*{Ce que le sujet évalue}
\begin{center}
\small\renewcommand{\arraystretch}{1.22}
\begin{tabular}{@{}>{\raggedright\arraybackslash}p{7.6cm}>{\raggedright\arraybackslash}p{5.4cm}@{}}
\toprule
\textbf{Compétence} & \textbf{Questions}\\\midrule
Démontrer qu'une famille est libre (combinaison nulle, évaluation, composition) & 1.2 · 2.3(b) · 2.4(a) · P.3\\
Démontrer qu'une famille est génératrice, trouver une base d'un sous-espace & 1.3 · 2.2(a) · 3.2\\
Utiliser la dimension (« $n$ vecteurs en dimension $n$ ») & 1.1 · 2.2(b) · 2.4(a) · 3.3 · P.3\\
Coordonnées dans une base, changement de base & 2.1(b) · 2.3(c) · 2.4(c) · 3.4 · P.4\\
Sous-espaces, somme directe, supplémentaires, projecteur & 1.4 · 3.1 · 3.3 · 3.4\\
Application linéaire : noyau, image, rang, théorème du rang & 1.6 · P.2\\
Calculer un déterminant (opérations, développement, factorisation) & 2.1(a) · 4.1 · 4.2\\
Propriétés du déterminant ($\det(\lambda A)$, $\det A^T$, $\det(AB)$) & 1.5 · 4.3 · P.6\\
Déterminant, inversibilité, liberté, systèmes de Cramer & 2.1 · 4.1 · 4.2 · P.6\\
« Trouver tous les… » : prouver les deux sens & 2.1(a) · P.5\\\bottomrule
\end{tabular}
\end{center}

\lienDL{Les questions 2.3(b), 4.2 et P.3 à P.5 reprennent les difficultés du DL~3 sous une autre forme : liberté par évaluation
(comme III.3), Vandermonde, famille $(x,Ax,A^2x)$ avec $A^3=0$ (comme II.1 et II.4.a). Les formulations sont différentes : il
faut reconnaître la technique.}

\begin{center}\small\renewcommand{\arraystretch}{1.2}
\begin{tabular}{@{}lcccccc@{}}\toprule
& Ex.~1 & Ex.~2 & Ex.~3 & Ex.~4 & Problème & Total\\\midrule
Points & 3 & 5 & 3,5 & 4 & 4,5 & 20\\
Temps conseillé & 15 min & 50 min & 35 min & 35 min & 45 min & 3 h\\\bottomrule
\end{tabular}
\end{center}

% ---------------------------------------------------------------------
\clearpage
\Ex{1}{Vrai ou faux}{3 points}

\Q{1. Vrai.} $\dim\R^3=3$ : une famille libre de $\R^3$ a au plus 3 vecteurs. Une famille de 4 vecteurs est donc liée.

\Q{2. Faux.} Contre-exemple : $u=(1,0,0)$, $v=(0,1,0)$, $w=(1,0,0)$. Les familles $(u,v)$ et $(v,w)$ sont libres, mais $u-w=0$ :
$(u,v,w)$ est liée.
\piege{la liberté ne se vérifie pas « deux à deux ». Trois vecteurs deux à deux non colinéaires peuvent être liés, par exemple
$(1,0)$, $(0,1)$, $(1,1)$ dans $\R^2$.}

\Q{3. Vrai.} $\dim\R^3=3$ et la famille a 3 vecteurs : pour une telle famille, génératrice $\iff$ libre $\iff$ base.

\Q{4. Faux.} Dans $\R^2$, $F=\vect\big((1,0)\big)$ et $G=\vect\big((0,1)\big)$ : $(1,0)$ et $(0,1)$ sont dans $F\cup G$, mais pas
leur somme $(1,1)$. $F\cup G$ n'est pas stable par addition.

\Q{5. Faux.} Multiplier $A$ par $2$ multiplie chacune des 3 colonnes par $2$, et le déterminant est linéaire par rapport à
\emph{chaque} colonne : $\det(2A)=2^3\det A=8\det A$. Par exemple $\det(2I_3)=8$.
\retenir{$\det(\lambda A)=\lambda^n\det A$ pour $A\in\Mn{\K}$.}

\Q{6. Vrai.} Par le théorème du rang, $\dim\Ker f=4-\rg f$. Or $\Im f\subset\R^3$, donc $\rg f\leq3$ et $\dim\Ker f\geq1$ :
$\Ker f\neq\{0\}$, $f$ n'est pas injective.

% ---------------------------------------------------------------------
\Ex{2}{Familles libres, génératrices, bases}{5 points}

\Partie{A. Une famille de $\R^3$ avec paramètre}

\Q{2.1(a)} \choix{3 vecteurs en dimension 3 avec un paramètre : un seul déterminant donne toutes les valeurs de $m$ d'un coup.
Résoudre un système avec $m$ obligerait à discuter à chaque pivot.}
\[\det(u_1,u_2,u_3)=\begin{vmatrix}1&0&2\\2&1&m\\-1&1&1\end{vmatrix}
\underset{C_3\leftarrow C_3-2C_1}{=}\begin{vmatrix}1&0&0\\2&1&m-4\\-1&1&3\end{vmatrix}=1\times\begin{vmatrix}1&m-4\\1&3\end{vmatrix}=7-m.\]
Le déterminant est non nul si et seulement si $m\neq7$. \textbf{$(u_1,u_2,u_3)$ est une base de $\R^3$ si et seulement si $m\neq7$.}
Les deux sens sont prouvés en même temps, car « base $\iff\det\neq0$ ».

\Q{2.1(b)} Pour $m=1$, on cherche $(a,b,c)$ tel que $au_1+bu_2+cu_3=(1,1,1)$ :
\[\begin{cases}a+2c=1\\2a+b+c=1\\-a+b+c=1\end{cases}\]
La différence des deux dernières lignes donne $3a=0$, donc $a=0$, puis $c=\frac12$ et $b=\frac12$. \textbf{Les coordonnées de
$(1,1,1)$ dans $\mathcal B'$ sont $\left(0,\frac12,\frac12\right)$.} Vérification : $\frac12(u_2+u_3)=\frac12(2,2,2)=(1,1,1)$.
\piege{les coordonnées dépendent de la base : $(1,1,1)$ est la colonne dans la base canonique, $(0,\frac12,\frac12)$ dans
$\mathcal B'$. Toujours vérifier en recombinant.}

\Q{2.1(c)} Pour $m=7$, on cherche une relation $au_1+bu_2+u_3=0$ : la première coordonnée donne $a=-2$, la troisième $b=-3$, et
la deuxième est alors vérifiée. Donc $u_3=2u_1+3u_2$. Comme $u_1$ et $u_2$ ne sont pas colinéaires,
$\rg(u_1,u_2,u_3)=\dim\vect(u_1,u_2)=2$.

\Partie{B. Un sous-espace de $\R^4$}

\Q{2.2(a)} \choix{un sous-espace défini par des équations : on résout le système, et on écrit la solution générale comme
combinaison linéaire. Cela donne directement une famille génératrice.}

$(x,y,z,t)\in F\iff x=y-2z$ et $t=-x-y-z=-2y+z$. Donc
\[F=\{\,y(1,1,0,-2)+z(-2,0,1,1),\ y,z\in\R\,\}=\vect(f_1,f_2),\]
avec $f_1=(1,1,0,-2)$ et $f_2=(-2,0,1,1)$.
\emph{Génératrice :} c'est l'égalité ci-dessus. \emph{Libre :} si $yf_1+zf_2=0$, les deuxième et troisième coordonnées donnent
$y=0$ et $z=0$. \textbf{$(f_1,f_2)$ est une base de $F$, et $\dim F=2$.}
\autre{$F$ est défini par 2 équations indépendantes dans $\R^4$, donc $\dim F=4-2=2$ (théorème du rang appliqué à
$(x,y,z,t)\mapsto(x+y+z+t,\,x-y+2z)$, de rang 2).}

\Q{2.2(b)} \choix{on connaît maintenant $\dim F=2$. Pour 2 vecteurs dans un espace de dimension 2, il suffit qu'ils soient
dans $F$ et libres : inutile de prouver qu'ils engendrent $F$.}

$w_1=(3,1,-1,-3)$ : $3+1-1-3=0$ et $3-1-2=0$, donc $w_1\in F$. $w_2=(0,2,1,-3)$ : $0+2+1-3=0$ et $0-2+2=0$, donc $w_2\in F$. Ils
ne sont pas colinéaires (la première coordonnée de $w_2$ est nulle, pas celle de $w_1$, et $w_2\neq0$). \textbf{$(w_1,w_2)$ est une
famille libre de 2 vecteurs de $F$, avec $\dim F=2$ : c'est une base de $F$.}
\piege{oublier de vérifier que $w_1,w_2\in F$. Deux vecteurs libres de $\R^4$ ne forment pas une base de $F$ s'ils n'y sont pas.}

\Partie{C. Une famille de polynômes}

\Q{2.3(a)}
\begin{center}\small\renewcommand{\arraystretch}{1.2}
\begin{tabular}{@{}cccc@{}}\toprule
& $x=0$ & $x=1$ & $x=2$\\\midrule
$P_1=(X-1)(X-2)$ & $2$ & $0$ & $0$\\
$P_2=X(X-2)$ & $0$ & $-1$ & $0$\\
$P_3=X(X-1)$ & $0$ & $0$ & $2$\\\bottomrule
\end{tabular}
\end{center}

\Q{2.3(b)} \choix{les trois polynômes ont le même degré : l'argument « degrés distincts » ne marche pas. Le tableau montre
que chaque $P_i$ ne vit qu'en un point. On évalue la relation en $0$, $1$, $2$ : chaque évaluation isole un coefficient.}

Soient $a,b,c$ tels que $aP_1+bP_2+cP_3=0$. On évalue en $0$ : $2a=0$. En $1$ : $-b=0$. En $2$ : $2c=0$. Donc $a=b=c=0$ : la
famille est libre. Elle a 3 vecteurs et $\dim\R_2[X]=3$ : \textbf{c'est une base de $\R_2[X]$}.
\lienDL{c'est la même idée qu'en III.3 : une relation entre vecteurs devient une information sur des valeurs, qu'on lit point
par point.}

\Q{2.3(c)} D'après (b), $P=aP_1+bP_2+cP_3$ pour des réels $a,b,c$. On évalue en $0$, $1$, $2$ :
$P(0)=2a$, $P(1)=-b$, $P(2)=2c$. Donc
\[P=\frac{P(0)}2\,P_1-P(1)\,P_2+\frac{P(2)}2\,P_3.\]

\Q{2.3(d)} D'après (c), si $P(0)=1$, $P(1)=0$, $P(2)=3$, nécessairement $P=\frac12P_1+\frac32P_3=2X^2-3X+1$. Réciproquement,
ce polynôme vérifie $P(0)=1$, $P(1)=0$, $P(2)=8-6+1=3$. \textbf{L'unique solution est $P=2X^2-3X+1$.}

\Partie{D. Un raisonnement sans coordonnées}

\Q{2.4(a)} \choix{on ne connaît pas les vecteurs, seulement que $(u,v,w)$ est une base. On écrit une combinaison nulle et on la
réécrit en fonction de $u,v,w$ pour utiliser leur liberté.}

Soient $\alpha,\beta,\gamma$ tels que $\alpha(u+v)+\beta(v+w)+\gamma(w+u)=0$, soit $(\alpha+\gamma)u+(\alpha+\beta)v+(\beta+\gamma)w=0$.
Comme $(u,v,w)$ est libre : $\alpha+\gamma=\alpha+\beta=\beta+\gamma=0$. En additionnant, $\alpha+\beta+\gamma=0$, puis
$\beta=\gamma=\alpha=0$. La famille est libre, a 3 vecteurs, et $\dim E=3$ : \textbf{c'est une base de $E$.}
\autre{la matrice de la famille dans la base $(u,v,w)$ est $\begin{pmatrix}1&0&1\\1&1&0\\0&1&1\end{pmatrix}$, de déterminant
$2\neq0$.}

\Q{2.4(b)} $(u-v)+(v-w)+(w-u)=0$ : c'est une relation à coefficients non tous nuls. \textbf{La famille est liée.}

\Q{2.4(c)} $(u+v)-(v+w)+(w+u)=2u$, donc $u=\frac12(u+v)-\frac12(v+w)+\frac12(w+u)$ : \textbf{coordonnées
$\left(\frac12,-\frac12,\frac12\right)$}.

% ---------------------------------------------------------------------
\Ex{3}{Sous-espaces supplémentaires et projecteur}{3,5 points}

\Q{3.1} L'application $\varphi:P\mapsto\big(P(1),P'(1)\big)$, de $E$ dans $\R^2$, est linéaire (l'évaluation et la dérivation le
sont). $F=\Ker\varphi$ est donc un sous-espace vectoriel de $E$.
\autre{vérifier directement : le polynôme nul est dans $F$, et si $P,Q\in F$, alors $(\lambda P+\mu Q)(1)=0$ et
$(\lambda P+\mu Q)'(1)=0$.}

\Q{3.2(a)} Les polynômes $1$, $X-1$, $(X-1)^2$, $(X-1)^3$ sont non nuls et de degrés $0,1,2,3$ distincts : la famille est libre.
Elle a 4 vecteurs et $\dim E=4$ : \textbf{$\mathcal B_1$ est une base de $E$}.

\Q{3.2(b)} \choix{la formule de Taylor donne directement les coordonnées dans $\mathcal B_1$, et les conditions qui
définissent $F$ portent exactement sur les deux premières.}

Par la formule de Taylor, les coordonnées de $P$ dans $\mathcal B_1$ sont $\left(P(1),\,P'(1),\,\frac{P''(1)}2,\,\frac{P'''(1)}6\right)$.
Donc $P\in F$ si et seulement si ses deux premières coordonnées sont nulles, c'est-à-dire si et seulement si
$P\in\vect\big((X-1)^2,(X-1)^3\big)$. La famille $\big((X-1)^2,(X-1)^3\big)$ engendre $F$, et elle est libre car extraite de la
base $\mathcal B_1$. \textbf{C'est une base de $F$, et $\dim F=2$.}

\Q{3.3} \choix{deux méthodes possibles. La plus rapide : intersection nulle et somme des dimensions.}

Soit $P\in F\cap G$. Comme $P\in G=\R_1[X]$, $P=a+b(X-1)$, donc $P(1)=a$ et $P'(1)=b$. Comme $P\in F$, $a=b=0$ : $P=0$. Ainsi
$F\cap G=\{0\}$, et $\dim F+\dim G=2+2=4=\dim E$. \textbf{Donc $E=F\oplus G$.}
\autre{$G=\vect\big(1,X-1\big)$ et $F=\vect\big((X-1)^2,(X-1)^3\big)$ : en réunissant ces deux bases, on obtient la base
$\mathcal B_1$ de $E$, donc $E=F\oplus G$.}

\Q{3.4(a)} D'après la formule de Taylor,
\[P=\underbrace{P(1)+P'(1)(X-1)}_{\in\,G}+\underbrace{\frac{P''(1)}2(X-1)^2+\frac{P'''(1)}6(X-1)^3}_{\in\,F}.\]
La décomposition de $P$ selon $E=F\oplus G$ est unique ; la composante dans $G$ est $q(P)$. \textbf{Donc
$q(P)=P(1)+P'(1)(X-1)$.} Pour $P=X^3$ : $q(X^3)=1+3(X-1)=3X-2$. Contrôle : $X^3-(3X-2)=(X-1)^2(X+2)$ est bien dans $F$.

\Q{3.4(b)} $q$ laisse fixes les vecteurs de $G$ et annule ceux de $F$ :
\[\Mat_{\mathcal B_1}(q)=\begin{pmatrix}1&0&0&0\\0&1&0&0\\0&0&0&0\\0&0&0&0\end{pmatrix}.\]
Dans $\mathcal B_0$ : $q(1)=1$, $q(X)=X$, $q(X^2)=1+2(X-1)=2X-1$, $q(X^3)=3X-2$, donc
\[\Mat_{\mathcal B_0}(q)=\begin{pmatrix}1&0&-1&-2\\0&1&2&3\\0&0&0&0\\0&0&0&0\end{pmatrix}.\]
On vérifie que cette matrice $Q$ vérifie $Q^2=Q$, comme tout projecteur.
\retenir{dans une base adaptée à $E=F\oplus G$, la matrice d'un projecteur est diagonale avec des $1$ et des $0$. Choisir la
bonne base simplifie tout : c'est l'idée de toute la suite du cours.}

% ---------------------------------------------------------------------
\Ex{4}{Calculs de déterminants}{4 points}

\Q{4.1} \choix{la somme de chaque ligne vaut $a+3b$ : on additionne toutes les colonnes dans la première pour factoriser.}
\begin{align*}D_4&=\begin{vmatrix}a&b&b&b\\b&a&b&b\\b&b&a&b\\b&b&b&a\end{vmatrix}
\underset{C_1\leftarrow C_1+C_2+C_3+C_4}{=}(a+3b)\begin{vmatrix}1&b&b&b\\1&a&b&b\\1&b&a&b\\1&b&b&a\end{vmatrix}\\
&\underset{L_i\leftarrow L_i-L_1}{=}(a+3b)\begin{vmatrix}1&b&b&b\\0&a-b&0&0\\0&0&a-b&0\\0&0&0&a-b\end{vmatrix}.\end{align*}
La dernière matrice est triangulaire : \textbf{$D_4=(a+3b)(a-b)^3$}. La matrice est inversible si et seulement si $D_4\neq0$,
c'est-à-dire \textbf{$a\neq b$ et $a\neq-3b$}.
\piege{développer directement un $4\times4$ : 24 termes et presque sûrement une erreur. Toujours chercher d'abord une
structure (somme constante, colonnes proches, zéros).}

\Q{4.2(a)} On fait $L_2\leftarrow L_2-L_1$ et $L_3\leftarrow L_3-L_1$, puis on factorise :
\[V=\begin{vmatrix}1&a&a^2\\0&b-a&b^2-a^2\\0&c-a&c^2-a^2\end{vmatrix}=(b-a)(c-a)\begin{vmatrix}1&b+a\\1&c+a\end{vmatrix}=(b-a)(c-a)(c-b).\]
\gain{$b^2-a^2=(b-a)(b+a)$ : on fait apparaître les facteurs $(b-a)$ et $(c-a)$ avant de développer.}

\Q{4.2(b)} $P=x_0+x_1X+x_2X^2$ vérifie les trois conditions si et seulement si
\[\begin{pmatrix}1&a&a^2\\1&b&b^2\\1&c&c^2\end{pmatrix}\begin{pmatrix}x_0\\x_1\\x_2\end{pmatrix}=\begin{pmatrix}\alpha\\\beta\\\gamma\end{pmatrix}.\]
Le déterminant de ce système est $V\neq0$ car $a,b,c$ sont distincts : c'est un système de Cramer, il a une unique solution.
\textbf{Il existe un unique $P\in\R_2[X]$ qui convient.}
\lienDL{c'est le même phénomène qu'en 2.3, où $a=0$, $b=1$, $c=2$, et qu'en III.3 du DL : trois valeurs distinctes
« séparent » les polynômes de degré au plus 2. Ici on le prouve par le déterminant, en 2.3 par évaluation.}

\Q{4.3} \choix{on ne peut rien calculer : on utilise les propriétés, $\det A^T=\det A$ et $\det(\lambda A)=\lambda^n\det A$.}

$A^T=-A$, donc $\det A=\det A^T=\det(-A)=(-1)^n\det A=-\det A$ car $n$ est impair. Ainsi $2\det A=0$ : \textbf{$\det A=0$, et
$A$ n'est pas inversible.} Pour $n$ pair, c'est faux : $\begin{pmatrix}0&1\\-1&0\end{pmatrix}$ est antisymétrique, de
déterminant $1$.
\piege{écrire $\det(-A)=-\det A$ sans justifier : c'est vrai seulement parce que $n$ est impair.}

% ---------------------------------------------------------------------
\Pb{une matrice nilpotente}{4,5 points}

\Q{P.1}
\[A^2=\begin{pmatrix}0&0&0\\1&1&1\\-1&-1&-1\end{pmatrix},\qquad A^3=A^2A=0.\]
(Chaque ligne de $A^2$ est proportionnelle à $(1,1,1)$, et $(1,1,1)A=(0,0,0)$.)

\Q{P.2} \emph{Noyau.} $AX=0$ s'écrit $-x-y-z=0$, $x+2y+2z=0$, $-y-z=0$. La dernière équation donne $z=-y$, puis la première
$x=0$ ; la deuxième est alors vérifiée. \textbf{$\Ker A=\vect\big((0,1,-1)\big)$}, de dimension 1.

\emph{Image.} Elle est engendrée par les colonnes $C_1=(-1,1,0)$, $C_2=(-1,2,-1)$, $C_3=C_2$. $C_1$ et $C_2$ ne sont pas
colinéaires : \textbf{$\Im A=\vect(C_1,C_2)$, et $\rg A=2$}. Théorème du rang : $1+2=3$. \checkmark

\emph{Équation.} $C_1$ et $C_2$ vérifient $x+y+z=0$, et ce plan est de dimension 2 : \textbf{$\Im A$ a pour équation $x+y+z=0$}.

\emph{Inclusion.} $(0,1,-1)$ vérifie $0+1-1=0$, donc $\Ker A\subset\Im A$.

\Q{P.3} \choix{l'énoncé interdit le déterminant : on utilise $A^3=0$. Appliquer $A^2$ à une relation fait disparaître tous
les termes sauf un.}

On a $x_0=e_1$, $Ax_0=(-1,1,0)$, $A^2x_0=(0,1,-1)\neq0$. Soient $c_0,c_1,c_2$ tels que $c_0x_0+c_1Ax_0+c_2A^2x_0=0$.
\begin{itemize}[leftmargin=1.2em,itemsep=1pt]
\item On multiplie par $A^2$ : comme $A^3=A^4=0$, il reste $c_0A^2x_0=0$, et $A^2x_0\neq0$, donc $c_0=0$.
\item On multiplie alors par $A$ : il reste $c_1A^2x_0=0$, donc $c_1=0$.
\item Il reste $c_2A^2x_0=0$, donc $c_2=0$.
\end{itemize}
La famille est libre, de 3 vecteurs en dimension 3 : \textbf{$\mathcal B$ est une base de $\R^3$.}
\lienDL{c'est exactement la question II.1 du DL, avec une autre matrice.}

\Q{P.4} $f(x_0)=Ax_0$ est le 2\textsuperscript{e} vecteur de $\mathcal B$, $f(Ax_0)=A^2x_0$ le 3\textsuperscript{e}, et
$f(A^2x_0)=A^3x_0=0$ :
\[N=\Mat_{\mathcal B}(f)=\begin{pmatrix}0&0&0\\1&0&0\\0&1&0\end{pmatrix},\qquad
P=\begin{pmatrix}1&-1&0\\0&1&1\\0&0&-1\end{pmatrix},\qquad A=PNP^{-1}.\]
$P$ a pour colonnes les coordonnées de $x_0$, $Ax_0$, $A^2x_0$ dans la base canonique. La relation vient de la formule de
changement de base $N=P^{-1}AP$.

\Q{P.5} \choix{« tous les $x$ » : deux sens, et la question P.3 fait le premier.}

Soit $x=(x_1,x_2,x_3)$. On a $A^2x=(x_1+x_2+x_3)(0,1,-1)$.
\begin{itemize}[leftmargin=1.2em,itemsep=1pt]
\item Si $x_1+x_2+x_3\neq0$, alors $A^2x\neq0$ et le raisonnement de P.3 s'applique mot pour mot : la famille est une base.
\item Si $x_1+x_2+x_3=0$, alors $A^2x=0$ : la famille contient le vecteur nul, elle est liée.
\end{itemize}
\textbf{$(x,Ax,A^2x)$ est une base si et seulement si $x_1+x_2+x_3\neq0$}, c'est-à-dire $x\notin\Ker A^2$.
\lienDL{c'est la question II.4.a du DL. La condition ne se devine pas : elle vient du calcul de $A^2x$.}

\Q{P.6} \choix{calculer directement $\det(A+\lambda I_3)$ marche, mais la relation $A=PNP^{-1}$ donne le résultat sans calcul :
deux matrices semblables ont le même déterminant.}

$A+\lambda I_3=PNP^{-1}+\lambda PP^{-1}=P(N+\lambda I_3)P^{-1}$, donc
\[\det(A+\lambda I_3)=\det P\,\det(N+\lambda I_3)\,\det P^{-1}=\det(N+\lambda I_3)=\lambda^3,\]
car $N+\lambda I_3$ est triangulaire, de diagonale $(\lambda,\lambda,\lambda)$. Avec $\lambda=1$ : $\det(A+I_3)=1\neq0$,
\textbf{$A+I_3$ est inversible}. Avec $\lambda=0$ : $\det A=0$, \textbf{$A$ n'est pas inversible} (on le savait : $\rg A=2$).
\gain{on peut même donner l'inverse : $(I_3+A)(I_3-A+A^2)=I_3+A^3=I_3$.}

% ---------------------------------------------------------------------
\clearpage
\markboth{Diagnostic}{}
\subsection*{Après la correction : quoi revoir ?}
\begin{center}
\small\renewcommand{\arraystretch}{1.25}
\begin{tabular}{@{}>{\raggedright\arraybackslash}p{5.6cm}>{\raggedright\arraybackslash}p{7.4cm}@{}}
\toprule
\textbf{Si tu as perdu des points en…} & \textbf{Revoir}\\\midrule
1.2, 2.3(b), 2.4(a), P.3 (liberté) & Livret d'algèbre linéaire, méthode R.2 ; livret du DL, outils 1 et 2\\
2.2, 3.2 (génératrice, base d'un sous-espace) & Livret d'algèbre linéaire, rappels « Espaces vectoriels » et « Dimension finie »\\
2.1(b), 2.4(c), 3.4, P.4 (coordonnées, matrices) & Livret d'algèbre linéaire, proposition « changement de base » ; livret du DL, outil 3\\
3.3, 3.4 (supplémentaires, projecteur) & Livret d'algèbre linéaire, §1.2\\
P.2 (noyau, image, rang) & Livret d'algèbre linéaire, méthode R.3 « Calculer un rang, un noyau, une image »\\
4.1, 4.2 (calculs) & Livret des déterminants, §1.2 et §1.5 (somme constante, Vandermonde)\\
1.5, 4.3, P.6 (propriétés) & Livret des déterminants, §1.3\\
2.1(a), P.5 (« tous les… ») & Livret du DL, outil 4\\\bottomrule
\end{tabular}
\end{center}

\begin{essentiel}
\begin{itemize}[leftmargin=1.1em,itemsep=3pt]
\item Libre : combinaison nulle, puis le bon geste pour isoler les coefficients (identifier, évaluer, composer par $A^k$).
\item $n$ vecteurs en dimension $n$ : une seule propriété à prouver. Pour une base d'un sous-espace $F$, vérifier aussi que les
vecteurs sont dans $F$.
\item Un paramètre et $n$ vecteurs en dimension $n$ : le déterminant donne toutes les valeurs d'un coup.
\item Un déterminant se calcule par opérations sur les lignes et colonnes avant de développer.
\item $\det(\lambda A)=\lambda^n\det A$, $\det A^T=\det A$, deux matrices semblables ont le même déterminant.
\end{itemize}
\end{essentiel}

\else
% =====================================================================
%                               SUJET
% =====================================================================
\markboth{Examen blanc · Algèbre linéaire et déterminants}{}
\begin{center}
\begin{tcolorbox}[enhanced,colback=Fond,colframe=Fond,arc=0pt,boxrule=0pt,left=10pt,right=10pt,top=8pt,bottom=8pt,width=\linewidth]
\small
\textbf{\sffamily\color{Bleu}Consignes.} Le sujet comporte quatre exercices et un problème, indépendants : on peut les traiter
dans l'ordre de son choix. Toute réponse doit être \emph{justifiée} : citer le résultat du cours utilisé et vérifier ses
hypothèses. Les résultats sont à encadrer. Une réponse juste sans justification ne rapporte pas de points ; une démarche correcte
avec une erreur de calcul en rapporte. La clarté et la précision de la rédaction sont prises en compte.

\medskip
\begin{tabular}{@{}lcccccc@{}}
\textbf{\sffamily\color{Bleu}Barème} & Ex.~1 : 3 & Ex.~2 : 5 & Ex.~3 : 3,5 & Ex.~4 : 4 & Problème : 4,5 & \textbf{Total : 20}\\
\end{tabular}
\end{tcolorbox}
\end{center}

\Ex{1}{Vrai ou faux}{3 points}
Pour chaque affirmation, dire si elle est vraie ou fausse, et le justifier : une preuve courte si elle est vraie, un
contre-exemple si elle est fausse. \pts{0,5 point par affirmation}
\begin{enumerate}
\item Toute famille de quatre vecteurs de $\R^3$ est liée.
\item Si $(u,v)$ et $(v,w)$ sont deux familles libres d'un espace vectoriel $E$, alors $(u,v,w)$ est libre.
\item Une famille de trois vecteurs qui engendre $\R^3$ est une base de $\R^3$.
\item La réunion de deux sous-espaces vectoriels de $\R^2$ est un sous-espace vectoriel de $\R^2$.
\item Pour toute matrice $A\in\mathcal M_3(\R)$, $\det(2A)=2\det A$.
\item Une application linéaire de $\R^4$ dans $\R^3$ n'est jamais injective.
\end{enumerate}

\Ex{2}{Familles libres, génératrices, bases}{5 points}

\Partie{A. Une famille de $\R^3$ avec paramètre}
Pour $m\in\R$, on pose $u_1=(1,2,-1)$, $u_2=(0,1,1)$ et $u_3=(2,m,1)$.
\begin{enumerate}
\item
\begin{enumerate}
\item Déterminer les valeurs de $m$ pour lesquelles $(u_1,u_2,u_3)$ est une base de $\R^3$. \pts{0,75 pt}
\item On prend $m=1$ et on note $\mathcal B'=(u_1,u_2,u_3)$. Déterminer les coordonnées du vecteur $(1,1,1)$ dans $\mathcal B'$.
\pts{0,5 pt}
\item Pour la valeur de $m$ exclue au (a), écrire $u_3$ comme combinaison linéaire de $u_1$ et $u_2$, et donner le rang de
$(u_1,u_2,u_3)$. \pts{0,25 pt}
\end{enumerate}
\end{enumerate}

\Partie{B. Un sous-espace de $\R^4$}
Soit $F=\{(x,y,z,t)\in\R^4,\ x+y+z+t=0\ \text{et}\ x-y+2z=0\}$.
\begin{enumerate}[resume]
\item
\begin{enumerate}
\item Déterminer une famille génératrice de $F$, puis montrer que c'est une base de $F$. En déduire $\dim F$. \pts{0,75 pt}
\item Montrer que $w_1=(3,1,-1,-3)$ et $w_2=(0,2,1,-3)$ forment une base de $F$. \pts{0,5 pt}
\end{enumerate}
\end{enumerate}

\Partie{C. Une famille de polynômes}
Dans $\R_2[X]$, on pose $P_1=(X-1)(X-2)$, $P_2=X(X-2)$ et $P_3=X(X-1)$.
\begin{enumerate}[resume]
\item
\begin{enumerate}
\item Calculer $P_i(0)$, $P_i(1)$ et $P_i(2)$ pour $i=1,2,3$. \pts{0,25 pt}
\item Montrer que $(P_1,P_2,P_3)$ est une base de $\R_2[X]$. \pts{0,5 pt}
\item Soit $P\in\R_2[X]$. Exprimer les coordonnées de $P$ dans cette base à l'aide de $P(0)$, $P(1)$ et $P(2)$. \pts{0,25 pt}
\item En déduire l'unique polynôme $P\in\R_2[X]$ tel que $P(0)=1$, $P(1)=0$ et $P(2)=3$. \pts{0,25 pt}
\end{enumerate}
\end{enumerate}

\Partie{D. Un raisonnement sans coordonnées}
Soit $E$ un espace vectoriel de dimension 3 et $(u,v,w)$ une base de $E$.
\begin{enumerate}[resume]
\item
\begin{enumerate}
\item Montrer que $(u+v,\ v+w,\ w+u)$ est une base de $E$. \pts{0,5 pt}
\item La famille $(u-v,\ v-w,\ w-u)$ est-elle une base de $E$ ? \pts{0,25 pt}
\item Donner les coordonnées de $u$ dans la base $(u+v,\ v+w,\ w+u)$. \pts{0,25 pt}
\end{enumerate}
\end{enumerate}

\Ex{3}{Sous-espaces supplémentaires et projecteur}{3,5 points}
Soit $E=\R_3[X]$, de base canonique $\mathcal B_0=(1,X,X^2,X^3)$. On pose
\[F=\{P\in E,\ P(1)=0\ \text{et}\ P'(1)=0\}\qquad\text{et}\qquad G=\R_1[X].\]
On rappelle la formule de Taylor : pour tout $P\in E$,
$P=P(1)+P'(1)(X-1)+\frac{P''(1)}{2}(X-1)^2+\frac{P'''(1)}{6}(X-1)^3$.
\begin{enumerate}
\item Montrer que $F$ est un sous-espace vectoriel de $E$. \pts{0,5 pt}
\item
\begin{enumerate}
\item Montrer que $\mathcal B_1=\big(1,\ X-1,\ (X-1)^2,\ (X-1)^3\big)$ est une base de $E$. \pts{0,5 pt}
\item Donner les coordonnées d'un polynôme $P$ dans $\mathcal B_1$. En déduire une base de $F$ et sa dimension. \pts{0,75 pt}
\end{enumerate}
\item Montrer que $E=F\oplus G$. \pts{0,75 pt}
\item On note $q$ le projecteur sur $G$ parallèlement à $F$.
\begin{enumerate}
\item Montrer que $q(P)=P(1)+P'(1)(X-1)$ pour tout $P\in E$. Calculer $q(X^3)$. \pts{0,5 pt}
\item Écrire la matrice de $q$ dans la base $\mathcal B_1$, puis dans la base $\mathcal B_0$. \pts{0,5 pt}
\end{enumerate}
\end{enumerate}

\Ex{4}{Calculs de déterminants}{4 points}
\begin{enumerate}
\item Soient $a,b\in\R$. Calculer, sous forme factorisée, \pts{1,5 pt}
\[D_4=\begin{vmatrix}a&b&b&b\\b&a&b&b\\b&b&a&b\\b&b&b&a\end{vmatrix}.\]
Pour quelles valeurs de $a$ et $b$ la matrice correspondante est-elle inversible ?
\item Soient $a,b,c$ trois réels.
\begin{enumerate}
\item Calculer, sous forme factorisée, $V=\begin{vmatrix}1&a&a^2\\1&b&b^2\\1&c&c^2\end{vmatrix}$. \pts{1 pt}
\item On suppose $a,b,c$ deux à deux distincts. Soient $\alpha,\beta,\gamma\in\R$. Montrer qu'il existe un unique polynôme
$P\in\R_2[X]$ tel que $P(a)=\alpha$, $P(b)=\beta$ et $P(c)=\gamma$. \pts{0,5 pt}
\end{enumerate}
\item Soit $n$ un entier impair et $A\in\mathcal M_n(\R)$ une matrice antisymétrique ($A^T=-A$). Montrer que $A$ n'est pas
inversible. Le résultat reste-t-il vrai si $n$ est pair ? \pts{1 pt}
\end{enumerate}

\Pb{une matrice nilpotente}{4,5 points}
On considère
\[A=\begin{pmatrix}-1&-1&-1\\1&2&2\\0&-1&-1\end{pmatrix}\]
et $f$ l'endomorphisme de $\R^3$ canoniquement associé à $A$. On note $(e_1,e_2,e_3)$ la base canonique de $\R^3$.
\begin{enumerate}
\item Calculer $A^2$ et $A^3$. \pts{0,5 pt}
\item Déterminer une base de $\Ker A$ et une base de $\Im A$. Vérifier le théorème du rang, donner une équation de $\Im A$, et
montrer que $\Ker A\subset\Im A$. \pts{1 pt}
\item On pose $x_0=e_1$. Montrer, \emph{sans calculer de déterminant}, que $\mathcal B=(x_0,Ax_0,A^2x_0)$ est une base de
$\R^3$. \pts{1 pt}
\item Écrire la matrice $N$ de $f$ dans la base $\mathcal B$, la matrice de passage $P$ de la base canonique à $\mathcal B$, et
la relation entre $A$, $N$ et $P$. \pts{0,5 pt}
\item Déterminer tous les vecteurs $x=(x_1,x_2,x_3)$ de $\R^3$ tels que $(x,Ax,A^2x)$ soit une base de $\R^3$. \pts{1 pt}
\item Pour $\lambda\in\R$, calculer $\det(A+\lambda I_3)$ en utilisant la question 4. En déduire que $A+I_3$ est inversible et
que $A$ ne l'est pas. \pts{0,5 pt}
\end{enumerate}

\vfill
\begin{center}\sffamily\small\color{Gris}— Fin du sujet —\end{center}
\fi
\end{document}
"
\documentclass[a4paper,11pt]{article}

% ============================================================
%  Le Collé — Algèbre linéaire (compléments de deuxième année)
%  © 2026 Ahmed Amine Ouqlani
% ============================================================

% ---------- Polices ----------
\usepackage[utf8]{inputenc}
\usepackage[T1]{fontenc}
\usepackage[sc]{mathpazo}          % texte et mathématiques : Palatino
\usepackage[scaled=0.88]{helvet}   % titres et repères : linéale
\linespread{1.07}
\usepackage{amsmath,amssymb,amsthm,mathrsfs,stmaryrd}
\usepackage{microtype}

% ---------- Mise en page ----------
\usepackage[a4paper,top=2.3cm,bottom=2.3cm,left=2.4cm,right=2.4cm,
            headheight=14pt,headsep=0.8cm,footskip=1.1cm]{geometry}
\setlength{\parindent}{0pt}
\setlength{\parskip}{0.45em plus 0.15em minus 0.1em}
\raggedbottom
\emergencystretch=1.5em
\clubpenalty=10000 \widowpenalty=10000 \displaywidowpenalty=10000 \brokenpenalty=10000

% ---------- Couleurs : un bleu de nuit, un ocre, des gris ----------
\usepackage{xcolor}
\definecolor{Ink}{HTML}{1C2533}
\definecolor{Bleu}{HTML}{1F3A5F}
\definecolor{BleuMoyen}{HTML}{4A6A8F}
\definecolor{BleuClair}{HTML}{9DB2C9}
\definecolor{Fond}{HTML}{F2F4F7}
\definecolor{Ocre}{HTML}{B07A26}
\definecolor{Brique}{HTML}{A23B2A}
\definecolor{Gris}{HTML}{5E6672}
\definecolor{Filet}{HTML}{C5CCD6}

% ---------- Packages ----------
\usepackage{enumitem}
\setlist{itemsep=0.25em,topsep=0.35em,parsep=0pt}
\setlist[itemize]{label=\textcolor{BleuMoyen}{\small\textbullet},leftmargin=1.4em}
\setlist[enumerate]{leftmargin=1.8em,label=\textsf{\small\color{Bleu}\arabic*.}}
\usepackage{booktabs,array,multicol}
\usepackage{needspace,float}
\usepackage[font=small,labelfont={sf,color=Bleu},labelsep=period,skip=6pt]{caption}
\usepackage{tikz}
\usetikzlibrary{arrows.meta,calc,decorations.markings}
\usepackage{pgfplots}
\pgfplotsset{compat=1.17}
\usepackage{tkz-tab}
\usepackage[most]{tcolorbox}
\usepackage{titlesec,titletoc}
\usepackage{fancyhdr,lastpage,etoolbox}
\usepackage[colorlinks=true,linkcolor=Bleu,urlcolor=Bleu,citecolor=Bleu,
            pdftitle={Examen blanc — Algèbre linéaire et déterminants},pdfauthor={Ahmed Amine Ouqlani}]{hyperref}

\numberwithin{equation}{section}
\numberwithin{figure}{section}
\numberwithin{table}{section}

% ---------- Titres ----------
% numéro du chapitre dans le parcours PT (affiché), distinct de la numérotation interne
\newcommand{\grandnum}{\thesection}\newcommand{\ongletcouleur}{Bleu}
% Chapitres (sections) : grand numéro, titre en linéale, filet ocre
\titleformat{\section}[display]
  {\normalfont\sffamily\color{Ink}}
  {{\small\color{Gris}\textls[200]{\MakeUppercase{Chapitre}}}\\[-0.2em]%
   {\fontsize{64}{64}\selectfont\color{BleuClair}\grandnum}}
  {0.4em}
  {\fontsize{25}{30}\selectfont\bfseries}
  [\vspace{0.5em}{\color{Ocre}\rule{3.2cm}{1.4pt}}]
\titlespacing*{\section}{0pt}{0pt}{1.6em}
\titleformat{name=\section,numberless}[display]
  {\normalfont\sffamily\color{Ink}}{}{0pt}{\fontsize{25}{30}\selectfont\bfseries}
  [\vspace{0.5em}{\color{Ocre}\rule{3.2cm}{1.4pt}}]
\titlespacing*{name=\section,numberless}{0pt}{0pt}{1.6em}

\titleformat{\subsection}
  {\needspace{5\baselineskip}\normalfont\sffamily\Large\bfseries\color{Ink}}
  {\color{Bleu}\thesubsection}{0.8em}{}
\titlespacing*{\subsection}{0pt}{1.8em}{0.5em}
\titleformat{\subsubsection}
  {\needspace{5\baselineskip}\normalfont\sffamily\large\bfseries\color{Ink}}
  {\color{Bleu}\thesubsubsection}{0.7em}{}
\titlespacing*{\subsubsection}{0pt}{1.6em}{0.4em}
\titleformat{\paragraph}[runin]{\normalfont\sffamily\bfseries\color{Ink}}{}{0pt}{}[.]
\titlespacing*{\paragraph}{0pt}{0.9em}{0.6em}
\setcounter{secnumdepth}{3}
\setcounter{tocdepth}{2}

% ---------- Table des matières ----------
\titlecontents{section}[0pt]{\addvspace{0.7em}\sffamily\bfseries\color{Ink}}
  {}{}{\hfill\contentspage}[\addvspace{0.2em}]
\titlecontents{subsection}[2.2em]{\small\linespread{1}\selectfont}
  {\makebox[2.6em][l]{\color{Gris}\thecontentslabel}}{}{\titlerule*[0.7pc]{.}\contentspage}

% ---------- En-têtes et pieds de page ----------
\pagestyle{fancy}
\fancyhf{}
\renewcommand{\sectionmark}[1]{\markboth{#1}{}}
\fancyhead[L]{\footnotesize\sffamily\color{Gris}\leftmark}
\fancyhead[R]{\footnotesize\sffamily\scshape\color{Gris}Le Collé}
\fancyfoot[C]{\small\sffamily\color{Ink}\thepage}
\fancyfoot[L]{\scriptsize\sffamily\hyperlink{toc}{\color{Filet}sommaire}}
\renewcommand{\headrulewidth}{0.4pt}
\renewcommand{\headrule}{{\color{Filet}\hrule height\headrulewidth width\headwidth}}
\renewcommand{\footrulewidth}{0pt}
\fancypagestyle{plain}{\fancyhf{}\fancyfoot[C]{\small\sffamily\color{Ink}\thepage}\renewcommand{\headrulewidth}{0pt}}

% ---------- Énoncés ----------
\newtheoremstyle{enonce}{0pt}{0pt}{}{}{\sffamily\bfseries\color{Bleu}}{.}{0.5em}
  {\thmname{#1}\thmnumber{~#2}\thmnote{\ \normalfont\sffamily\color{Ink}(#3)}}
\newtheoremstyle{exemple}{0.9em}{0.9em}{}{}{\sffamily\bfseries\color{BleuMoyen}}{.}{0.5em}
  {\thmname{#1}\thmnumber{~#2}\thmnote{\ \normalfont\sffamily\color{Ink}(#3)}}
\newtheoremstyle{remarque}{0.7em}{0.7em}{}{}{\sffamily\itshape}{.}{0.5em}{\thmname{#1}\thmnote{\ (#3)}}
\theoremstyle{enonce}
\newtheorem{definition}{Définition}[section]
\newtheorem{theoreme}[definition]{Théorème}
\newtheorem{proposition}[definition]{Proposition}
\newtheorem{corollaire}[definition]{Corollaire}
\theoremstyle{exemple}
\newtheorem{exemple}[definition]{Exemple}
\theoremstyle{remarque}
\newtheorem*{remarque}{Remarque}

\tcbset{
  defi/.style={enhanced,breakable,colback=white,colframe=white,boxrule=0pt,arc=0pt,
     borderline west={1.2pt}{0pt}{Bleu},left=11pt,right=4pt,top=3pt,bottom=3pt,
     before skip=1.1em,after skip=1.1em},
  thm/.style={enhanced,breakable,colback=Fond,colframe=Fond,boxrule=0pt,arc=0pt,
     borderline west={2.4pt}{0pt}{Bleu},left=11pt,right=8pt,top=6pt,bottom=6pt,
     before skip=1.2em,after skip=1.2em}}
\tcolorboxenvironment{definition}{defi}
\tcolorboxenvironment{theoreme}{thm}
\tcolorboxenvironment{proposition}{thm}
\tcolorboxenvironment{corollaire}{thm}

\renewcommand{\proofname}{Démonstration}
\renewcommand{\qedsymbol}{\textcolor{BleuClair}{$\blacksquare$}}
\makeatletter
\renewenvironment{proof}[1][\proofname]{\par\pushQED{\qed}\normalfont
  \topsep0.5em\relax\trivlist\item[\hskip\labelsep\sffamily\itshape\color{Gris}#1.]\ignorespaces}
  {\popQED\endtrivlist\@endpefalse}
\makeatother

% ---------- Méthode : filets haut et bas, titre en bandeau ----------
\makeatletter\newcommand{\methlabelfix}{\protected@edef\@currentlabel{\thetcbcounter}}\makeatother
\newtcolorbox[auto counter,number within=section]{methode}[2][]{%
  enhanced,unbreakable,colback=white,colframe=white,boxrule=0pt,arc=0pt,
  borderline north={0.9pt}{0pt}{Bleu},borderline south={0.9pt}{0pt}{Bleu},
  left=4pt,right=4pt,top=8pt,bottom=8pt,before skip=1.4em,after skip=1.4em,
  before upper={\methlabelfix\phantomsection\addcontentsline{met}{metentry}{\protect\makebox[2.6em][l]{\thetcbcounter}#2}%
     {\sffamily\bfseries\small\color{white}\colorbox{Bleu}{\strut\ MÉTHODE~\thetcbcounter\ }}%
     \ \ {\sffamily\bfseries\color{Bleu}#2}\par\smallskip},#1}

% ---------- Repères pédagogiques : un filet de couleur, un mot-clé ----------
\newcommand{\definirbloc}[3]{% nom, couleur, libellé
  \newtcolorbox{#1}[1][]{enhanced,breakable,blanker,
    borderline west={1.4pt}{0pt}{#2},left=10pt,right=2pt,top=1pt,bottom=1pt,
    before skip=1em,after skip=1em,
    before upper={{\sffamily\bfseries\footnotesize\color{#2}\textls[120]{\MakeUppercase{#3}}}\quad},##1}}
\definirbloc{reflexe}{Bleu}{Réflexe}
\definirbloc{astuce}{Ocre}{Astuce}
\definirbloc{attention}{Brique}{Attention}
\definirbloc{enpratique}{BleuMoyen}{En pratique}
\definirbloc{plusloin}{Gris}{Pour aller plus loin}
\definirbloc{strategie}{Bleu}{Stratégie}
\definirbloc{lecture}{BleuMoyen}{Lire le schéma}
% compatibilité avec les parties reprises
\newenvironment{garde}{\begin{attention}}{\end{attention}}
\newenvironment{complementHP}{\begin{plusloin}}{\end{plusloin}}
\newenvironment{interpretation}{\begin{enpratique}}{\end{enpratique}}

% ---------- Exercices, corrigés ----------
\newcounter{exo}[section]
\renewcommand{\theexo}{\thesection.\arabic{exo}}
\newcommand{\niveau}[1]{{\color{Ocre}\footnotesize
  \ifcase#1\or$\bullet\circ\circ\circ\circ$\or$\bullet\bullet\circ\circ\circ$\or$\bullet\bullet\bullet\circ\circ$%
  \or$\bullet\bullet\bullet\bullet\circ$\or$\bullet\bullet\bullet\bullet\bullet$\fi}}
\newcommand{\serie}[1]{\par\needspace{8\baselineskip}\vspace{1.5em}%
  \noindent{\sffamily\bfseries\color{Bleu}#1}\par\nobreak\vspace{-0.4em}%
  \noindent{\color{Filet}\rule{\linewidth}{0.4pt}}\par\nobreak\vspace{0.1em}}
\newenvironment{exo}[3]{%
  \par\needspace{5\baselineskip}\vspace{0.75em}\refstepcounter{exo}\label{e:#1}%
  \def\exolab{#1}%
  \noindent{\sffamily\bfseries\color{Bleu}Exercice~\theexo}\quad\niveau{#2}%
  \hfill{\footnotesize\sffamily\color{Gris}#3}\par\nobreak\vspace{0.15em}}%
  {\par\nopagebreak\noindent\hfill{\scriptsize\sffamily\hyperref[c:\exolab]{\color{Gris}corrigé p.~\pageref{c:\exolab}}}\par}
\newenvironment{corrige}[1]{%
  \par\needspace{6\baselineskip}\vspace{1em}\phantomsection\label{c:#1}%
  \noindent{\sffamily\bfseries\color{Bleu}Exercice~\ref{e:#1}}\quad
  {\scriptsize\sffamily\hyperref[e:#1]{\color{Gris}énoncé p.~\pageref{e:#1}}}\par\nobreak\vspace{0.2em}}%
  {\hfill\textcolor{BleuClair}{$\square$}\par}
\newcommand{\etape}[1]{\par\smallskip\noindent{\sffamily\bfseries\small\color{BleuMoyen}#1.}\ }
\newenvironment{solution}{\par\noindent{\sffamily\bfseries\small\color{BleuMoyen}Solution.}\ }{\hfill\textcolor{BleuClair}{$\square$}\par}
\newcommand{\chercher}{\par\smallskip\noindent{\color{Filet}\rule[0.5ex]{1.2cm}{0.4pt}}\
  {\footnotesize\sffamily\itshape\color{Gris}laisser chercher l'étudiant avant de poursuivre}\par\smallskip}

% ---------- Guide de séance ----------

\newcommand{\marge}[1]{\leavevmode\marginpar{\raggedleft\scriptsize\sffamily\color{Gris}\textls[100]{\MakeUppercase{#1}}}\ignorespaces}
\newcommand{\Etape}[4]{%
  \subsection[Étape #1 --- #3 {\color{Gris}(#2)}]{#3}\label{etape:#1}%
  \noindent{\small\sffamily\color{Gris}{\color{Ocre}\bfseries Étape #1}\quad·\quad #2\quad·\quad #4}\par\medskip}
\newcommand{\question}[1]{\par\smallskip\noindent{\sffamily\bfseries\small\color{BleuMoyen}Question.}\ \textit{#1}\par\smallskip}
\newcommand{\rappel}[1]{\par\smallskip\noindent\phantomsection\addcontentsline{rap}{rapentry}{#1}%
  {\sffamily\bfseries\scriptsize\colorbox{Fond}{\color{Gris}\strut RAPPEL}}\ {\sffamily\small\color{Gris}#1.}\ }
\newcommand{\pourquoi}{\par\noindent{\sffamily\bfseries\small\color{Ocre}Pourquoi ?}\ }

% Ouverture de chapitre : \ouverture{objectifs}{prérequis}{durée}
\newcommand{\ouverture}[3]{%
  \begin{tcolorbox}[enhanced,colback=Fond,colframe=Fond,arc=0pt,boxrule=0pt,left=12pt,right=12pt,top=10pt,bottom=10pt,
     before skip=1.2em,after skip=1.6em]
  \begin{minipage}[t]{0.56\linewidth}
  {\sffamily\bfseries\small\color{Bleu}OBJECTIFS}\par\vspace{0.3em}\small #1
  \end{minipage}\hfill
  \begin{minipage}[t]{0.38\linewidth}
  {\sffamily\bfseries\small\color{Bleu}PRÉREQUIS}\par\vspace{0.3em}\small #2
  \ifx\relax#3\relax\else\par\vspace{0.6em}{\sffamily\bfseries\small\color{Bleu}DURÉE}\par\vspace{0.2em}\small #3\fi
  \end{minipage}
  \end{tcolorbox}}
% Fin de chapitre
\newenvironment{essentiel}{\par\needspace{12\baselineskip}\vspace{1.5em}%
  \begin{tcolorbox}[enhanced,unbreakable,colback=white,colframe=white,boxrule=0pt,arc=0pt,
   borderline north={1.4pt}{0pt}{Ocre},left=2pt,right=2pt,top=8pt,bottom=4pt]
  {\sffamily\bfseries\color{Ocre}L'ESSENTIEL DU CHAPITRE}\par\medskip}{\end{tcolorbox}}

% ---------- Notations ----------
\newcommand{\K}{\mathbb{K}}
\newcommand{\R}{\mathbb{R}}
\newcommand{\C}{\mathbb{C}}
\newcommand{\N}{\mathbb{N}}
\newcommand{\Z}{\mathbb{Z}}
\newcommand{\Mn}[1]{\mathcal{M}_{n}(#1)}
\newcommand{\GLn}[1]{\mathrm{GL}_n(#1)}
\newcommand{\LE}{\mathcal{L}(E)}
\newcommand{\rg}{\operatorname{rg}}
\newcommand{\vect}{\operatorname{Vect}}
\newcommand{\Mat}{\operatorname{Mat}}
\newcommand{\Id}{\operatorname{Id}}
\newcommand{\sh}{\operatorname{sh}}
\newcommand{\ch}{\operatorname{ch}}
\renewcommand{\th}{\operatorname{th}}
\DeclareMathOperator{\argsh}{argsh}
\newcommand{\vv}[1]{\overrightarrow{#1}}
\newcommand{\ve}[1]{\vec{#1}}
\newcommand{\ient}[2]{\llbracket #1,#2\rrbracket}
\newcommand{\dt}[1]{\left|\begin{matrix}#1\end{matrix}\right|}
\newcommand{\dis}{\displaystyle}

% ---------- Figures ----------
\pgfplotsset{trig format plots=rad}
\pgfplotsset{courbe/.style={axis lines=middle,axis equal image,
  axis line style={Gris,-{Stealth[length=5pt]}},
  xlabel={$x$},ylabel={$y$},ticklabel style={font=\scriptsize\sffamily,color=Gris},
  every axis x label/.style={at={(ticklabel* cs:1)},anchor=west,font=\small},
  every axis y label/.style={at={(ticklabel* cs:1)},anchor=south,font=\small},
  samples=300,smooth,width=0.6\linewidth}}
\tikzset{arcetude/.style={Bleu,line width=1.25pt},
         reste/.style={BleuClair,line width=1pt},
         tgte/.style={Ocre,line width=0.9pt},
         asy/.style={Gris,dashed,line width=0.6pt},
         pt/.style={circle,fill=Ink,inner sep=1.2pt}}

% ---------- Navigation ----------
\makeatletter
\newcommand*\l@rapentry{\@dottedtocline{1}{0em}{0em}}
\newcommand*\l@metentry{\@dottedtocline{1}{0em}{0em}}
\newcommand{\listemethodes}{\@starttoc{met}}
\newcommand{\listerappels}{\@starttoc{rap}}
\makeatother
% onglets de chapitre sur le bord droit
\usepackage{eso-pic}
\newif\ifonglets \ongletsfalse
\newcommand{\dessineronglet}{%
  \edef\ongcol{\ongletcouleur}%
  \edef\ongnum{\grandnum}%
  \pgfmathsetmacro{\ongy}{3.2+2.3*(\ongnum-1)}%
  \begin{tikzpicture}[remember picture,overlay]
  \fill[\ongcol] ([xshift=-0.75cm,yshift=-\ongy cm]current page.north east)
     rectangle ([yshift=-\ongy cm-1.9cm]current page.north east);
  \node[text=white,font=\sffamily\bfseries\large] at ([xshift=-0.37cm,yshift=-\ongy cm-0.95cm]current page.north east)
     {\ongnum};
  \end{tikzpicture}}
\AddToShipoutPictureBG{\ifonglets\ifnum\value{section}>0 \dessineronglet\fi\fi}
% blocs de navigation et d'enseignement
\definirbloc{commencer}{Gris}{Si vous commencez ici}
\definirbloc{horsprogramme}{Brique}{Hors programme}
\newcommand{\HP}{\,{\sffamily\bfseries\scriptsize\textcolor{Brique}{[HP]}}}
\newcommand{\source}[1]{{\footnotesize\sffamily\color{Gris}\,(#1)}}
\newcommand{\coupdepouce}[1]{\par\vspace{-0.2em}\noindent{\footnotesize\sffamily\color{Gris}%
  \textbf{Coup de pouce} (à donner si l'étudiant bloque) : #1}\par}
\newcommand{\pointarret}[1]{\par\vspace{1em}\noindent{\color{Filet}\rule[0.5ex]{\linewidth}{0.4pt}}\par\nobreak
  \noindent{\footnotesize\sffamily\color{Ocre}\bfseries POINT D'ARRÊT POSSIBLE}\quad{\footnotesize\sffamily\color{Gris}#1}\par\vspace{0.6em}}
\newcommand{\demarche}[1]{\noindent{\sffamily\bfseries\small\color{BleuMoyen}Démarche.}\ \textit{#1}\par\smallskip}

% ---------- Corrigés : idées, astuces, pièges, à retenir ----------
\newcommand{\choix}[1]{\par\noindent{\sffamily\bfseries\small\color{BleuMoyen}Choix de la méthode.}\ #1\par\smallskip}
\newcommand{\gain}[1]{\par\noindent{\sffamily\bfseries\small\color{Ocre}Astuce.}\ #1\par}
\newcommand{\piege}[1]{\par\noindent{\sffamily\bfseries\small\color{Brique}Piège.}\ #1\par}
\newtcolorbox{retenirbox}{enhanced,breakable,colback=Fond,colframe=Fond,arc=0pt,boxrule=0pt,
  borderline west={2pt}{0pt}{Ocre},left=8pt,right=6pt,top=4pt,bottom=4pt,before skip=0.7em,after skip=0.5em}
\newcommand{\retenir}[1]{\begin{retenirbox}{\sffamily\bfseries\footnotesize\color{Ocre}À RETENIR}\quad #1\end{retenirbox}}
\newcommand{\TD}[1]{{\color{Ocre}\bfseries TD\,#1}\ ·\ }
\newcommand{\q}[1]{\par\smallskip\noindent{\sffamily\bfseries\color{Bleu}#1}\ }
\providecommand{\iddots}{\mathinner{\mkern1mu\raise1pt\hbox{.}\mkern2mu\raise4pt\hbox{.}\mkern2mu\raise7pt\hbox{.}\mkern1mu}}


% =====================================================================
%  Examen blanc — Algèbre linéaire et déterminants (Le Collé, PT)
%  Une source, deux PDF : \def\CORRIGE{} avant \input donne le corrigé.
% =====================================================================
\providecommand{\Ker}{\operatorname{Ker}}\renewcommand{\Im}{\operatorname{Im}}
\fancyfoot[L]{}\fancyhead[L]{\footnotesize\sffamily\color{Gris}Examen blanc · Algèbre linéaire et déterminants\ifdefined\CORRIGE\ · corrigé\fi}
\newcommand{\pts}[1]{\hfill{\footnotesize\sffamily\color{Gris}(#1)}}
\newcommand{\Ex}[3]{\par\needspace{6\baselineskip}\vspace{1.3em}\noindent
  {\sffamily\bfseries\large\color{Bleu}Exercice #1 \,·\, #2}\hfill{\sffamily\small\color{Ocre}#3}\par
  \vspace{2pt}{\color{Filet}\rule{\linewidth}{0.4pt}}\par\smallskip}
\newcommand{\Pb}[2]{\par\needspace{6\baselineskip}\vspace{1.3em}\noindent
  {\sffamily\bfseries\large\color{Bleu}Problème \,·\, #1}\hfill{\sffamily\small\color{Ocre}#2}\par
  \vspace{2pt}{\color{Filet}\rule{\linewidth}{0.4pt}}\par\smallskip}
\newcommand{\Partie}[1]{\par\needspace{4\baselineskip}\medskip\noindent{\sffamily\bfseries\color{BleuMoyen}#1}\par\nobreak\smallskip}
\newcommand{\Q}[1]{\par\needspace{3\baselineskip}\medskip\noindent{\sffamily\bfseries\color{Bleu}#1}\enspace}
\newcommand{\autre}[1]{\par\noindent{\sffamily\bfseries\small\color{BleuMoyen}Autre méthode.}\ #1\par}
\newcommand{\lienDL}[1]{\par\noindent{\sffamily\bfseries\small\color{Ocre}Lien avec le DL.}\ #1\par}
\setlist[enumerate,1]{label=\arabic*.,leftmargin=1.6em,itemsep=4pt,topsep=4pt}
\setlist[enumerate,2]{label=(\alph*),leftmargin=1.6em,itemsep=2pt,topsep=2pt}

\begin{document}
\pagestyle{fancy}
\thispagestyle{plain}

% ------------------------------------------------------------------ en-tête
\begin{tikzpicture}[remember picture,overlay]
\fill[Bleu] (current page.north west) rectangle ([yshift=-4.2cm]current page.north east);
\node[anchor=north west,text=white,font=\sffamily\footnotesize,inner sep=0] at ([xshift=2.2cm,yshift=-1.1cm]current page.north west)
  {\textls[120]{MATHÉMATIQUES · CPGE PT · EXAMEN BLANC}};
\node[anchor=north west,text=white,font=\sffamily\bfseries\LARGE,inner sep=0] at ([xshift=2.2cm,yshift=-1.75cm]current page.north west)
  {Algèbre linéaire et déterminants};
\node[anchor=north west,text=white,font=\sffamily\normalsize,inner sep=0] at ([xshift=2.2cm,yshift=-2.75cm]current page.north west)
  {\ifdefined\CORRIGE Corrigé détaillé et diagnostic\else Durée : 3 heures \quad·\quad Calculatrice interdite\fi};
\node[anchor=north east,text=white,font=\sffamily\scshape\large,inner sep=0] at ([xshift=-2.2cm,yshift=-1.05cm]current page.north east)
  {Le Collé};
\end{tikzpicture}
\vspace*{2.9cm}

\ifdefined\CORRIGE
% =====================================================================
%                              CORRIGÉ
% =====================================================================
\markboth{Corrigé · Examen blanc}{}
\paragraph{Comment utiliser ce corrigé} Ne l'ouvre qu'après avoir composé les 3 heures. Corrige d'abord ta copie avec un
stylo d'une autre couleur, question par question, en comptant les points. Pour chaque question ratée, lis le \emph{choix de la
méthode} : c'est lui qu'il faut retenir, plus que le calcul. La grille de la dernière page dit quoi revoir.

\subsection*{Ce que le sujet évalue}
\begin{center}
\small\renewcommand{\arraystretch}{1.22}
\begin{tabular}{@{}>{\raggedright\arraybackslash}p{7.6cm}>{\raggedright\arraybackslash}p{5.4cm}@{}}
\toprule
\textbf{Compétence} & \textbf{Questions}\\\midrule
Démontrer qu'une famille est libre (combinaison nulle, évaluation, composition) & 1.2 · 2.3(b) · 2.4(a) · P.3\\
Démontrer qu'une famille est génératrice, trouver une base d'un sous-espace & 1.3 · 2.2(a) · 3.2\\
Utiliser la dimension (« $n$ vecteurs en dimension $n$ ») & 1.1 · 2.2(b) · 2.4(a) · 3.3 · P.3\\
Coordonnées dans une base, changement de base & 2.1(b) · 2.3(c) · 2.4(c) · 3.4 · P.4\\
Sous-espaces, somme directe, supplémentaires, projecteur & 1.4 · 3.1 · 3.3 · 3.4\\
Application linéaire : noyau, image, rang, théorème du rang & 1.6 · P.2\\
Calculer un déterminant (opérations, développement, factorisation) & 2.1(a) · 4.1 · 4.2\\
Propriétés du déterminant ($\det(\lambda A)$, $\det A^T$, $\det(AB)$) & 1.5 · 4.3 · P.6\\
Déterminant, inversibilité, liberté, systèmes de Cramer & 2.1 · 4.1 · 4.2 · P.6\\
« Trouver tous les… » : prouver les deux sens & 2.1(a) · P.5\\\bottomrule
\end{tabular}
\end{center}

\lienDL{Les questions 2.3(b), 4.2 et P.3 à P.5 reprennent les difficultés du DL~3 sous une autre forme : liberté par évaluation
(comme III.3), Vandermonde, famille $(x,Ax,A^2x)$ avec $A^3=0$ (comme II.1 et II.4.a). Les formulations sont différentes : il
faut reconnaître la technique.}

\begin{center}\small\renewcommand{\arraystretch}{1.2}
\begin{tabular}{@{}lcccccc@{}}\toprule
& Ex.~1 & Ex.~2 & Ex.~3 & Ex.~4 & Problème & Total\\\midrule
Points & 3 & 5 & 3,5 & 4 & 4,5 & 20\\
Temps conseillé & 15 min & 50 min & 35 min & 35 min & 45 min & 3 h\\\bottomrule
\end{tabular}
\end{center}

% ---------------------------------------------------------------------
\clearpage
\Ex{1}{Vrai ou faux}{3 points}

\Q{1. Vrai.} $\dim\R^3=3$ : une famille libre de $\R^3$ a au plus 3 vecteurs. Une famille de 4 vecteurs est donc liée.

\Q{2. Faux.} Contre-exemple : $u=(1,0,0)$, $v=(0,1,0)$, $w=(1,0,0)$. Les familles $(u,v)$ et $(v,w)$ sont libres, mais $u-w=0$ :
$(u,v,w)$ est liée.
\piege{la liberté ne se vérifie pas « deux à deux ». Trois vecteurs deux à deux non colinéaires peuvent être liés, par exemple
$(1,0)$, $(0,1)$, $(1,1)$ dans $\R^2$.}

\Q{3. Vrai.} $\dim\R^3=3$ et la famille a 3 vecteurs : pour une telle famille, génératrice $\iff$ libre $\iff$ base.

\Q{4. Faux.} Dans $\R^2$, $F=\vect\big((1,0)\big)$ et $G=\vect\big((0,1)\big)$ : $(1,0)$ et $(0,1)$ sont dans $F\cup G$, mais pas
leur somme $(1,1)$. $F\cup G$ n'est pas stable par addition.

\Q{5. Faux.} Multiplier $A$ par $2$ multiplie chacune des 3 colonnes par $2$, et le déterminant est linéaire par rapport à
\emph{chaque} colonne : $\det(2A)=2^3\det A=8\det A$. Par exemple $\det(2I_3)=8$.
\retenir{$\det(\lambda A)=\lambda^n\det A$ pour $A\in\Mn{\K}$.}

\Q{6. Vrai.} Par le théorème du rang, $\dim\Ker f=4-\rg f$. Or $\Im f\subset\R^3$, donc $\rg f\leq3$ et $\dim\Ker f\geq1$ :
$\Ker f\neq\{0\}$, $f$ n'est pas injective.

% ---------------------------------------------------------------------
\Ex{2}{Familles libres, génératrices, bases}{5 points}

\Partie{A. Une famille de $\R^3$ avec paramètre}

\Q{2.1(a)} \choix{3 vecteurs en dimension 3 avec un paramètre : un seul déterminant donne toutes les valeurs de $m$ d'un coup.
Résoudre un système avec $m$ obligerait à discuter à chaque pivot.}
\[\det(u_1,u_2,u_3)=\begin{vmatrix}1&0&2\\2&1&m\\-1&1&1\end{vmatrix}
\underset{C_3\leftarrow C_3-2C_1}{=}\begin{vmatrix}1&0&0\\2&1&m-4\\-1&1&3\end{vmatrix}=1\times\begin{vmatrix}1&m-4\\1&3\end{vmatrix}=7-m.\]
Le déterminant est non nul si et seulement si $m\neq7$. \textbf{$(u_1,u_2,u_3)$ est une base de $\R^3$ si et seulement si $m\neq7$.}
Les deux sens sont prouvés en même temps, car « base $\iff\det\neq0$ ».

\Q{2.1(b)} Pour $m=1$, on cherche $(a,b,c)$ tel que $au_1+bu_2+cu_3=(1,1,1)$ :
\[\begin{cases}a+2c=1\\2a+b+c=1\\-a+b+c=1\end{cases}\]
La différence des deux dernières lignes donne $3a=0$, donc $a=0$, puis $c=\frac12$ et $b=\frac12$. \textbf{Les coordonnées de
$(1,1,1)$ dans $\mathcal B'$ sont $\left(0,\frac12,\frac12\right)$.} Vérification : $\frac12(u_2+u_3)=\frac12(2,2,2)=(1,1,1)$.
\piege{les coordonnées dépendent de la base : $(1,1,1)$ est la colonne dans la base canonique, $(0,\frac12,\frac12)$ dans
$\mathcal B'$. Toujours vérifier en recombinant.}

\Q{2.1(c)} Pour $m=7$, on cherche une relation $au_1+bu_2+u_3=0$ : la première coordonnée donne $a=-2$, la troisième $b=-3$, et
la deuxième est alors vérifiée. Donc $u_3=2u_1+3u_2$. Comme $u_1$ et $u_2$ ne sont pas colinéaires,
$\rg(u_1,u_2,u_3)=\dim\vect(u_1,u_2)=2$.

\Partie{B. Un sous-espace de $\R^4$}

\Q{2.2(a)} \choix{un sous-espace défini par des équations : on résout le système, et on écrit la solution générale comme
combinaison linéaire. Cela donne directement une famille génératrice.}

$(x,y,z,t)\in F\iff x=y-2z$ et $t=-x-y-z=-2y+z$. Donc
\[F=\{\,y(1,1,0,-2)+z(-2,0,1,1),\ y,z\in\R\,\}=\vect(f_1,f_2),\]
avec $f_1=(1,1,0,-2)$ et $f_2=(-2,0,1,1)$.
\emph{Génératrice :} c'est l'égalité ci-dessus. \emph{Libre :} si $yf_1+zf_2=0$, les deuxième et troisième coordonnées donnent
$y=0$ et $z=0$. \textbf{$(f_1,f_2)$ est une base de $F$, et $\dim F=2$.}
\autre{$F$ est défini par 2 équations indépendantes dans $\R^4$, donc $\dim F=4-2=2$ (théorème du rang appliqué à
$(x,y,z,t)\mapsto(x+y+z+t,\,x-y+2z)$, de rang 2).}

\Q{2.2(b)} \choix{on connaît maintenant $\dim F=2$. Pour 2 vecteurs dans un espace de dimension 2, il suffit qu'ils soient
dans $F$ et libres : inutile de prouver qu'ils engendrent $F$.}

$w_1=(3,1,-1,-3)$ : $3+1-1-3=0$ et $3-1-2=0$, donc $w_1\in F$. $w_2=(0,2,1,-3)$ : $0+2+1-3=0$ et $0-2+2=0$, donc $w_2\in F$. Ils
ne sont pas colinéaires (la première coordonnée de $w_2$ est nulle, pas celle de $w_1$, et $w_2\neq0$). \textbf{$(w_1,w_2)$ est une
famille libre de 2 vecteurs de $F$, avec $\dim F=2$ : c'est une base de $F$.}
\piege{oublier de vérifier que $w_1,w_2\in F$. Deux vecteurs libres de $\R^4$ ne forment pas une base de $F$ s'ils n'y sont pas.}

\Partie{C. Une famille de polynômes}

\Q{2.3(a)}
\begin{center}\small\renewcommand{\arraystretch}{1.2}
\begin{tabular}{@{}cccc@{}}\toprule
& $x=0$ & $x=1$ & $x=2$\\\midrule
$P_1=(X-1)(X-2)$ & $2$ & $0$ & $0$\\
$P_2=X(X-2)$ & $0$ & $-1$ & $0$\\
$P_3=X(X-1)$ & $0$ & $0$ & $2$\\\bottomrule
\end{tabular}
\end{center}

\Q{2.3(b)} \choix{les trois polynômes ont le même degré : l'argument « degrés distincts » ne marche pas. Le tableau montre
que chaque $P_i$ ne vit qu'en un point. On évalue la relation en $0$, $1$, $2$ : chaque évaluation isole un coefficient.}

Soient $a,b,c$ tels que $aP_1+bP_2+cP_3=0$. On évalue en $0$ : $2a=0$. En $1$ : $-b=0$. En $2$ : $2c=0$. Donc $a=b=c=0$ : la
famille est libre. Elle a 3 vecteurs et $\dim\R_2[X]=3$ : \textbf{c'est une base de $\R_2[X]$}.
\lienDL{c'est la même idée qu'en III.3 : une relation entre vecteurs devient une information sur des valeurs, qu'on lit point
par point.}

\Q{2.3(c)} D'après (b), $P=aP_1+bP_2+cP_3$ pour des réels $a,b,c$. On évalue en $0$, $1$, $2$ :
$P(0)=2a$, $P(1)=-b$, $P(2)=2c$. Donc
\[P=\frac{P(0)}2\,P_1-P(1)\,P_2+\frac{P(2)}2\,P_3.\]

\Q{2.3(d)} D'après (c), si $P(0)=1$, $P(1)=0$, $P(2)=3$, nécessairement $P=\frac12P_1+\frac32P_3=2X^2-3X+1$. Réciproquement,
ce polynôme vérifie $P(0)=1$, $P(1)=0$, $P(2)=8-6+1=3$. \textbf{L'unique solution est $P=2X^2-3X+1$.}

\Partie{D. Un raisonnement sans coordonnées}

\Q{2.4(a)} \choix{on ne connaît pas les vecteurs, seulement que $(u,v,w)$ est une base. On écrit une combinaison nulle et on la
réécrit en fonction de $u,v,w$ pour utiliser leur liberté.}

Soient $\alpha,\beta,\gamma$ tels que $\alpha(u+v)+\beta(v+w)+\gamma(w+u)=0$, soit $(\alpha+\gamma)u+(\alpha+\beta)v+(\beta+\gamma)w=0$.
Comme $(u,v,w)$ est libre : $\alpha+\gamma=\alpha+\beta=\beta+\gamma=0$. En additionnant, $\alpha+\beta+\gamma=0$, puis
$\beta=\gamma=\alpha=0$. La famille est libre, a 3 vecteurs, et $\dim E=3$ : \textbf{c'est une base de $E$.}
\autre{la matrice de la famille dans la base $(u,v,w)$ est $\begin{pmatrix}1&0&1\\1&1&0\\0&1&1\end{pmatrix}$, de déterminant
$2\neq0$.}

\Q{2.4(b)} $(u-v)+(v-w)+(w-u)=0$ : c'est une relation à coefficients non tous nuls. \textbf{La famille est liée.}

\Q{2.4(c)} $(u+v)-(v+w)+(w+u)=2u$, donc $u=\frac12(u+v)-\frac12(v+w)+\frac12(w+u)$ : \textbf{coordonnées
$\left(\frac12,-\frac12,\frac12\right)$}.

% ---------------------------------------------------------------------
\Ex{3}{Sous-espaces supplémentaires et projecteur}{3,5 points}

\Q{3.1} L'application $\varphi:P\mapsto\big(P(1),P'(1)\big)$, de $E$ dans $\R^2$, est linéaire (l'évaluation et la dérivation le
sont). $F=\Ker\varphi$ est donc un sous-espace vectoriel de $E$.
\autre{vérifier directement : le polynôme nul est dans $F$, et si $P,Q\in F$, alors $(\lambda P+\mu Q)(1)=0$ et
$(\lambda P+\mu Q)'(1)=0$.}

\Q{3.2(a)} Les polynômes $1$, $X-1$, $(X-1)^2$, $(X-1)^3$ sont non nuls et de degrés $0,1,2,3$ distincts : la famille est libre.
Elle a 4 vecteurs et $\dim E=4$ : \textbf{$\mathcal B_1$ est une base de $E$}.

\Q{3.2(b)} \choix{la formule de Taylor donne directement les coordonnées dans $\mathcal B_1$, et les conditions qui
définissent $F$ portent exactement sur les deux premières.}

Par la formule de Taylor, les coordonnées de $P$ dans $\mathcal B_1$ sont $\left(P(1),\,P'(1),\,\frac{P''(1)}2,\,\frac{P'''(1)}6\right)$.
Donc $P\in F$ si et seulement si ses deux premières coordonnées sont nulles, c'est-à-dire si et seulement si
$P\in\vect\big((X-1)^2,(X-1)^3\big)$. La famille $\big((X-1)^2,(X-1)^3\big)$ engendre $F$, et elle est libre car extraite de la
base $\mathcal B_1$. \textbf{C'est une base de $F$, et $\dim F=2$.}

\Q{3.3} \choix{deux méthodes possibles. La plus rapide : intersection nulle et somme des dimensions.}

Soit $P\in F\cap G$. Comme $P\in G=\R_1[X]$, $P=a+b(X-1)$, donc $P(1)=a$ et $P'(1)=b$. Comme $P\in F$, $a=b=0$ : $P=0$. Ainsi
$F\cap G=\{0\}$, et $\dim F+\dim G=2+2=4=\dim E$. \textbf{Donc $E=F\oplus G$.}
\autre{$G=\vect\big(1,X-1\big)$ et $F=\vect\big((X-1)^2,(X-1)^3\big)$ : en réunissant ces deux bases, on obtient la base
$\mathcal B_1$ de $E$, donc $E=F\oplus G$.}

\Q{3.4(a)} D'après la formule de Taylor,
\[P=\underbrace{P(1)+P'(1)(X-1)}_{\in\,G}+\underbrace{\frac{P''(1)}2(X-1)^2+\frac{P'''(1)}6(X-1)^3}_{\in\,F}.\]
La décomposition de $P$ selon $E=F\oplus G$ est unique ; la composante dans $G$ est $q(P)$. \textbf{Donc
$q(P)=P(1)+P'(1)(X-1)$.} Pour $P=X^3$ : $q(X^3)=1+3(X-1)=3X-2$. Contrôle : $X^3-(3X-2)=(X-1)^2(X+2)$ est bien dans $F$.

\Q{3.4(b)} $q$ laisse fixes les vecteurs de $G$ et annule ceux de $F$ :
\[\Mat_{\mathcal B_1}(q)=\begin{pmatrix}1&0&0&0\\0&1&0&0\\0&0&0&0\\0&0&0&0\end{pmatrix}.\]
Dans $\mathcal B_0$ : $q(1)=1$, $q(X)=X$, $q(X^2)=1+2(X-1)=2X-1$, $q(X^3)=3X-2$, donc
\[\Mat_{\mathcal B_0}(q)=\begin{pmatrix}1&0&-1&-2\\0&1&2&3\\0&0&0&0\\0&0&0&0\end{pmatrix}.\]
On vérifie que cette matrice $Q$ vérifie $Q^2=Q$, comme tout projecteur.
\retenir{dans une base adaptée à $E=F\oplus G$, la matrice d'un projecteur est diagonale avec des $1$ et des $0$. Choisir la
bonne base simplifie tout : c'est l'idée de toute la suite du cours.}

% ---------------------------------------------------------------------
\Ex{4}{Calculs de déterminants}{4 points}

\Q{4.1} \choix{la somme de chaque ligne vaut $a+3b$ : on additionne toutes les colonnes dans la première pour factoriser.}
\begin{align*}D_4&=\begin{vmatrix}a&b&b&b\\b&a&b&b\\b&b&a&b\\b&b&b&a\end{vmatrix}
\underset{C_1\leftarrow C_1+C_2+C_3+C_4}{=}(a+3b)\begin{vmatrix}1&b&b&b\\1&a&b&b\\1&b&a&b\\1&b&b&a\end{vmatrix}\\
&\underset{L_i\leftarrow L_i-L_1}{=}(a+3b)\begin{vmatrix}1&b&b&b\\0&a-b&0&0\\0&0&a-b&0\\0&0&0&a-b\end{vmatrix}.\end{align*}
La dernière matrice est triangulaire : \textbf{$D_4=(a+3b)(a-b)^3$}. La matrice est inversible si et seulement si $D_4\neq0$,
c'est-à-dire \textbf{$a\neq b$ et $a\neq-3b$}.
\piege{développer directement un $4\times4$ : 24 termes et presque sûrement une erreur. Toujours chercher d'abord une
structure (somme constante, colonnes proches, zéros).}

\Q{4.2(a)} On fait $L_2\leftarrow L_2-L_1$ et $L_3\leftarrow L_3-L_1$, puis on factorise :
\[V=\begin{vmatrix}1&a&a^2\\0&b-a&b^2-a^2\\0&c-a&c^2-a^2\end{vmatrix}=(b-a)(c-a)\begin{vmatrix}1&b+a\\1&c+a\end{vmatrix}=(b-a)(c-a)(c-b).\]
\gain{$b^2-a^2=(b-a)(b+a)$ : on fait apparaître les facteurs $(b-a)$ et $(c-a)$ avant de développer.}

\Q{4.2(b)} $P=x_0+x_1X+x_2X^2$ vérifie les trois conditions si et seulement si
\[\begin{pmatrix}1&a&a^2\\1&b&b^2\\1&c&c^2\end{pmatrix}\begin{pmatrix}x_0\\x_1\\x_2\end{pmatrix}=\begin{pmatrix}\alpha\\\beta\\\gamma\end{pmatrix}.\]
Le déterminant de ce système est $V\neq0$ car $a,b,c$ sont distincts : c'est un système de Cramer, il a une unique solution.
\textbf{Il existe un unique $P\in\R_2[X]$ qui convient.}
\lienDL{c'est le même phénomène qu'en 2.3, où $a=0$, $b=1$, $c=2$, et qu'en III.3 du DL : trois valeurs distinctes
« séparent » les polynômes de degré au plus 2. Ici on le prouve par le déterminant, en 2.3 par évaluation.}

\Q{4.3} \choix{on ne peut rien calculer : on utilise les propriétés, $\det A^T=\det A$ et $\det(\lambda A)=\lambda^n\det A$.}

$A^T=-A$, donc $\det A=\det A^T=\det(-A)=(-1)^n\det A=-\det A$ car $n$ est impair. Ainsi $2\det A=0$ : \textbf{$\det A=0$, et
$A$ n'est pas inversible.} Pour $n$ pair, c'est faux : $\begin{pmatrix}0&1\\-1&0\end{pmatrix}$ est antisymétrique, de
déterminant $1$.
\piege{écrire $\det(-A)=-\det A$ sans justifier : c'est vrai seulement parce que $n$ est impair.}

% ---------------------------------------------------------------------
\Pb{une matrice nilpotente}{4,5 points}

\Q{P.1}
\[A^2=\begin{pmatrix}0&0&0\\1&1&1\\-1&-1&-1\end{pmatrix},\qquad A^3=A^2A=0.\]
(Chaque ligne de $A^2$ est proportionnelle à $(1,1,1)$, et $(1,1,1)A=(0,0,0)$.)

\Q{P.2} \emph{Noyau.} $AX=0$ s'écrit $-x-y-z=0$, $x+2y+2z=0$, $-y-z=0$. La dernière équation donne $z=-y$, puis la première
$x=0$ ; la deuxième est alors vérifiée. \textbf{$\Ker A=\vect\big((0,1,-1)\big)$}, de dimension 1.

\emph{Image.} Elle est engendrée par les colonnes $C_1=(-1,1,0)$, $C_2=(-1,2,-1)$, $C_3=C_2$. $C_1$ et $C_2$ ne sont pas
colinéaires : \textbf{$\Im A=\vect(C_1,C_2)$, et $\rg A=2$}. Théorème du rang : $1+2=3$. \checkmark

\emph{Équation.} $C_1$ et $C_2$ vérifient $x+y+z=0$, et ce plan est de dimension 2 : \textbf{$\Im A$ a pour équation $x+y+z=0$}.

\emph{Inclusion.} $(0,1,-1)$ vérifie $0+1-1=0$, donc $\Ker A\subset\Im A$.

\Q{P.3} \choix{l'énoncé interdit le déterminant : on utilise $A^3=0$. Appliquer $A^2$ à une relation fait disparaître tous
les termes sauf un.}

On a $x_0=e_1$, $Ax_0=(-1,1,0)$, $A^2x_0=(0,1,-1)\neq0$. Soient $c_0,c_1,c_2$ tels que $c_0x_0+c_1Ax_0+c_2A^2x_0=0$.
\begin{itemize}[leftmargin=1.2em,itemsep=1pt]
\item On multiplie par $A^2$ : comme $A^3=A^4=0$, il reste $c_0A^2x_0=0$, et $A^2x_0\neq0$, donc $c_0=0$.
\item On multiplie alors par $A$ : il reste $c_1A^2x_0=0$, donc $c_1=0$.
\item Il reste $c_2A^2x_0=0$, donc $c_2=0$.
\end{itemize}
La famille est libre, de 3 vecteurs en dimension 3 : \textbf{$\mathcal B$ est une base de $\R^3$.}
\lienDL{c'est exactement la question II.1 du DL, avec une autre matrice.}

\Q{P.4} $f(x_0)=Ax_0$ est le 2\textsuperscript{e} vecteur de $\mathcal B$, $f(Ax_0)=A^2x_0$ le 3\textsuperscript{e}, et
$f(A^2x_0)=A^3x_0=0$ :
\[N=\Mat_{\mathcal B}(f)=\begin{pmatrix}0&0&0\\1&0&0\\0&1&0\end{pmatrix},\qquad
P=\begin{pmatrix}1&-1&0\\0&1&1\\0&0&-1\end{pmatrix},\qquad A=PNP^{-1}.\]
$P$ a pour colonnes les coordonnées de $x_0$, $Ax_0$, $A^2x_0$ dans la base canonique. La relation vient de la formule de
changement de base $N=P^{-1}AP$.

\Q{P.5} \choix{« tous les $x$ » : deux sens, et la question P.3 fait le premier.}

Soit $x=(x_1,x_2,x_3)$. On a $A^2x=(x_1+x_2+x_3)(0,1,-1)$.
\begin{itemize}[leftmargin=1.2em,itemsep=1pt]
\item Si $x_1+x_2+x_3\neq0$, alors $A^2x\neq0$ et le raisonnement de P.3 s'applique mot pour mot : la famille est une base.
\item Si $x_1+x_2+x_3=0$, alors $A^2x=0$ : la famille contient le vecteur nul, elle est liée.
\end{itemize}
\textbf{$(x,Ax,A^2x)$ est une base si et seulement si $x_1+x_2+x_3\neq0$}, c'est-à-dire $x\notin\Ker A^2$.
\lienDL{c'est la question II.4.a du DL. La condition ne se devine pas : elle vient du calcul de $A^2x$.}

\Q{P.6} \choix{calculer directement $\det(A+\lambda I_3)$ marche, mais la relation $A=PNP^{-1}$ donne le résultat sans calcul :
deux matrices semblables ont le même déterminant.}

$A+\lambda I_3=PNP^{-1}+\lambda PP^{-1}=P(N+\lambda I_3)P^{-1}$, donc
\[\det(A+\lambda I_3)=\det P\,\det(N+\lambda I_3)\,\det P^{-1}=\det(N+\lambda I_3)=\lambda^3,\]
car $N+\lambda I_3$ est triangulaire, de diagonale $(\lambda,\lambda,\lambda)$. Avec $\lambda=1$ : $\det(A+I_3)=1\neq0$,
\textbf{$A+I_3$ est inversible}. Avec $\lambda=0$ : $\det A=0$, \textbf{$A$ n'est pas inversible} (on le savait : $\rg A=2$).
\gain{on peut même donner l'inverse : $(I_3+A)(I_3-A+A^2)=I_3+A^3=I_3$.}

% ---------------------------------------------------------------------
\clearpage
\markboth{Diagnostic}{}
\subsection*{Après la correction : quoi revoir ?}
\begin{center}
\small\renewcommand{\arraystretch}{1.25}
\begin{tabular}{@{}>{\raggedright\arraybackslash}p{5.6cm}>{\raggedright\arraybackslash}p{7.4cm}@{}}
\toprule
\textbf{Si tu as perdu des points en…} & \textbf{Revoir}\\\midrule
1.2, 2.3(b), 2.4(a), P.3 (liberté) & Livret d'algèbre linéaire, méthode R.2 ; livret du DL, outils 1 et 2\\
2.2, 3.2 (génératrice, base d'un sous-espace) & Livret d'algèbre linéaire, rappels « Espaces vectoriels » et « Dimension finie »\\
2.1(b), 2.4(c), 3.4, P.4 (coordonnées, matrices) & Livret d'algèbre linéaire, proposition « changement de base » ; livret du DL, outil 3\\
3.3, 3.4 (supplémentaires, projecteur) & Livret d'algèbre linéaire, §1.2\\
P.2 (noyau, image, rang) & Livret d'algèbre linéaire, méthode R.3 « Calculer un rang, un noyau, une image »\\
4.1, 4.2 (calculs) & Livret des déterminants, §1.2 et §1.5 (somme constante, Vandermonde)\\
1.5, 4.3, P.6 (propriétés) & Livret des déterminants, §1.3\\
2.1(a), P.5 (« tous les… ») & Livret du DL, outil 4\\\bottomrule
\end{tabular}
\end{center}

\begin{essentiel}
\begin{itemize}[leftmargin=1.1em,itemsep=3pt]
\item Libre : combinaison nulle, puis le bon geste pour isoler les coefficients (identifier, évaluer, composer par $A^k$).
\item $n$ vecteurs en dimension $n$ : une seule propriété à prouver. Pour une base d'un sous-espace $F$, vérifier aussi que les
vecteurs sont dans $F$.
\item Un paramètre et $n$ vecteurs en dimension $n$ : le déterminant donne toutes les valeurs d'un coup.
\item Un déterminant se calcule par opérations sur les lignes et colonnes avant de développer.
\item $\det(\lambda A)=\lambda^n\det A$, $\det A^T=\det A$, deux matrices semblables ont le même déterminant.
\end{itemize}
\end{essentiel}

\else
% =====================================================================
%                               SUJET
% =====================================================================
\markboth{Examen blanc · Algèbre linéaire et déterminants}{}
\begin{center}
\begin{tcolorbox}[enhanced,colback=Fond,colframe=Fond,arc=0pt,boxrule=0pt,left=10pt,right=10pt,top=8pt,bottom=8pt,width=\linewidth]
\small
\textbf{\sffamily\color{Bleu}Consignes.} Le sujet comporte quatre exercices et un problème, indépendants : on peut les traiter
dans l'ordre de son choix. Toute réponse doit être \emph{justifiée} : citer le résultat du cours utilisé et vérifier ses
hypothèses. Les résultats sont à encadrer. Une réponse juste sans justification ne rapporte pas de points ; une démarche correcte
avec une erreur de calcul en rapporte. La clarté et la précision de la rédaction sont prises en compte.

\medskip
\begin{tabular}{@{}lcccccc@{}}
\textbf{\sffamily\color{Bleu}Barème} & Ex.~1 : 3 & Ex.~2 : 5 & Ex.~3 : 3,5 & Ex.~4 : 4 & Problème : 4,5 & \textbf{Total : 20}\\
\end{tabular}
\end{tcolorbox}
\end{center}

\Ex{1}{Vrai ou faux}{3 points}
Pour chaque affirmation, dire si elle est vraie ou fausse, et le justifier : une preuve courte si elle est vraie, un
contre-exemple si elle est fausse. \pts{0,5 point par affirmation}
\begin{enumerate}
\item Toute famille de quatre vecteurs de $\R^3$ est liée.
\item Si $(u,v)$ et $(v,w)$ sont deux familles libres d'un espace vectoriel $E$, alors $(u,v,w)$ est libre.
\item Une famille de trois vecteurs qui engendre $\R^3$ est une base de $\R^3$.
\item La réunion de deux sous-espaces vectoriels de $\R^2$ est un sous-espace vectoriel de $\R^2$.
\item Pour toute matrice $A\in\mathcal M_3(\R)$, $\det(2A)=2\det A$.
\item Une application linéaire de $\R^4$ dans $\R^3$ n'est jamais injective.
\end{enumerate}

\Ex{2}{Familles libres, génératrices, bases}{5 points}

\Partie{A. Une famille de $\R^3$ avec paramètre}
Pour $m\in\R$, on pose $u_1=(1,2,-1)$, $u_2=(0,1,1)$ et $u_3=(2,m,1)$.
\begin{enumerate}
\item
\begin{enumerate}
\item Déterminer les valeurs de $m$ pour lesquelles $(u_1,u_2,u_3)$ est une base de $\R^3$. \pts{0,75 pt}
\item On prend $m=1$ et on note $\mathcal B'=(u_1,u_2,u_3)$. Déterminer les coordonnées du vecteur $(1,1,1)$ dans $\mathcal B'$.
\pts{0,5 pt}
\item Pour la valeur de $m$ exclue au (a), écrire $u_3$ comme combinaison linéaire de $u_1$ et $u_2$, et donner le rang de
$(u_1,u_2,u_3)$. \pts{0,25 pt}
\end{enumerate}
\end{enumerate}

\Partie{B. Un sous-espace de $\R^4$}
Soit $F=\{(x,y,z,t)\in\R^4,\ x+y+z+t=0\ \text{et}\ x-y+2z=0\}$.
\begin{enumerate}[resume]
\item
\begin{enumerate}
\item Déterminer une famille génératrice de $F$, puis montrer que c'est une base de $F$. En déduire $\dim F$. \pts{0,75 pt}
\item Montrer que $w_1=(3,1,-1,-3)$ et $w_2=(0,2,1,-3)$ forment une base de $F$. \pts{0,5 pt}
\end{enumerate}
\end{enumerate}

\Partie{C. Une famille de polynômes}
Dans $\R_2[X]$, on pose $P_1=(X-1)(X-2)$, $P_2=X(X-2)$ et $P_3=X(X-1)$.
\begin{enumerate}[resume]
\item
\begin{enumerate}
\item Calculer $P_i(0)$, $P_i(1)$ et $P_i(2)$ pour $i=1,2,3$. \pts{0,25 pt}
\item Montrer que $(P_1,P_2,P_3)$ est une base de $\R_2[X]$. \pts{0,5 pt}
\item Soit $P\in\R_2[X]$. Exprimer les coordonnées de $P$ dans cette base à l'aide de $P(0)$, $P(1)$ et $P(2)$. \pts{0,25 pt}
\item En déduire l'unique polynôme $P\in\R_2[X]$ tel que $P(0)=1$, $P(1)=0$ et $P(2)=3$. \pts{0,25 pt}
\end{enumerate}
\end{enumerate}

\Partie{D. Un raisonnement sans coordonnées}
Soit $E$ un espace vectoriel de dimension 3 et $(u,v,w)$ une base de $E$.
\begin{enumerate}[resume]
\item
\begin{enumerate}
\item Montrer que $(u+v,\ v+w,\ w+u)$ est une base de $E$. \pts{0,5 pt}
\item La famille $(u-v,\ v-w,\ w-u)$ est-elle une base de $E$ ? \pts{0,25 pt}
\item Donner les coordonnées de $u$ dans la base $(u+v,\ v+w,\ w+u)$. \pts{0,25 pt}
\end{enumerate}
\end{enumerate}

\Ex{3}{Sous-espaces supplémentaires et projecteur}{3,5 points}
Soit $E=\R_3[X]$, de base canonique $\mathcal B_0=(1,X,X^2,X^3)$. On pose
\[F=\{P\in E,\ P(1)=0\ \text{et}\ P'(1)=0\}\qquad\text{et}\qquad G=\R_1[X].\]
On rappelle la formule de Taylor : pour tout $P\in E$,
$P=P(1)+P'(1)(X-1)+\frac{P''(1)}{2}(X-1)^2+\frac{P'''(1)}{6}(X-1)^3$.
\begin{enumerate}
\item Montrer que $F$ est un sous-espace vectoriel de $E$. \pts{0,5 pt}
\item
\begin{enumerate}
\item Montrer que $\mathcal B_1=\big(1,\ X-1,\ (X-1)^2,\ (X-1)^3\big)$ est une base de $E$. \pts{0,5 pt}
\item Donner les coordonnées d'un polynôme $P$ dans $\mathcal B_1$. En déduire une base de $F$ et sa dimension. \pts{0,75 pt}
\end{enumerate}
\item Montrer que $E=F\oplus G$. \pts{0,75 pt}
\item On note $q$ le projecteur sur $G$ parallèlement à $F$.
\begin{enumerate}
\item Montrer que $q(P)=P(1)+P'(1)(X-1)$ pour tout $P\in E$. Calculer $q(X^3)$. \pts{0,5 pt}
\item Écrire la matrice de $q$ dans la base $\mathcal B_1$, puis dans la base $\mathcal B_0$. \pts{0,5 pt}
\end{enumerate}
\end{enumerate}

\Ex{4}{Calculs de déterminants}{4 points}
\begin{enumerate}
\item Soient $a,b\in\R$. Calculer, sous forme factorisée, \pts{1,5 pt}
\[D_4=\begin{vmatrix}a&b&b&b\\b&a&b&b\\b&b&a&b\\b&b&b&a\end{vmatrix}.\]
Pour quelles valeurs de $a$ et $b$ la matrice correspondante est-elle inversible ?
\item Soient $a,b,c$ trois réels.
\begin{enumerate}
\item Calculer, sous forme factorisée, $V=\begin{vmatrix}1&a&a^2\\1&b&b^2\\1&c&c^2\end{vmatrix}$. \pts{1 pt}
\item On suppose $a,b,c$ deux à deux distincts. Soient $\alpha,\beta,\gamma\in\R$. Montrer qu'il existe un unique polynôme
$P\in\R_2[X]$ tel que $P(a)=\alpha$, $P(b)=\beta$ et $P(c)=\gamma$. \pts{0,5 pt}
\end{enumerate}
\item Soit $n$ un entier impair et $A\in\mathcal M_n(\R)$ une matrice antisymétrique ($A^T=-A$). Montrer que $A$ n'est pas
inversible. Le résultat reste-t-il vrai si $n$ est pair ? \pts{1 pt}
\end{enumerate}

\Pb{une matrice nilpotente}{4,5 points}
On considère
\[A=\begin{pmatrix}-1&-1&-1\\1&2&2\\0&-1&-1\end{pmatrix}\]
et $f$ l'endomorphisme de $\R^3$ canoniquement associé à $A$. On note $(e_1,e_2,e_3)$ la base canonique de $\R^3$.
\begin{enumerate}
\item Calculer $A^2$ et $A^3$. \pts{0,5 pt}
\item Déterminer une base de $\Ker A$ et une base de $\Im A$. Vérifier le théorème du rang, donner une équation de $\Im A$, et
montrer que $\Ker A\subset\Im A$. \pts{1 pt}
\item On pose $x_0=e_1$. Montrer, \emph{sans calculer de déterminant}, que $\mathcal B=(x_0,Ax_0,A^2x_0)$ est une base de
$\R^3$. \pts{1 pt}
\item Écrire la matrice $N$ de $f$ dans la base $\mathcal B$, la matrice de passage $P$ de la base canonique à $\mathcal B$, et
la relation entre $A$, $N$ et $P$. \pts{0,5 pt}
\item Déterminer tous les vecteurs $x=(x_1,x_2,x_3)$ de $\R^3$ tels que $(x,Ax,A^2x)$ soit une base de $\R^3$. \pts{1 pt}
\item Pour $\lambda\in\R$, calculer $\det(A+\lambda I_3)$ en utilisant la question 4. En déduire que $A+I_3$ est inversible et
que $A$ ne l'est pas. \pts{0,5 pt}
\end{enumerate}

\vfill
\begin{center}\sffamily\small\color{Gris}— Fin du sujet —\end{center}
\fi
\end{document}
"
\documentclass[a4paper,11pt]{article}

% ============================================================
%  Le Collé — Algèbre linéaire (compléments de deuxième année)
%  © 2026 Ahmed Amine Ouqlani
% ============================================================

% ---------- Polices ----------
\usepackage[utf8]{inputenc}
\usepackage[T1]{fontenc}
\usepackage[sc]{mathpazo}          % texte et mathématiques : Palatino
\usepackage[scaled=0.88]{helvet}   % titres et repères : linéale
\linespread{1.07}
\usepackage{amsmath,amssymb,amsthm,mathrsfs,stmaryrd}
\usepackage{microtype}

% ---------- Mise en page ----------
\usepackage[a4paper,top=2.3cm,bottom=2.3cm,left=2.4cm,right=2.4cm,
            headheight=14pt,headsep=0.8cm,footskip=1.1cm]{geometry}
\setlength{\parindent}{0pt}
\setlength{\parskip}{0.45em plus 0.15em minus 0.1em}
\raggedbottom
\emergencystretch=1.5em
\clubpenalty=10000 \widowpenalty=10000 \displaywidowpenalty=10000 \brokenpenalty=10000

% ---------- Couleurs : un bleu de nuit, un ocre, des gris ----------
\usepackage{xcolor}
\definecolor{Ink}{HTML}{1C2533}
\definecolor{Bleu}{HTML}{1F3A5F}
\definecolor{BleuMoyen}{HTML}{4A6A8F}
\definecolor{BleuClair}{HTML}{9DB2C9}
\definecolor{Fond}{HTML}{F2F4F7}
\definecolor{Ocre}{HTML}{B07A26}
\definecolor{Brique}{HTML}{A23B2A}
\definecolor{Gris}{HTML}{5E6672}
\definecolor{Filet}{HTML}{C5CCD6}

% ---------- Packages ----------
\usepackage{enumitem}
\setlist{itemsep=0.25em,topsep=0.35em,parsep=0pt}
\setlist[itemize]{label=\textcolor{BleuMoyen}{\small\textbullet},leftmargin=1.4em}
\setlist[enumerate]{leftmargin=1.8em,label=\textsf{\small\color{Bleu}\arabic*.}}
\usepackage{booktabs,array,multicol}
\usepackage{needspace,float}
\usepackage[font=small,labelfont={sf,color=Bleu},labelsep=period,skip=6pt]{caption}
\usepackage{tikz}
\usetikzlibrary{arrows.meta,calc,decorations.markings}
\usepackage{pgfplots}
\pgfplotsset{compat=1.17}
\usepackage{tkz-tab}
\usepackage[most]{tcolorbox}
\usepackage{titlesec,titletoc}
\usepackage{fancyhdr,lastpage,etoolbox}
\usepackage[colorlinks=true,linkcolor=Bleu,urlcolor=Bleu,citecolor=Bleu,
            pdftitle={Examen blanc — Algèbre linéaire et déterminants},pdfauthor={Ahmed Amine Ouqlani}]{hyperref}

\numberwithin{equation}{section}
\numberwithin{figure}{section}
\numberwithin{table}{section}

% ---------- Titres ----------
% numéro du chapitre dans le parcours PT (affiché), distinct de la numérotation interne
\newcommand{\grandnum}{\thesection}\newcommand{\ongletcouleur}{Bleu}
% Chapitres (sections) : grand numéro, titre en linéale, filet ocre
\titleformat{\section}[display]
  {\normalfont\sffamily\color{Ink}}
  {{\small\color{Gris}\textls[200]{\MakeUppercase{Chapitre}}}\\[-0.2em]%
   {\fontsize{64}{64}\selectfont\color{BleuClair}\grandnum}}
  {0.4em}
  {\fontsize{25}{30}\selectfont\bfseries}
  [\vspace{0.5em}{\color{Ocre}\rule{3.2cm}{1.4pt}}]
\titlespacing*{\section}{0pt}{0pt}{1.6em}
\titleformat{name=\section,numberless}[display]
  {\normalfont\sffamily\color{Ink}}{}{0pt}{\fontsize{25}{30}\selectfont\bfseries}
  [\vspace{0.5em}{\color{Ocre}\rule{3.2cm}{1.4pt}}]
\titlespacing*{name=\section,numberless}{0pt}{0pt}{1.6em}

\titleformat{\subsection}
  {\needspace{5\baselineskip}\normalfont\sffamily\Large\bfseries\color{Ink}}
  {\color{Bleu}\thesubsection}{0.8em}{}
\titlespacing*{\subsection}{0pt}{1.8em}{0.5em}
\titleformat{\subsubsection}
  {\needspace{5\baselineskip}\normalfont\sffamily\large\bfseries\color{Ink}}
  {\color{Bleu}\thesubsubsection}{0.7em}{}
\titlespacing*{\subsubsection}{0pt}{1.6em}{0.4em}
\titleformat{\paragraph}[runin]{\normalfont\sffamily\bfseries\color{Ink}}{}{0pt}{}[.]
\titlespacing*{\paragraph}{0pt}{0.9em}{0.6em}
\setcounter{secnumdepth}{3}
\setcounter{tocdepth}{2}

% ---------- Table des matières ----------
\titlecontents{section}[0pt]{\addvspace{0.7em}\sffamily\bfseries\color{Ink}}
  {}{}{\hfill\contentspage}[\addvspace{0.2em}]
\titlecontents{subsection}[2.2em]{\small\linespread{1}\selectfont}
  {\makebox[2.6em][l]{\color{Gris}\thecontentslabel}}{}{\titlerule*[0.7pc]{.}\contentspage}

% ---------- En-têtes et pieds de page ----------
\pagestyle{fancy}
\fancyhf{}
\renewcommand{\sectionmark}[1]{\markboth{#1}{}}
\fancyhead[L]{\footnotesize\sffamily\color{Gris}\leftmark}
\fancyhead[R]{\footnotesize\sffamily\scshape\color{Gris}Le Collé}
\fancyfoot[C]{\small\sffamily\color{Ink}\thepage}
\fancyfoot[L]{\scriptsize\sffamily\hyperlink{toc}{\color{Filet}sommaire}}
\renewcommand{\headrulewidth}{0.4pt}
\renewcommand{\headrule}{{\color{Filet}\hrule height\headrulewidth width\headwidth}}
\renewcommand{\footrulewidth}{0pt}
\fancypagestyle{plain}{\fancyhf{}\fancyfoot[C]{\small\sffamily\color{Ink}\thepage}\renewcommand{\headrulewidth}{0pt}}

% ---------- Énoncés ----------
\newtheoremstyle{enonce}{0pt}{0pt}{}{}{\sffamily\bfseries\color{Bleu}}{.}{0.5em}
  {\thmname{#1}\thmnumber{~#2}\thmnote{\ \normalfont\sffamily\color{Ink}(#3)}}
\newtheoremstyle{exemple}{0.9em}{0.9em}{}{}{\sffamily\bfseries\color{BleuMoyen}}{.}{0.5em}
  {\thmname{#1}\thmnumber{~#2}\thmnote{\ \normalfont\sffamily\color{Ink}(#3)}}
\newtheoremstyle{remarque}{0.7em}{0.7em}{}{}{\sffamily\itshape}{.}{0.5em}{\thmname{#1}\thmnote{\ (#3)}}
\theoremstyle{enonce}
\newtheorem{definition}{Définition}[section]
\newtheorem{theoreme}[definition]{Théorème}
\newtheorem{proposition}[definition]{Proposition}
\newtheorem{corollaire}[definition]{Corollaire}
\theoremstyle{exemple}
\newtheorem{exemple}[definition]{Exemple}
\theoremstyle{remarque}
\newtheorem*{remarque}{Remarque}

\tcbset{
  defi/.style={enhanced,breakable,colback=white,colframe=white,boxrule=0pt,arc=0pt,
     borderline west={1.2pt}{0pt}{Bleu},left=11pt,right=4pt,top=3pt,bottom=3pt,
     before skip=1.1em,after skip=1.1em},
  thm/.style={enhanced,breakable,colback=Fond,colframe=Fond,boxrule=0pt,arc=0pt,
     borderline west={2.4pt}{0pt}{Bleu},left=11pt,right=8pt,top=6pt,bottom=6pt,
     before skip=1.2em,after skip=1.2em}}
\tcolorboxenvironment{definition}{defi}
\tcolorboxenvironment{theoreme}{thm}
\tcolorboxenvironment{proposition}{thm}
\tcolorboxenvironment{corollaire}{thm}

\renewcommand{\proofname}{Démonstration}
\renewcommand{\qedsymbol}{\textcolor{BleuClair}{$\blacksquare$}}
\makeatletter
\renewenvironment{proof}[1][\proofname]{\par\pushQED{\qed}\normalfont
  \topsep0.5em\relax\trivlist\item[\hskip\labelsep\sffamily\itshape\color{Gris}#1.]\ignorespaces}
  {\popQED\endtrivlist\@endpefalse}
\makeatother

% ---------- Méthode : filets haut et bas, titre en bandeau ----------
\makeatletter\newcommand{\methlabelfix}{\protected@edef\@currentlabel{\thetcbcounter}}\makeatother
\newtcolorbox[auto counter,number within=section]{methode}[2][]{%
  enhanced,unbreakable,colback=white,colframe=white,boxrule=0pt,arc=0pt,
  borderline north={0.9pt}{0pt}{Bleu},borderline south={0.9pt}{0pt}{Bleu},
  left=4pt,right=4pt,top=8pt,bottom=8pt,before skip=1.4em,after skip=1.4em,
  before upper={\methlabelfix\phantomsection\addcontentsline{met}{metentry}{\protect\makebox[2.6em][l]{\thetcbcounter}#2}%
     {\sffamily\bfseries\small\color{white}\colorbox{Bleu}{\strut\ MÉTHODE~\thetcbcounter\ }}%
     \ \ {\sffamily\bfseries\color{Bleu}#2}\par\smallskip},#1}

% ---------- Repères pédagogiques : un filet de couleur, un mot-clé ----------
\newcommand{\definirbloc}[3]{% nom, couleur, libellé
  \newtcolorbox{#1}[1][]{enhanced,breakable,blanker,
    borderline west={1.4pt}{0pt}{#2},left=10pt,right=2pt,top=1pt,bottom=1pt,
    before skip=1em,after skip=1em,
    before upper={{\sffamily\bfseries\footnotesize\color{#2}\textls[120]{\MakeUppercase{#3}}}\quad},##1}}
\definirbloc{reflexe}{Bleu}{Réflexe}
\definirbloc{astuce}{Ocre}{Astuce}
\definirbloc{attention}{Brique}{Attention}
\definirbloc{enpratique}{BleuMoyen}{En pratique}
\definirbloc{plusloin}{Gris}{Pour aller plus loin}
\definirbloc{strategie}{Bleu}{Stratégie}
\definirbloc{lecture}{BleuMoyen}{Lire le schéma}
% compatibilité avec les parties reprises
\newenvironment{garde}{\begin{attention}}{\end{attention}}
\newenvironment{complementHP}{\begin{plusloin}}{\end{plusloin}}
\newenvironment{interpretation}{\begin{enpratique}}{\end{enpratique}}

% ---------- Exercices, corrigés ----------
\newcounter{exo}[section]
\renewcommand{\theexo}{\thesection.\arabic{exo}}
\newcommand{\niveau}[1]{{\color{Ocre}\footnotesize
  \ifcase#1\or$\bullet\circ\circ\circ\circ$\or$\bullet\bullet\circ\circ\circ$\or$\bullet\bullet\bullet\circ\circ$%
  \or$\bullet\bullet\bullet\bullet\circ$\or$\bullet\bullet\bullet\bullet\bullet$\fi}}
\newcommand{\serie}[1]{\par\needspace{8\baselineskip}\vspace{1.5em}%
  \noindent{\sffamily\bfseries\color{Bleu}#1}\par\nobreak\vspace{-0.4em}%
  \noindent{\color{Filet}\rule{\linewidth}{0.4pt}}\par\nobreak\vspace{0.1em}}
\newenvironment{exo}[3]{%
  \par\needspace{5\baselineskip}\vspace{0.75em}\refstepcounter{exo}\label{e:#1}%
  \def\exolab{#1}%
  \noindent{\sffamily\bfseries\color{Bleu}Exercice~\theexo}\quad\niveau{#2}%
  \hfill{\footnotesize\sffamily\color{Gris}#3}\par\nobreak\vspace{0.15em}}%
  {\par\nopagebreak\noindent\hfill{\scriptsize\sffamily\hyperref[c:\exolab]{\color{Gris}corrigé p.~\pageref{c:\exolab}}}\par}
\newenvironment{corrige}[1]{%
  \par\needspace{6\baselineskip}\vspace{1em}\phantomsection\label{c:#1}%
  \noindent{\sffamily\bfseries\color{Bleu}Exercice~\ref{e:#1}}\quad
  {\scriptsize\sffamily\hyperref[e:#1]{\color{Gris}énoncé p.~\pageref{e:#1}}}\par\nobreak\vspace{0.2em}}%
  {\hfill\textcolor{BleuClair}{$\square$}\par}
\newcommand{\etape}[1]{\par\smallskip\noindent{\sffamily\bfseries\small\color{BleuMoyen}#1.}\ }
\newenvironment{solution}{\par\noindent{\sffamily\bfseries\small\color{BleuMoyen}Solution.}\ }{\hfill\textcolor{BleuClair}{$\square$}\par}
\newcommand{\chercher}{\par\smallskip\noindent{\color{Filet}\rule[0.5ex]{1.2cm}{0.4pt}}\
  {\footnotesize\sffamily\itshape\color{Gris}laisser chercher l'étudiant avant de poursuivre}\par\smallskip}

% ---------- Guide de séance ----------

\newcommand{\marge}[1]{\leavevmode\marginpar{\raggedleft\scriptsize\sffamily\color{Gris}\textls[100]{\MakeUppercase{#1}}}\ignorespaces}
\newcommand{\Etape}[4]{%
  \subsection[Étape #1 --- #3 {\color{Gris}(#2)}]{#3}\label{etape:#1}%
  \noindent{\small\sffamily\color{Gris}{\color{Ocre}\bfseries Étape #1}\quad·\quad #2\quad·\quad #4}\par\medskip}
\newcommand{\question}[1]{\par\smallskip\noindent{\sffamily\bfseries\small\color{BleuMoyen}Question.}\ \textit{#1}\par\smallskip}
\newcommand{\rappel}[1]{\par\smallskip\noindent\phantomsection\addcontentsline{rap}{rapentry}{#1}%
  {\sffamily\bfseries\scriptsize\colorbox{Fond}{\color{Gris}\strut RAPPEL}}\ {\sffamily\small\color{Gris}#1.}\ }
\newcommand{\pourquoi}{\par\noindent{\sffamily\bfseries\small\color{Ocre}Pourquoi ?}\ }

% Ouverture de chapitre : \ouverture{objectifs}{prérequis}{durée}
\newcommand{\ouverture}[3]{%
  \begin{tcolorbox}[enhanced,colback=Fond,colframe=Fond,arc=0pt,boxrule=0pt,left=12pt,right=12pt,top=10pt,bottom=10pt,
     before skip=1.2em,after skip=1.6em]
  \begin{minipage}[t]{0.56\linewidth}
  {\sffamily\bfseries\small\color{Bleu}OBJECTIFS}\par\vspace{0.3em}\small #1
  \end{minipage}\hfill
  \begin{minipage}[t]{0.38\linewidth}
  {\sffamily\bfseries\small\color{Bleu}PRÉREQUIS}\par\vspace{0.3em}\small #2
  \ifx\relax#3\relax\else\par\vspace{0.6em}{\sffamily\bfseries\small\color{Bleu}DURÉE}\par\vspace{0.2em}\small #3\fi
  \end{minipage}
  \end{tcolorbox}}
% Fin de chapitre
\newenvironment{essentiel}{\par\needspace{12\baselineskip}\vspace{1.5em}%
  \begin{tcolorbox}[enhanced,unbreakable,colback=white,colframe=white,boxrule=0pt,arc=0pt,
   borderline north={1.4pt}{0pt}{Ocre},left=2pt,right=2pt,top=8pt,bottom=4pt]
  {\sffamily\bfseries\color{Ocre}L'ESSENTIEL DU CHAPITRE}\par\medskip}{\end{tcolorbox}}

% ---------- Notations ----------
\newcommand{\K}{\mathbb{K}}
\newcommand{\R}{\mathbb{R}}
\newcommand{\C}{\mathbb{C}}
\newcommand{\N}{\mathbb{N}}
\newcommand{\Z}{\mathbb{Z}}
\newcommand{\Mn}[1]{\mathcal{M}_{n}(#1)}
\newcommand{\GLn}[1]{\mathrm{GL}_n(#1)}
\newcommand{\LE}{\mathcal{L}(E)}
\newcommand{\rg}{\operatorname{rg}}
\newcommand{\vect}{\operatorname{Vect}}
\newcommand{\Mat}{\operatorname{Mat}}
\newcommand{\Id}{\operatorname{Id}}
\newcommand{\sh}{\operatorname{sh}}
\newcommand{\ch}{\operatorname{ch}}
\renewcommand{\th}{\operatorname{th}}
\DeclareMathOperator{\argsh}{argsh}
\newcommand{\vv}[1]{\overrightarrow{#1}}
\newcommand{\ve}[1]{\vec{#1}}
\newcommand{\ient}[2]{\llbracket #1,#2\rrbracket}
\newcommand{\dt}[1]{\left|\begin{matrix}#1\end{matrix}\right|}
\newcommand{\dis}{\displaystyle}

% ---------- Figures ----------
\pgfplotsset{trig format plots=rad}
\pgfplotsset{courbe/.style={axis lines=middle,axis equal image,
  axis line style={Gris,-{Stealth[length=5pt]}},
  xlabel={$x$},ylabel={$y$},ticklabel style={font=\scriptsize\sffamily,color=Gris},
  every axis x label/.style={at={(ticklabel* cs:1)},anchor=west,font=\small},
  every axis y label/.style={at={(ticklabel* cs:1)},anchor=south,font=\small},
  samples=300,smooth,width=0.6\linewidth}}
\tikzset{arcetude/.style={Bleu,line width=1.25pt},
         reste/.style={BleuClair,line width=1pt},
         tgte/.style={Ocre,line width=0.9pt},
         asy/.style={Gris,dashed,line width=0.6pt},
         pt/.style={circle,fill=Ink,inner sep=1.2pt}}

% ---------- Navigation ----------
\makeatletter
\newcommand*\l@rapentry{\@dottedtocline{1}{0em}{0em}}
\newcommand*\l@metentry{\@dottedtocline{1}{0em}{0em}}
\newcommand{\listemethodes}{\@starttoc{met}}
\newcommand{\listerappels}{\@starttoc{rap}}
\makeatother
% onglets de chapitre sur le bord droit
\usepackage{eso-pic}
\newif\ifonglets \ongletsfalse
\newcommand{\dessineronglet}{%
  \edef\ongcol{\ongletcouleur}%
  \edef\ongnum{\grandnum}%
  \pgfmathsetmacro{\ongy}{3.2+2.3*(\ongnum-1)}%
  \begin{tikzpicture}[remember picture,overlay]
  \fill[\ongcol] ([xshift=-0.75cm,yshift=-\ongy cm]current page.north east)
     rectangle ([yshift=-\ongy cm-1.9cm]current page.north east);
  \node[text=white,font=\sffamily\bfseries\large] at ([xshift=-0.37cm,yshift=-\ongy cm-0.95cm]current page.north east)
     {\ongnum};
  \end{tikzpicture}}
\AddToShipoutPictureBG{\ifonglets\ifnum\value{section}>0 \dessineronglet\fi\fi}
% blocs de navigation et d'enseignement
\definirbloc{commencer}{Gris}{Si vous commencez ici}
\definirbloc{horsprogramme}{Brique}{Hors programme}
\newcommand{\HP}{\,{\sffamily\bfseries\scriptsize\textcolor{Brique}{[HP]}}}
\newcommand{\source}[1]{{\footnotesize\sffamily\color{Gris}\,(#1)}}
\newcommand{\coupdepouce}[1]{\par\vspace{-0.2em}\noindent{\footnotesize\sffamily\color{Gris}%
  \textbf{Coup de pouce} (à donner si l'étudiant bloque) : #1}\par}
\newcommand{\pointarret}[1]{\par\vspace{1em}\noindent{\color{Filet}\rule[0.5ex]{\linewidth}{0.4pt}}\par\nobreak
  \noindent{\footnotesize\sffamily\color{Ocre}\bfseries POINT D'ARRÊT POSSIBLE}\quad{\footnotesize\sffamily\color{Gris}#1}\par\vspace{0.6em}}
\newcommand{\demarche}[1]{\noindent{\sffamily\bfseries\small\color{BleuMoyen}Démarche.}\ \textit{#1}\par\smallskip}

% ---------- Corrigés : idées, astuces, pièges, à retenir ----------
\newcommand{\choix}[1]{\par\noindent{\sffamily\bfseries\small\color{BleuMoyen}Choix de la méthode.}\ #1\par\smallskip}
\newcommand{\gain}[1]{\par\noindent{\sffamily\bfseries\small\color{Ocre}Astuce.}\ #1\par}
\newcommand{\piege}[1]{\par\noindent{\sffamily\bfseries\small\color{Brique}Piège.}\ #1\par}
\newtcolorbox{retenirbox}{enhanced,breakable,colback=Fond,colframe=Fond,arc=0pt,boxrule=0pt,
  borderline west={2pt}{0pt}{Ocre},left=8pt,right=6pt,top=4pt,bottom=4pt,before skip=0.7em,after skip=0.5em}
\newcommand{\retenir}[1]{\begin{retenirbox}{\sffamily\bfseries\footnotesize\color{Ocre}À RETENIR}\quad #1\end{retenirbox}}
\newcommand{\TD}[1]{{\color{Ocre}\bfseries TD\,#1}\ ·\ }
\newcommand{\q}[1]{\par\smallskip\noindent{\sffamily\bfseries\color{Bleu}#1}\ }
\providecommand{\iddots}{\mathinner{\mkern1mu\raise1pt\hbox{.}\mkern2mu\raise4pt\hbox{.}\mkern2mu\raise7pt\hbox{.}\mkern1mu}}


% =====================================================================
%  Examen blanc — Algèbre linéaire et déterminants (Le Collé, PT)
%  Une source, deux PDF : \def\CORRIGE{} avant \input donne le corrigé.
% =====================================================================
\providecommand{\Ker}{\operatorname{Ker}}\renewcommand{\Im}{\operatorname{Im}}
\fancyfoot[L]{}\fancyhead[L]{\footnotesize\sffamily\color{Gris}Examen blanc · Algèbre linéaire et déterminants\ifdefined\CORRIGE\ · corrigé\fi}
\newcommand{\pts}[1]{\hfill{\footnotesize\sffamily\color{Gris}(#1)}}
\newcommand{\Ex}[3]{\par\needspace{6\baselineskip}\vspace{1.3em}\noindent
  {\sffamily\bfseries\large\color{Bleu}Exercice #1 \,·\, #2}\hfill{\sffamily\small\color{Ocre}#3}\par
  \vspace{2pt}{\color{Filet}\rule{\linewidth}{0.4pt}}\par\smallskip}
\newcommand{\Pb}[2]{\par\needspace{6\baselineskip}\vspace{1.3em}\noindent
  {\sffamily\bfseries\large\color{Bleu}Problème \,·\, #1}\hfill{\sffamily\small\color{Ocre}#2}\par
  \vspace{2pt}{\color{Filet}\rule{\linewidth}{0.4pt}}\par\smallskip}
\newcommand{\Partie}[1]{\par\needspace{4\baselineskip}\medskip\noindent{\sffamily\bfseries\color{BleuMoyen}#1}\par\nobreak\smallskip}
\newcommand{\Q}[1]{\par\needspace{3\baselineskip}\medskip\noindent{\sffamily\bfseries\color{Bleu}#1}\enspace}
\newcommand{\autre}[1]{\par\noindent{\sffamily\bfseries\small\color{BleuMoyen}Autre méthode.}\ #1\par}
\newcommand{\lienDL}[1]{\par\noindent{\sffamily\bfseries\small\color{Ocre}Lien avec le DL.}\ #1\par}
\setlist[enumerate,1]{label=\arabic*.,leftmargin=1.6em,itemsep=4pt,topsep=4pt}
\setlist[enumerate,2]{label=(\alph*),leftmargin=1.6em,itemsep=2pt,topsep=2pt}

\begin{document}
\pagestyle{fancy}
\thispagestyle{plain}

% ------------------------------------------------------------------ en-tête
\begin{tikzpicture}[remember picture,overlay]
\fill[Bleu] (current page.north west) rectangle ([yshift=-4.2cm]current page.north east);
\node[anchor=north west,text=white,font=\sffamily\footnotesize,inner sep=0] at ([xshift=2.2cm,yshift=-1.1cm]current page.north west)
  {\textls[120]{MATHÉMATIQUES · CPGE PT · EXAMEN BLANC}};
\node[anchor=north west,text=white,font=\sffamily\bfseries\LARGE,inner sep=0] at ([xshift=2.2cm,yshift=-1.75cm]current page.north west)
  {Algèbre linéaire et déterminants};
\node[anchor=north west,text=white,font=\sffamily\normalsize,inner sep=0] at ([xshift=2.2cm,yshift=-2.75cm]current page.north west)
  {\ifdefined\CORRIGE Corrigé détaillé et diagnostic\else Durée : 3 heures \quad·\quad Calculatrice interdite\fi};
\node[anchor=north east,text=white,font=\sffamily\scshape\large,inner sep=0] at ([xshift=-2.2cm,yshift=-1.05cm]current page.north east)
  {Le Collé};
\end{tikzpicture}
\vspace*{2.9cm}

\ifdefined\CORRIGE
% =====================================================================
%                              CORRIGÉ
% =====================================================================
\markboth{Corrigé · Examen blanc}{}
\paragraph{Comment utiliser ce corrigé} Ne l'ouvre qu'après avoir composé les 3 heures. Corrige d'abord ta copie avec un
stylo d'une autre couleur, question par question, en comptant les points. Pour chaque question ratée, lis le \emph{choix de la
méthode} : c'est lui qu'il faut retenir, plus que le calcul. La grille de la dernière page dit quoi revoir.

\subsection*{Ce que le sujet évalue}
\begin{center}
\small\renewcommand{\arraystretch}{1.22}
\begin{tabular}{@{}>{\raggedright\arraybackslash}p{7.6cm}>{\raggedright\arraybackslash}p{5.4cm}@{}}
\toprule
\textbf{Compétence} & \textbf{Questions}\\\midrule
Démontrer qu'une famille est libre (combinaison nulle, évaluation, composition) & 1.2 · 2.3(b) · 2.4(a) · P.3\\
Démontrer qu'une famille est génératrice, trouver une base d'un sous-espace & 1.3 · 2.2(a) · 3.2\\
Utiliser la dimension (« $n$ vecteurs en dimension $n$ ») & 1.1 · 2.2(b) · 2.4(a) · 3.3 · P.3\\
Coordonnées dans une base, changement de base & 2.1(b) · 2.3(c) · 2.4(c) · 3.4 · P.4\\
Sous-espaces, somme directe, supplémentaires, projecteur & 1.4 · 3.1 · 3.3 · 3.4\\
Application linéaire : noyau, image, rang, théorème du rang & 1.6 · P.2\\
Calculer un déterminant (opérations, développement, factorisation) & 2.1(a) · 4.1 · 4.2\\
Propriétés du déterminant ($\det(\lambda A)$, $\det A^T$, $\det(AB)$) & 1.5 · 4.3 · P.6\\
Déterminant, inversibilité, liberté, systèmes de Cramer & 2.1 · 4.1 · 4.2 · P.6\\
« Trouver tous les… » : prouver les deux sens & 2.1(a) · P.5\\\bottomrule
\end{tabular}
\end{center}

\lienDL{Les questions 2.3(b), 4.2 et P.3 à P.5 reprennent les difficultés du DL~3 sous une autre forme : liberté par évaluation
(comme III.3), Vandermonde, famille $(x,Ax,A^2x)$ avec $A^3=0$ (comme II.1 et II.4.a). Les formulations sont différentes : il
faut reconnaître la technique.}

\begin{center}\small\renewcommand{\arraystretch}{1.2}
\begin{tabular}{@{}lcccccc@{}}\toprule
& Ex.~1 & Ex.~2 & Ex.~3 & Ex.~4 & Problème & Total\\\midrule
Points & 3 & 5 & 3,5 & 4 & 4,5 & 20\\
Temps conseillé & 15 min & 50 min & 35 min & 35 min & 45 min & 3 h\\\bottomrule
\end{tabular}
\end{center}

% ---------------------------------------------------------------------
\clearpage
\Ex{1}{Vrai ou faux}{3 points}

\Q{1. Vrai.} $\dim\R^3=3$ : une famille libre de $\R^3$ a au plus 3 vecteurs. Une famille de 4 vecteurs est donc liée.

\Q{2. Faux.} Contre-exemple : $u=(1,0,0)$, $v=(0,1,0)$, $w=(1,0,0)$. Les familles $(u,v)$ et $(v,w)$ sont libres, mais $u-w=0$ :
$(u,v,w)$ est liée.
\piege{la liberté ne se vérifie pas « deux à deux ». Trois vecteurs deux à deux non colinéaires peuvent être liés, par exemple
$(1,0)$, $(0,1)$, $(1,1)$ dans $\R^2$.}

\Q{3. Vrai.} $\dim\R^3=3$ et la famille a 3 vecteurs : pour une telle famille, génératrice $\iff$ libre $\iff$ base.

\Q{4. Faux.} Dans $\R^2$, $F=\vect\big((1,0)\big)$ et $G=\vect\big((0,1)\big)$ : $(1,0)$ et $(0,1)$ sont dans $F\cup G$, mais pas
leur somme $(1,1)$. $F\cup G$ n'est pas stable par addition.

\Q{5. Faux.} Multiplier $A$ par $2$ multiplie chacune des 3 colonnes par $2$, et le déterminant est linéaire par rapport à
\emph{chaque} colonne : $\det(2A)=2^3\det A=8\det A$. Par exemple $\det(2I_3)=8$.
\retenir{$\det(\lambda A)=\lambda^n\det A$ pour $A\in\Mn{\K}$.}

\Q{6. Vrai.} Par le théorème du rang, $\dim\Ker f=4-\rg f$. Or $\Im f\subset\R^3$, donc $\rg f\leq3$ et $\dim\Ker f\geq1$ :
$\Ker f\neq\{0\}$, $f$ n'est pas injective.

% ---------------------------------------------------------------------
\Ex{2}{Familles libres, génératrices, bases}{5 points}

\Partie{A. Une famille de $\R^3$ avec paramètre}

\Q{2.1(a)} \choix{3 vecteurs en dimension 3 avec un paramètre : un seul déterminant donne toutes les valeurs de $m$ d'un coup.
Résoudre un système avec $m$ obligerait à discuter à chaque pivot.}
\[\det(u_1,u_2,u_3)=\begin{vmatrix}1&0&2\\2&1&m\\-1&1&1\end{vmatrix}
\underset{C_3\leftarrow C_3-2C_1}{=}\begin{vmatrix}1&0&0\\2&1&m-4\\-1&1&3\end{vmatrix}=1\times\begin{vmatrix}1&m-4\\1&3\end{vmatrix}=7-m.\]
Le déterminant est non nul si et seulement si $m\neq7$. \textbf{$(u_1,u_2,u_3)$ est une base de $\R^3$ si et seulement si $m\neq7$.}
Les deux sens sont prouvés en même temps, car « base $\iff\det\neq0$ ».

\Q{2.1(b)} Pour $m=1$, on cherche $(a,b,c)$ tel que $au_1+bu_2+cu_3=(1,1,1)$ :
\[\begin{cases}a+2c=1\\2a+b+c=1\\-a+b+c=1\end{cases}\]
La différence des deux dernières lignes donne $3a=0$, donc $a=0$, puis $c=\frac12$ et $b=\frac12$. \textbf{Les coordonnées de
$(1,1,1)$ dans $\mathcal B'$ sont $\left(0,\frac12,\frac12\right)$.} Vérification : $\frac12(u_2+u_3)=\frac12(2,2,2)=(1,1,1)$.
\piege{les coordonnées dépendent de la base : $(1,1,1)$ est la colonne dans la base canonique, $(0,\frac12,\frac12)$ dans
$\mathcal B'$. Toujours vérifier en recombinant.}

\Q{2.1(c)} Pour $m=7$, on cherche une relation $au_1+bu_2+u_3=0$ : la première coordonnée donne $a=-2$, la troisième $b=-3$, et
la deuxième est alors vérifiée. Donc $u_3=2u_1+3u_2$. Comme $u_1$ et $u_2$ ne sont pas colinéaires,
$\rg(u_1,u_2,u_3)=\dim\vect(u_1,u_2)=2$.

\Partie{B. Un sous-espace de $\R^4$}

\Q{2.2(a)} \choix{un sous-espace défini par des équations : on résout le système, et on écrit la solution générale comme
combinaison linéaire. Cela donne directement une famille génératrice.}

$(x,y,z,t)\in F\iff x=y-2z$ et $t=-x-y-z=-2y+z$. Donc
\[F=\{\,y(1,1,0,-2)+z(-2,0,1,1),\ y,z\in\R\,\}=\vect(f_1,f_2),\]
avec $f_1=(1,1,0,-2)$ et $f_2=(-2,0,1,1)$.
\emph{Génératrice :} c'est l'égalité ci-dessus. \emph{Libre :} si $yf_1+zf_2=0$, les deuxième et troisième coordonnées donnent
$y=0$ et $z=0$. \textbf{$(f_1,f_2)$ est une base de $F$, et $\dim F=2$.}
\autre{$F$ est défini par 2 équations indépendantes dans $\R^4$, donc $\dim F=4-2=2$ (théorème du rang appliqué à
$(x,y,z,t)\mapsto(x+y+z+t,\,x-y+2z)$, de rang 2).}

\Q{2.2(b)} \choix{on connaît maintenant $\dim F=2$. Pour 2 vecteurs dans un espace de dimension 2, il suffit qu'ils soient
dans $F$ et libres : inutile de prouver qu'ils engendrent $F$.}

$w_1=(3,1,-1,-3)$ : $3+1-1-3=0$ et $3-1-2=0$, donc $w_1\in F$. $w_2=(0,2,1,-3)$ : $0+2+1-3=0$ et $0-2+2=0$, donc $w_2\in F$. Ils
ne sont pas colinéaires (la première coordonnée de $w_2$ est nulle, pas celle de $w_1$, et $w_2\neq0$). \textbf{$(w_1,w_2)$ est une
famille libre de 2 vecteurs de $F$, avec $\dim F=2$ : c'est une base de $F$.}
\piege{oublier de vérifier que $w_1,w_2\in F$. Deux vecteurs libres de $\R^4$ ne forment pas une base de $F$ s'ils n'y sont pas.}

\Partie{C. Une famille de polynômes}

\Q{2.3(a)}
\begin{center}\small\renewcommand{\arraystretch}{1.2}
\begin{tabular}{@{}cccc@{}}\toprule
& $x=0$ & $x=1$ & $x=2$\\\midrule
$P_1=(X-1)(X-2)$ & $2$ & $0$ & $0$\\
$P_2=X(X-2)$ & $0$ & $-1$ & $0$\\
$P_3=X(X-1)$ & $0$ & $0$ & $2$\\\bottomrule
\end{tabular}
\end{center}

\Q{2.3(b)} \choix{les trois polynômes ont le même degré : l'argument « degrés distincts » ne marche pas. Le tableau montre
que chaque $P_i$ ne vit qu'en un point. On évalue la relation en $0$, $1$, $2$ : chaque évaluation isole un coefficient.}

Soient $a,b,c$ tels que $aP_1+bP_2+cP_3=0$. On évalue en $0$ : $2a=0$. En $1$ : $-b=0$. En $2$ : $2c=0$. Donc $a=b=c=0$ : la
famille est libre. Elle a 3 vecteurs et $\dim\R_2[X]=3$ : \textbf{c'est une base de $\R_2[X]$}.
\lienDL{c'est la même idée qu'en III.3 : une relation entre vecteurs devient une information sur des valeurs, qu'on lit point
par point.}

\Q{2.3(c)} D'après (b), $P=aP_1+bP_2+cP_3$ pour des réels $a,b,c$. On évalue en $0$, $1$, $2$ :
$P(0)=2a$, $P(1)=-b$, $P(2)=2c$. Donc
\[P=\frac{P(0)}2\,P_1-P(1)\,P_2+\frac{P(2)}2\,P_3.\]

\Q{2.3(d)} D'après (c), si $P(0)=1$, $P(1)=0$, $P(2)=3$, nécessairement $P=\frac12P_1+\frac32P_3=2X^2-3X+1$. Réciproquement,
ce polynôme vérifie $P(0)=1$, $P(1)=0$, $P(2)=8-6+1=3$. \textbf{L'unique solution est $P=2X^2-3X+1$.}

\Partie{D. Un raisonnement sans coordonnées}

\Q{2.4(a)} \choix{on ne connaît pas les vecteurs, seulement que $(u,v,w)$ est une base. On écrit une combinaison nulle et on la
réécrit en fonction de $u,v,w$ pour utiliser leur liberté.}

Soient $\alpha,\beta,\gamma$ tels que $\alpha(u+v)+\beta(v+w)+\gamma(w+u)=0$, soit $(\alpha+\gamma)u+(\alpha+\beta)v+(\beta+\gamma)w=0$.
Comme $(u,v,w)$ est libre : $\alpha+\gamma=\alpha+\beta=\beta+\gamma=0$. En additionnant, $\alpha+\beta+\gamma=0$, puis
$\beta=\gamma=\alpha=0$. La famille est libre, a 3 vecteurs, et $\dim E=3$ : \textbf{c'est une base de $E$.}
\autre{la matrice de la famille dans la base $(u,v,w)$ est $\begin{pmatrix}1&0&1\\1&1&0\\0&1&1\end{pmatrix}$, de déterminant
$2\neq0$.}

\Q{2.4(b)} $(u-v)+(v-w)+(w-u)=0$ : c'est une relation à coefficients non tous nuls. \textbf{La famille est liée.}

\Q{2.4(c)} $(u+v)-(v+w)+(w+u)=2u$, donc $u=\frac12(u+v)-\frac12(v+w)+\frac12(w+u)$ : \textbf{coordonnées
$\left(\frac12,-\frac12,\frac12\right)$}.

% ---------------------------------------------------------------------
\Ex{3}{Sous-espaces supplémentaires et projecteur}{3,5 points}

\Q{3.1} L'application $\varphi:P\mapsto\big(P(1),P'(1)\big)$, de $E$ dans $\R^2$, est linéaire (l'évaluation et la dérivation le
sont). $F=\Ker\varphi$ est donc un sous-espace vectoriel de $E$.
\autre{vérifier directement : le polynôme nul est dans $F$, et si $P,Q\in F$, alors $(\lambda P+\mu Q)(1)=0$ et
$(\lambda P+\mu Q)'(1)=0$.}

\Q{3.2(a)} Les polynômes $1$, $X-1$, $(X-1)^2$, $(X-1)^3$ sont non nuls et de degrés $0,1,2,3$ distincts : la famille est libre.
Elle a 4 vecteurs et $\dim E=4$ : \textbf{$\mathcal B_1$ est une base de $E$}.

\Q{3.2(b)} \choix{la formule de Taylor donne directement les coordonnées dans $\mathcal B_1$, et les conditions qui
définissent $F$ portent exactement sur les deux premières.}

Par la formule de Taylor, les coordonnées de $P$ dans $\mathcal B_1$ sont $\left(P(1),\,P'(1),\,\frac{P''(1)}2,\,\frac{P'''(1)}6\right)$.
Donc $P\in F$ si et seulement si ses deux premières coordonnées sont nulles, c'est-à-dire si et seulement si
$P\in\vect\big((X-1)^2,(X-1)^3\big)$. La famille $\big((X-1)^2,(X-1)^3\big)$ engendre $F$, et elle est libre car extraite de la
base $\mathcal B_1$. \textbf{C'est une base de $F$, et $\dim F=2$.}

\Q{3.3} \choix{deux méthodes possibles. La plus rapide : intersection nulle et somme des dimensions.}

Soit $P\in F\cap G$. Comme $P\in G=\R_1[X]$, $P=a+b(X-1)$, donc $P(1)=a$ et $P'(1)=b$. Comme $P\in F$, $a=b=0$ : $P=0$. Ainsi
$F\cap G=\{0\}$, et $\dim F+\dim G=2+2=4=\dim E$. \textbf{Donc $E=F\oplus G$.}
\autre{$G=\vect\big(1,X-1\big)$ et $F=\vect\big((X-1)^2,(X-1)^3\big)$ : en réunissant ces deux bases, on obtient la base
$\mathcal B_1$ de $E$, donc $E=F\oplus G$.}

\Q{3.4(a)} D'après la formule de Taylor,
\[P=\underbrace{P(1)+P'(1)(X-1)}_{\in\,G}+\underbrace{\frac{P''(1)}2(X-1)^2+\frac{P'''(1)}6(X-1)^3}_{\in\,F}.\]
La décomposition de $P$ selon $E=F\oplus G$ est unique ; la composante dans $G$ est $q(P)$. \textbf{Donc
$q(P)=P(1)+P'(1)(X-1)$.} Pour $P=X^3$ : $q(X^3)=1+3(X-1)=3X-2$. Contrôle : $X^3-(3X-2)=(X-1)^2(X+2)$ est bien dans $F$.

\Q{3.4(b)} $q$ laisse fixes les vecteurs de $G$ et annule ceux de $F$ :
\[\Mat_{\mathcal B_1}(q)=\begin{pmatrix}1&0&0&0\\0&1&0&0\\0&0&0&0\\0&0&0&0\end{pmatrix}.\]
Dans $\mathcal B_0$ : $q(1)=1$, $q(X)=X$, $q(X^2)=1+2(X-1)=2X-1$, $q(X^3)=3X-2$, donc
\[\Mat_{\mathcal B_0}(q)=\begin{pmatrix}1&0&-1&-2\\0&1&2&3\\0&0&0&0\\0&0&0&0\end{pmatrix}.\]
On vérifie que cette matrice $Q$ vérifie $Q^2=Q$, comme tout projecteur.
\retenir{dans une base adaptée à $E=F\oplus G$, la matrice d'un projecteur est diagonale avec des $1$ et des $0$. Choisir la
bonne base simplifie tout : c'est l'idée de toute la suite du cours.}

% ---------------------------------------------------------------------
\Ex{4}{Calculs de déterminants}{4 points}

\Q{4.1} \choix{la somme de chaque ligne vaut $a+3b$ : on additionne toutes les colonnes dans la première pour factoriser.}
\begin{align*}D_4&=\begin{vmatrix}a&b&b&b\\b&a&b&b\\b&b&a&b\\b&b&b&a\end{vmatrix}
\underset{C_1\leftarrow C_1+C_2+C_3+C_4}{=}(a+3b)\begin{vmatrix}1&b&b&b\\1&a&b&b\\1&b&a&b\\1&b&b&a\end{vmatrix}\\
&\underset{L_i\leftarrow L_i-L_1}{=}(a+3b)\begin{vmatrix}1&b&b&b\\0&a-b&0&0\\0&0&a-b&0\\0&0&0&a-b\end{vmatrix}.\end{align*}
La dernière matrice est triangulaire : \textbf{$D_4=(a+3b)(a-b)^3$}. La matrice est inversible si et seulement si $D_4\neq0$,
c'est-à-dire \textbf{$a\neq b$ et $a\neq-3b$}.
\piege{développer directement un $4\times4$ : 24 termes et presque sûrement une erreur. Toujours chercher d'abord une
structure (somme constante, colonnes proches, zéros).}

\Q{4.2(a)} On fait $L_2\leftarrow L_2-L_1$ et $L_3\leftarrow L_3-L_1$, puis on factorise :
\[V=\begin{vmatrix}1&a&a^2\\0&b-a&b^2-a^2\\0&c-a&c^2-a^2\end{vmatrix}=(b-a)(c-a)\begin{vmatrix}1&b+a\\1&c+a\end{vmatrix}=(b-a)(c-a)(c-b).\]
\gain{$b^2-a^2=(b-a)(b+a)$ : on fait apparaître les facteurs $(b-a)$ et $(c-a)$ avant de développer.}

\Q{4.2(b)} $P=x_0+x_1X+x_2X^2$ vérifie les trois conditions si et seulement si
\[\begin{pmatrix}1&a&a^2\\1&b&b^2\\1&c&c^2\end{pmatrix}\begin{pmatrix}x_0\\x_1\\x_2\end{pmatrix}=\begin{pmatrix}\alpha\\\beta\\\gamma\end{pmatrix}.\]
Le déterminant de ce système est $V\neq0$ car $a,b,c$ sont distincts : c'est un système de Cramer, il a une unique solution.
\textbf{Il existe un unique $P\in\R_2[X]$ qui convient.}
\lienDL{c'est le même phénomène qu'en 2.3, où $a=0$, $b=1$, $c=2$, et qu'en III.3 du DL : trois valeurs distinctes
« séparent » les polynômes de degré au plus 2. Ici on le prouve par le déterminant, en 2.3 par évaluation.}

\Q{4.3} \choix{on ne peut rien calculer : on utilise les propriétés, $\det A^T=\det A$ et $\det(\lambda A)=\lambda^n\det A$.}

$A^T=-A$, donc $\det A=\det A^T=\det(-A)=(-1)^n\det A=-\det A$ car $n$ est impair. Ainsi $2\det A=0$ : \textbf{$\det A=0$, et
$A$ n'est pas inversible.} Pour $n$ pair, c'est faux : $\begin{pmatrix}0&1\\-1&0\end{pmatrix}$ est antisymétrique, de
déterminant $1$.
\piege{écrire $\det(-A)=-\det A$ sans justifier : c'est vrai seulement parce que $n$ est impair.}

% ---------------------------------------------------------------------
\Pb{une matrice nilpotente}{4,5 points}

\Q{P.1}
\[A^2=\begin{pmatrix}0&0&0\\1&1&1\\-1&-1&-1\end{pmatrix},\qquad A^3=A^2A=0.\]
(Chaque ligne de $A^2$ est proportionnelle à $(1,1,1)$, et $(1,1,1)A=(0,0,0)$.)

\Q{P.2} \emph{Noyau.} $AX=0$ s'écrit $-x-y-z=0$, $x+2y+2z=0$, $-y-z=0$. La dernière équation donne $z=-y$, puis la première
$x=0$ ; la deuxième est alors vérifiée. \textbf{$\Ker A=\vect\big((0,1,-1)\big)$}, de dimension 1.

\emph{Image.} Elle est engendrée par les colonnes $C_1=(-1,1,0)$, $C_2=(-1,2,-1)$, $C_3=C_2$. $C_1$ et $C_2$ ne sont pas
colinéaires : \textbf{$\Im A=\vect(C_1,C_2)$, et $\rg A=2$}. Théorème du rang : $1+2=3$. \checkmark

\emph{Équation.} $C_1$ et $C_2$ vérifient $x+y+z=0$, et ce plan est de dimension 2 : \textbf{$\Im A$ a pour équation $x+y+z=0$}.

\emph{Inclusion.} $(0,1,-1)$ vérifie $0+1-1=0$, donc $\Ker A\subset\Im A$.

\Q{P.3} \choix{l'énoncé interdit le déterminant : on utilise $A^3=0$. Appliquer $A^2$ à une relation fait disparaître tous
les termes sauf un.}

On a $x_0=e_1$, $Ax_0=(-1,1,0)$, $A^2x_0=(0,1,-1)\neq0$. Soient $c_0,c_1,c_2$ tels que $c_0x_0+c_1Ax_0+c_2A^2x_0=0$.
\begin{itemize}[leftmargin=1.2em,itemsep=1pt]
\item On multiplie par $A^2$ : comme $A^3=A^4=0$, il reste $c_0A^2x_0=0$, et $A^2x_0\neq0$, donc $c_0=0$.
\item On multiplie alors par $A$ : il reste $c_1A^2x_0=0$, donc $c_1=0$.
\item Il reste $c_2A^2x_0=0$, donc $c_2=0$.
\end{itemize}
La famille est libre, de 3 vecteurs en dimension 3 : \textbf{$\mathcal B$ est une base de $\R^3$.}
\lienDL{c'est exactement la question II.1 du DL, avec une autre matrice.}

\Q{P.4} $f(x_0)=Ax_0$ est le 2\textsuperscript{e} vecteur de $\mathcal B$, $f(Ax_0)=A^2x_0$ le 3\textsuperscript{e}, et
$f(A^2x_0)=A^3x_0=0$ :
\[N=\Mat_{\mathcal B}(f)=\begin{pmatrix}0&0&0\\1&0&0\\0&1&0\end{pmatrix},\qquad
P=\begin{pmatrix}1&-1&0\\0&1&1\\0&0&-1\end{pmatrix},\qquad A=PNP^{-1}.\]
$P$ a pour colonnes les coordonnées de $x_0$, $Ax_0$, $A^2x_0$ dans la base canonique. La relation vient de la formule de
changement de base $N=P^{-1}AP$.

\Q{P.5} \choix{« tous les $x$ » : deux sens, et la question P.3 fait le premier.}

Soit $x=(x_1,x_2,x_3)$. On a $A^2x=(x_1+x_2+x_3)(0,1,-1)$.
\begin{itemize}[leftmargin=1.2em,itemsep=1pt]
\item Si $x_1+x_2+x_3\neq0$, alors $A^2x\neq0$ et le raisonnement de P.3 s'applique mot pour mot : la famille est une base.
\item Si $x_1+x_2+x_3=0$, alors $A^2x=0$ : la famille contient le vecteur nul, elle est liée.
\end{itemize}
\textbf{$(x,Ax,A^2x)$ est une base si et seulement si $x_1+x_2+x_3\neq0$}, c'est-à-dire $x\notin\Ker A^2$.
\lienDL{c'est la question II.4.a du DL. La condition ne se devine pas : elle vient du calcul de $A^2x$.}

\Q{P.6} \choix{calculer directement $\det(A+\lambda I_3)$ marche, mais la relation $A=PNP^{-1}$ donne le résultat sans calcul :
deux matrices semblables ont le même déterminant.}

$A+\lambda I_3=PNP^{-1}+\lambda PP^{-1}=P(N+\lambda I_3)P^{-1}$, donc
\[\det(A+\lambda I_3)=\det P\,\det(N+\lambda I_3)\,\det P^{-1}=\det(N+\lambda I_3)=\lambda^3,\]
car $N+\lambda I_3$ est triangulaire, de diagonale $(\lambda,\lambda,\lambda)$. Avec $\lambda=1$ : $\det(A+I_3)=1\neq0$,
\textbf{$A+I_3$ est inversible}. Avec $\lambda=0$ : $\det A=0$, \textbf{$A$ n'est pas inversible} (on le savait : $\rg A=2$).
\gain{on peut même donner l'inverse : $(I_3+A)(I_3-A+A^2)=I_3+A^3=I_3$.}

% ---------------------------------------------------------------------
\clearpage
\markboth{Diagnostic}{}
\subsection*{Après la correction : quoi revoir ?}
\begin{center}
\small\renewcommand{\arraystretch}{1.25}
\begin{tabular}{@{}>{\raggedright\arraybackslash}p{5.6cm}>{\raggedright\arraybackslash}p{7.4cm}@{}}
\toprule
\textbf{Si tu as perdu des points en…} & \textbf{Revoir}\\\midrule
1.2, 2.3(b), 2.4(a), P.3 (liberté) & Livret d'algèbre linéaire, méthode R.2 ; livret du DL, outils 1 et 2\\
2.2, 3.2 (génératrice, base d'un sous-espace) & Livret d'algèbre linéaire, rappels « Espaces vectoriels » et « Dimension finie »\\
2.1(b), 2.4(c), 3.4, P.4 (coordonnées, matrices) & Livret d'algèbre linéaire, proposition « changement de base » ; livret du DL, outil 3\\
3.3, 3.4 (supplémentaires, projecteur) & Livret d'algèbre linéaire, §1.2\\
P.2 (noyau, image, rang) & Livret d'algèbre linéaire, méthode R.3 « Calculer un rang, un noyau, une image »\\
4.1, 4.2 (calculs) & Livret des déterminants, §1.2 et §1.5 (somme constante, Vandermonde)\\
1.5, 4.3, P.6 (propriétés) & Livret des déterminants, §1.3\\
2.1(a), P.5 (« tous les… ») & Livret du DL, outil 4\\\bottomrule
\end{tabular}
\end{center}

\begin{essentiel}
\begin{itemize}[leftmargin=1.1em,itemsep=3pt]
\item Libre : combinaison nulle, puis le bon geste pour isoler les coefficients (identifier, évaluer, composer par $A^k$).
\item $n$ vecteurs en dimension $n$ : une seule propriété à prouver. Pour une base d'un sous-espace $F$, vérifier aussi que les
vecteurs sont dans $F$.
\item Un paramètre et $n$ vecteurs en dimension $n$ : le déterminant donne toutes les valeurs d'un coup.
\item Un déterminant se calcule par opérations sur les lignes et colonnes avant de développer.
\item $\det(\lambda A)=\lambda^n\det A$, $\det A^T=\det A$, deux matrices semblables ont le même déterminant.
\end{itemize}
\end{essentiel}

\else
% =====================================================================
%                               SUJET
% =====================================================================
\markboth{Examen blanc · Algèbre linéaire et déterminants}{}
\begin{center}
\begin{tcolorbox}[enhanced,colback=Fond,colframe=Fond,arc=0pt,boxrule=0pt,left=10pt,right=10pt,top=8pt,bottom=8pt,width=\linewidth]
\small
\textbf{\sffamily\color{Bleu}Consignes.} Le sujet comporte quatre exercices et un problème, indépendants : on peut les traiter
dans l'ordre de son choix. Toute réponse doit être \emph{justifiée} : citer le résultat du cours utilisé et vérifier ses
hypothèses. Les résultats sont à encadrer. Une réponse juste sans justification ne rapporte pas de points ; une démarche correcte
avec une erreur de calcul en rapporte. La clarté et la précision de la rédaction sont prises en compte.

\medskip
\begin{tabular}{@{}lcccccc@{}}
\textbf{\sffamily\color{Bleu}Barème} & Ex.~1 : 3 & Ex.~2 : 5 & Ex.~3 : 3,5 & Ex.~4 : 4 & Problème : 4,5 & \textbf{Total : 20}\\
\end{tabular}
\end{tcolorbox}
\end{center}

\Ex{1}{Vrai ou faux}{3 points}
Pour chaque affirmation, dire si elle est vraie ou fausse, et le justifier : une preuve courte si elle est vraie, un
contre-exemple si elle est fausse. \pts{0,5 point par affirmation}
\begin{enumerate}
\item Toute famille de quatre vecteurs de $\R^3$ est liée.
\item Si $(u,v)$ et $(v,w)$ sont deux familles libres d'un espace vectoriel $E$, alors $(u,v,w)$ est libre.
\item Une famille de trois vecteurs qui engendre $\R^3$ est une base de $\R^3$.
\item La réunion de deux sous-espaces vectoriels de $\R^2$ est un sous-espace vectoriel de $\R^2$.
\item Pour toute matrice $A\in\mathcal M_3(\R)$, $\det(2A)=2\det A$.
\item Une application linéaire de $\R^4$ dans $\R^3$ n'est jamais injective.
\end{enumerate}

\Ex{2}{Familles libres, génératrices, bases}{5 points}

\Partie{A. Une famille de $\R^3$ avec paramètre}
Pour $m\in\R$, on pose $u_1=(1,2,-1)$, $u_2=(0,1,1)$ et $u_3=(2,m,1)$.
\begin{enumerate}
\item
\begin{enumerate}
\item Déterminer les valeurs de $m$ pour lesquelles $(u_1,u_2,u_3)$ est une base de $\R^3$. \pts{0,75 pt}
\item On prend $m=1$ et on note $\mathcal B'=(u_1,u_2,u_3)$. Déterminer les coordonnées du vecteur $(1,1,1)$ dans $\mathcal B'$.
\pts{0,5 pt}
\item Pour la valeur de $m$ exclue au (a), écrire $u_3$ comme combinaison linéaire de $u_1$ et $u_2$, et donner le rang de
$(u_1,u_2,u_3)$. \pts{0,25 pt}
\end{enumerate}
\end{enumerate}

\Partie{B. Un sous-espace de $\R^4$}
Soit $F=\{(x,y,z,t)\in\R^4,\ x+y+z+t=0\ \text{et}\ x-y+2z=0\}$.
\begin{enumerate}[resume]
\item
\begin{enumerate}
\item Déterminer une famille génératrice de $F$, puis montrer que c'est une base de $F$. En déduire $\dim F$. \pts{0,75 pt}
\item Montrer que $w_1=(3,1,-1,-3)$ et $w_2=(0,2,1,-3)$ forment une base de $F$. \pts{0,5 pt}
\end{enumerate}
\end{enumerate}

\Partie{C. Une famille de polynômes}
Dans $\R_2[X]$, on pose $P_1=(X-1)(X-2)$, $P_2=X(X-2)$ et $P_3=X(X-1)$.
\begin{enumerate}[resume]
\item
\begin{enumerate}
\item Calculer $P_i(0)$, $P_i(1)$ et $P_i(2)$ pour $i=1,2,3$. \pts{0,25 pt}
\item Montrer que $(P_1,P_2,P_3)$ est une base de $\R_2[X]$. \pts{0,5 pt}
\item Soit $P\in\R_2[X]$. Exprimer les coordonnées de $P$ dans cette base à l'aide de $P(0)$, $P(1)$ et $P(2)$. \pts{0,25 pt}
\item En déduire l'unique polynôme $P\in\R_2[X]$ tel que $P(0)=1$, $P(1)=0$ et $P(2)=3$. \pts{0,25 pt}
\end{enumerate}
\end{enumerate}

\Partie{D. Un raisonnement sans coordonnées}
Soit $E$ un espace vectoriel de dimension 3 et $(u,v,w)$ une base de $E$.
\begin{enumerate}[resume]
\item
\begin{enumerate}
\item Montrer que $(u+v,\ v+w,\ w+u)$ est une base de $E$. \pts{0,5 pt}
\item La famille $(u-v,\ v-w,\ w-u)$ est-elle une base de $E$ ? \pts{0,25 pt}
\item Donner les coordonnées de $u$ dans la base $(u+v,\ v+w,\ w+u)$. \pts{0,25 pt}
\end{enumerate}
\end{enumerate}

\Ex{3}{Sous-espaces supplémentaires et projecteur}{3,5 points}
Soit $E=\R_3[X]$, de base canonique $\mathcal B_0=(1,X,X^2,X^3)$. On pose
\[F=\{P\in E,\ P(1)=0\ \text{et}\ P'(1)=0\}\qquad\text{et}\qquad G=\R_1[X].\]
On rappelle la formule de Taylor : pour tout $P\in E$,
$P=P(1)+P'(1)(X-1)+\frac{P''(1)}{2}(X-1)^2+\frac{P'''(1)}{6}(X-1)^3$.
\begin{enumerate}
\item Montrer que $F$ est un sous-espace vectoriel de $E$. \pts{0,5 pt}
\item
\begin{enumerate}
\item Montrer que $\mathcal B_1=\big(1,\ X-1,\ (X-1)^2,\ (X-1)^3\big)$ est une base de $E$. \pts{0,5 pt}
\item Donner les coordonnées d'un polynôme $P$ dans $\mathcal B_1$. En déduire une base de $F$ et sa dimension. \pts{0,75 pt}
\end{enumerate}
\item Montrer que $E=F\oplus G$. \pts{0,75 pt}
\item On note $q$ le projecteur sur $G$ parallèlement à $F$.
\begin{enumerate}
\item Montrer que $q(P)=P(1)+P'(1)(X-1)$ pour tout $P\in E$. Calculer $q(X^3)$. \pts{0,5 pt}
\item Écrire la matrice de $q$ dans la base $\mathcal B_1$, puis dans la base $\mathcal B_0$. \pts{0,5 pt}
\end{enumerate}
\end{enumerate}

\Ex{4}{Calculs de déterminants}{4 points}
\begin{enumerate}
\item Soient $a,b\in\R$. Calculer, sous forme factorisée, \pts{1,5 pt}
\[D_4=\begin{vmatrix}a&b&b&b\\b&a&b&b\\b&b&a&b\\b&b&b&a\end{vmatrix}.\]
Pour quelles valeurs de $a$ et $b$ la matrice correspondante est-elle inversible ?
\item Soient $a,b,c$ trois réels.
\begin{enumerate}
\item Calculer, sous forme factorisée, $V=\begin{vmatrix}1&a&a^2\\1&b&b^2\\1&c&c^2\end{vmatrix}$. \pts{1 pt}
\item On suppose $a,b,c$ deux à deux distincts. Soient $\alpha,\beta,\gamma\in\R$. Montrer qu'il existe un unique polynôme
$P\in\R_2[X]$ tel que $P(a)=\alpha$, $P(b)=\beta$ et $P(c)=\gamma$. \pts{0,5 pt}
\end{enumerate}
\item Soit $n$ un entier impair et $A\in\mathcal M_n(\R)$ une matrice antisymétrique ($A^T=-A$). Montrer que $A$ n'est pas
inversible. Le résultat reste-t-il vrai si $n$ est pair ? \pts{1 pt}
\end{enumerate}

\Pb{une matrice nilpotente}{4,5 points}
On considère
\[A=\begin{pmatrix}-1&-1&-1\\1&2&2\\0&-1&-1\end{pmatrix}\]
et $f$ l'endomorphisme de $\R^3$ canoniquement associé à $A$. On note $(e_1,e_2,e_3)$ la base canonique de $\R^3$.
\begin{enumerate}
\item Calculer $A^2$ et $A^3$. \pts{0,5 pt}
\item Déterminer une base de $\Ker A$ et une base de $\Im A$. Vérifier le théorème du rang, donner une équation de $\Im A$, et
montrer que $\Ker A\subset\Im A$. \pts{1 pt}
\item On pose $x_0=e_1$. Montrer, \emph{sans calculer de déterminant}, que $\mathcal B=(x_0,Ax_0,A^2x_0)$ est une base de
$\R^3$. \pts{1 pt}
\item Écrire la matrice $N$ de $f$ dans la base $\mathcal B$, la matrice de passage $P$ de la base canonique à $\mathcal B$, et
la relation entre $A$, $N$ et $P$. \pts{0,5 pt}
\item Déterminer tous les vecteurs $x=(x_1,x_2,x_3)$ de $\R^3$ tels que $(x,Ax,A^2x)$ soit une base de $\R^3$. \pts{1 pt}
\item Pour $\lambda\in\R$, calculer $\det(A+\lambda I_3)$ en utilisant la question 4. En déduire que $A+I_3$ est inversible et
que $A$ ne l'est pas. \pts{0,5 pt}
\end{enumerate}

\vfill
\begin{center}\sffamily\small\color{Gris}— Fin du sujet —\end{center}
\fi
\end{document}
"
\documentclass[a4paper,11pt]{article}

% ============================================================
%  Le Collé — Algèbre linéaire (compléments de deuxième année)
%  © 2026 Ahmed Amine Ouqlani
% ============================================================

% ---------- Polices ----------
\usepackage[utf8]{inputenc}
\usepackage[T1]{fontenc}
\usepackage[sc]{mathpazo}          % texte et mathématiques : Palatino
\usepackage[scaled=0.88]{helvet}   % titres et repères : linéale
\linespread{1.07}
\usepackage{amsmath,amssymb,amsthm,mathrsfs,stmaryrd}
\usepackage{microtype}

% ---------- Mise en page ----------
\usepackage[a4paper,top=2.3cm,bottom=2.3cm,left=2.4cm,right=2.4cm,
            headheight=14pt,headsep=0.8cm,footskip=1.1cm]{geometry}
\setlength{\parindent}{0pt}
\setlength{\parskip}{0.45em plus 0.15em minus 0.1em}
\raggedbottom
\emergencystretch=1.5em
\clubpenalty=10000 \widowpenalty=10000 \displaywidowpenalty=10000 \brokenpenalty=10000

% ---------- Couleurs : un bleu de nuit, un ocre, des gris ----------
\usepackage{xcolor}
\definecolor{Ink}{HTML}{1C2533}
\definecolor{Bleu}{HTML}{1F3A5F}
\definecolor{BleuMoyen}{HTML}{4A6A8F}
\definecolor{BleuClair}{HTML}{9DB2C9}
\definecolor{Fond}{HTML}{F2F4F7}
\definecolor{Ocre}{HTML}{B07A26}
\definecolor{Brique}{HTML}{A23B2A}
\definecolor{Gris}{HTML}{5E6672}
\definecolor{Filet}{HTML}{C5CCD6}

% ---------- Packages ----------
\usepackage{enumitem}
\setlist{itemsep=0.25em,topsep=0.35em,parsep=0pt}
\setlist[itemize]{label=\textcolor{BleuMoyen}{\small\textbullet},leftmargin=1.4em}
\setlist[enumerate]{leftmargin=1.8em,label=\textsf{\small\color{Bleu}\arabic*.}}
\usepackage{booktabs,array,multicol}
\usepackage{needspace,float}
\usepackage[font=small,labelfont={sf,color=Bleu},labelsep=period,skip=6pt]{caption}
\usepackage{tikz}
\usetikzlibrary{arrows.meta,calc,decorations.markings}
\usepackage{pgfplots}
\pgfplotsset{compat=1.17}
\usepackage{tkz-tab}
\usepackage[most]{tcolorbox}
\usepackage{titlesec,titletoc}
\usepackage{fancyhdr,lastpage,etoolbox}
\usepackage[colorlinks=true,linkcolor=Bleu,urlcolor=Bleu,citecolor=Bleu,
            pdftitle={Examen blanc — Algèbre linéaire et déterminants},pdfauthor={Ahmed Amine Ouqlani}]{hyperref}

\numberwithin{equation}{section}
\numberwithin{figure}{section}
\numberwithin{table}{section}

% ---------- Titres ----------
% numéro du chapitre dans le parcours PT (affiché), distinct de la numérotation interne
\newcommand{\grandnum}{\thesection}\newcommand{\ongletcouleur}{Bleu}
% Chapitres (sections) : grand numéro, titre en linéale, filet ocre
\titleformat{\section}[display]
  {\normalfont\sffamily\color{Ink}}
  {{\small\color{Gris}\textls[200]{\MakeUppercase{Chapitre}}}\\[-0.2em]%
   {\fontsize{64}{64}\selectfont\color{BleuClair}\grandnum}}
  {0.4em}
  {\fontsize{25}{30}\selectfont\bfseries}
  [\vspace{0.5em}{\color{Ocre}\rule{3.2cm}{1.4pt}}]
\titlespacing*{\section}{0pt}{0pt}{1.6em}
\titleformat{name=\section,numberless}[display]
  {\normalfont\sffamily\color{Ink}}{}{0pt}{\fontsize{25}{30}\selectfont\bfseries}
  [\vspace{0.5em}{\color{Ocre}\rule{3.2cm}{1.4pt}}]
\titlespacing*{name=\section,numberless}{0pt}{0pt}{1.6em}

\titleformat{\subsection}
  {\needspace{5\baselineskip}\normalfont\sffamily\Large\bfseries\color{Ink}}
  {\color{Bleu}\thesubsection}{0.8em}{}
\titlespacing*{\subsection}{0pt}{1.8em}{0.5em}
\titleformat{\subsubsection}
  {\needspace{5\baselineskip}\normalfont\sffamily\large\bfseries\color{Ink}}
  {\color{Bleu}\thesubsubsection}{0.7em}{}
\titlespacing*{\subsubsection}{0pt}{1.6em}{0.4em}
\titleformat{\paragraph}[runin]{\normalfont\sffamily\bfseries\color{Ink}}{}{0pt}{}[.]
\titlespacing*{\paragraph}{0pt}{0.9em}{0.6em}
\setcounter{secnumdepth}{3}
\setcounter{tocdepth}{2}

% ---------- Table des matières ----------
\titlecontents{section}[0pt]{\addvspace{0.7em}\sffamily\bfseries\color{Ink}}
  {}{}{\hfill\contentspage}[\addvspace{0.2em}]
\titlecontents{subsection}[2.2em]{\small\linespread{1}\selectfont}
  {\makebox[2.6em][l]{\color{Gris}\thecontentslabel}}{}{\titlerule*[0.7pc]{.}\contentspage}

% ---------- En-têtes et pieds de page ----------
\pagestyle{fancy}
\fancyhf{}
\renewcommand{\sectionmark}[1]{\markboth{#1}{}}
\fancyhead[L]{\footnotesize\sffamily\color{Gris}\leftmark}
\fancyhead[R]{\footnotesize\sffamily\scshape\color{Gris}Le Collé}
\fancyfoot[C]{\small\sffamily\color{Ink}\thepage}
\fancyfoot[L]{\scriptsize\sffamily\hyperlink{toc}{\color{Filet}sommaire}}
\renewcommand{\headrulewidth}{0.4pt}
\renewcommand{\headrule}{{\color{Filet}\hrule height\headrulewidth width\headwidth}}
\renewcommand{\footrulewidth}{0pt}
\fancypagestyle{plain}{\fancyhf{}\fancyfoot[C]{\small\sffamily\color{Ink}\thepage}\renewcommand{\headrulewidth}{0pt}}

% ---------- Énoncés ----------
\newtheoremstyle{enonce}{0pt}{0pt}{}{}{\sffamily\bfseries\color{Bleu}}{.}{0.5em}
  {\thmname{#1}\thmnumber{~#2}\thmnote{\ \normalfont\sffamily\color{Ink}(#3)}}
\newtheoremstyle{exemple}{0.9em}{0.9em}{}{}{\sffamily\bfseries\color{BleuMoyen}}{.}{0.5em}
  {\thmname{#1}\thmnumber{~#2}\thmnote{\ \normalfont\sffamily\color{Ink}(#3)}}
\newtheoremstyle{remarque}{0.7em}{0.7em}{}{}{\sffamily\itshape}{.}{0.5em}{\thmname{#1}\thmnote{\ (#3)}}
\theoremstyle{enonce}
\newtheorem{definition}{Définition}[section]
\newtheorem{theoreme}[definition]{Théorème}
\newtheorem{proposition}[definition]{Proposition}
\newtheorem{corollaire}[definition]{Corollaire}
\theoremstyle{exemple}
\newtheorem{exemple}[definition]{Exemple}
\theoremstyle{remarque}
\newtheorem*{remarque}{Remarque}

\tcbset{
  defi/.style={enhanced,breakable,colback=white,colframe=white,boxrule=0pt,arc=0pt,
     borderline west={1.2pt}{0pt}{Bleu},left=11pt,right=4pt,top=3pt,bottom=3pt,
     before skip=1.1em,after skip=1.1em},
  thm/.style={enhanced,breakable,colback=Fond,colframe=Fond,boxrule=0pt,arc=0pt,
     borderline west={2.4pt}{0pt}{Bleu},left=11pt,right=8pt,top=6pt,bottom=6pt,
     before skip=1.2em,after skip=1.2em}}
\tcolorboxenvironment{definition}{defi}
\tcolorboxenvironment{theoreme}{thm}
\tcolorboxenvironment{proposition}{thm}
\tcolorboxenvironment{corollaire}{thm}

\renewcommand{\proofname}{Démonstration}
\renewcommand{\qedsymbol}{\textcolor{BleuClair}{$\blacksquare$}}
\makeatletter
\renewenvironment{proof}[1][\proofname]{\par\pushQED{\qed}\normalfont
  \topsep0.5em\relax\trivlist\item[\hskip\labelsep\sffamily\itshape\color{Gris}#1.]\ignorespaces}
  {\popQED\endtrivlist\@endpefalse}
\makeatother

% ---------- Méthode : filets haut et bas, titre en bandeau ----------
\makeatletter\newcommand{\methlabelfix}{\protected@edef\@currentlabel{\thetcbcounter}}\makeatother
\newtcolorbox[auto counter,number within=section]{methode}[2][]{%
  enhanced,unbreakable,colback=white,colframe=white,boxrule=0pt,arc=0pt,
  borderline north={0.9pt}{0pt}{Bleu},borderline south={0.9pt}{0pt}{Bleu},
  left=4pt,right=4pt,top=8pt,bottom=8pt,before skip=1.4em,after skip=1.4em,
  before upper={\methlabelfix\phantomsection\addcontentsline{met}{metentry}{\protect\makebox[2.6em][l]{\thetcbcounter}#2}%
     {\sffamily\bfseries\small\color{white}\colorbox{Bleu}{\strut\ MÉTHODE~\thetcbcounter\ }}%
     \ \ {\sffamily\bfseries\color{Bleu}#2}\par\smallskip},#1}

% ---------- Repères pédagogiques : un filet de couleur, un mot-clé ----------
\newcommand{\definirbloc}[3]{% nom, couleur, libellé
  \newtcolorbox{#1}[1][]{enhanced,breakable,blanker,
    borderline west={1.4pt}{0pt}{#2},left=10pt,right=2pt,top=1pt,bottom=1pt,
    before skip=1em,after skip=1em,
    before upper={{\sffamily\bfseries\footnotesize\color{#2}\textls[120]{\MakeUppercase{#3}}}\quad},##1}}
\definirbloc{reflexe}{Bleu}{Réflexe}
\definirbloc{astuce}{Ocre}{Astuce}
\definirbloc{attention}{Brique}{Attention}
\definirbloc{enpratique}{BleuMoyen}{En pratique}
\definirbloc{plusloin}{Gris}{Pour aller plus loin}
\definirbloc{strategie}{Bleu}{Stratégie}
\definirbloc{lecture}{BleuMoyen}{Lire le schéma}
% compatibilité avec les parties reprises
\newenvironment{garde}{\begin{attention}}{\end{attention}}
\newenvironment{complementHP}{\begin{plusloin}}{\end{plusloin}}
\newenvironment{interpretation}{\begin{enpratique}}{\end{enpratique}}

% ---------- Exercices, corrigés ----------
\newcounter{exo}[section]
\renewcommand{\theexo}{\thesection.\arabic{exo}}
\newcommand{\niveau}[1]{{\color{Ocre}\footnotesize
  \ifcase#1\or$\bullet\circ\circ\circ\circ$\or$\bullet\bullet\circ\circ\circ$\or$\bullet\bullet\bullet\circ\circ$%
  \or$\bullet\bullet\bullet\bullet\circ$\or$\bullet\bullet\bullet\bullet\bullet$\fi}}
\newcommand{\serie}[1]{\par\needspace{8\baselineskip}\vspace{1.5em}%
  \noindent{\sffamily\bfseries\color{Bleu}#1}\par\nobreak\vspace{-0.4em}%
  \noindent{\color{Filet}\rule{\linewidth}{0.4pt}}\par\nobreak\vspace{0.1em}}
\newenvironment{exo}[3]{%
  \par\needspace{5\baselineskip}\vspace{0.75em}\refstepcounter{exo}\label{e:#1}%
  \def\exolab{#1}%
  \noindent{\sffamily\bfseries\color{Bleu}Exercice~\theexo}\quad\niveau{#2}%
  \hfill{\footnotesize\sffamily\color{Gris}#3}\par\nobreak\vspace{0.15em}}%
  {\par\nopagebreak\noindent\hfill{\scriptsize\sffamily\hyperref[c:\exolab]{\color{Gris}corrigé p.~\pageref{c:\exolab}}}\par}
\newenvironment{corrige}[1]{%
  \par\needspace{6\baselineskip}\vspace{1em}\phantomsection\label{c:#1}%
  \noindent{\sffamily\bfseries\color{Bleu}Exercice~\ref{e:#1}}\quad
  {\scriptsize\sffamily\hyperref[e:#1]{\color{Gris}énoncé p.~\pageref{e:#1}}}\par\nobreak\vspace{0.2em}}%
  {\hfill\textcolor{BleuClair}{$\square$}\par}
\newcommand{\etape}[1]{\par\smallskip\noindent{\sffamily\bfseries\small\color{BleuMoyen}#1.}\ }
\newenvironment{solution}{\par\noindent{\sffamily\bfseries\small\color{BleuMoyen}Solution.}\ }{\hfill\textcolor{BleuClair}{$\square$}\par}
\newcommand{\chercher}{\par\smallskip\noindent{\color{Filet}\rule[0.5ex]{1.2cm}{0.4pt}}\
  {\footnotesize\sffamily\itshape\color{Gris}laisser chercher l'étudiant avant de poursuivre}\par\smallskip}

% ---------- Guide de séance ----------

\newcommand{\marge}[1]{\leavevmode\marginpar{\raggedleft\scriptsize\sffamily\color{Gris}\textls[100]{\MakeUppercase{#1}}}\ignorespaces}
\newcommand{\Etape}[4]{%
  \subsection[Étape #1 --- #3 {\color{Gris}(#2)}]{#3}\label{etape:#1}%
  \noindent{\small\sffamily\color{Gris}{\color{Ocre}\bfseries Étape #1}\quad·\quad #2\quad·\quad #4}\par\medskip}
\newcommand{\question}[1]{\par\smallskip\noindent{\sffamily\bfseries\small\color{BleuMoyen}Question.}\ \textit{#1}\par\smallskip}
\newcommand{\rappel}[1]{\par\smallskip\noindent\phantomsection\addcontentsline{rap}{rapentry}{#1}%
  {\sffamily\bfseries\scriptsize\colorbox{Fond}{\color{Gris}\strut RAPPEL}}\ {\sffamily\small\color{Gris}#1.}\ }
\newcommand{\pourquoi}{\par\noindent{\sffamily\bfseries\small\color{Ocre}Pourquoi ?}\ }

% Ouverture de chapitre : \ouverture{objectifs}{prérequis}{durée}
\newcommand{\ouverture}[3]{%
  \begin{tcolorbox}[enhanced,colback=Fond,colframe=Fond,arc=0pt,boxrule=0pt,left=12pt,right=12pt,top=10pt,bottom=10pt,
     before skip=1.2em,after skip=1.6em]
  \begin{minipage}[t]{0.56\linewidth}
  {\sffamily\bfseries\small\color{Bleu}OBJECTIFS}\par\vspace{0.3em}\small #1
  \end{minipage}\hfill
  \begin{minipage}[t]{0.38\linewidth}
  {\sffamily\bfseries\small\color{Bleu}PRÉREQUIS}\par\vspace{0.3em}\small #2
  \ifx\relax#3\relax\else\par\vspace{0.6em}{\sffamily\bfseries\small\color{Bleu}DURÉE}\par\vspace{0.2em}\small #3\fi
  \end{minipage}
  \end{tcolorbox}}
% Fin de chapitre
\newenvironment{essentiel}{\par\needspace{12\baselineskip}\vspace{1.5em}%
  \begin{tcolorbox}[enhanced,unbreakable,colback=white,colframe=white,boxrule=0pt,arc=0pt,
   borderline north={1.4pt}{0pt}{Ocre},left=2pt,right=2pt,top=8pt,bottom=4pt]
  {\sffamily\bfseries\color{Ocre}L'ESSENTIEL DU CHAPITRE}\par\medskip}{\end{tcolorbox}}

% ---------- Notations ----------
\newcommand{\K}{\mathbb{K}}
\newcommand{\R}{\mathbb{R}}
\newcommand{\C}{\mathbb{C}}
\newcommand{\N}{\mathbb{N}}
\newcommand{\Z}{\mathbb{Z}}
\newcommand{\Mn}[1]{\mathcal{M}_{n}(#1)}
\newcommand{\GLn}[1]{\mathrm{GL}_n(#1)}
\newcommand{\LE}{\mathcal{L}(E)}
\newcommand{\rg}{\operatorname{rg}}
\newcommand{\vect}{\operatorname{Vect}}
\newcommand{\Mat}{\operatorname{Mat}}
\newcommand{\Id}{\operatorname{Id}}
\newcommand{\sh}{\operatorname{sh}}
\newcommand{\ch}{\operatorname{ch}}
\renewcommand{\th}{\operatorname{th}}
\DeclareMathOperator{\argsh}{argsh}
\newcommand{\vv}[1]{\overrightarrow{#1}}
\newcommand{\ve}[1]{\vec{#1}}
\newcommand{\ient}[2]{\llbracket #1,#2\rrbracket}
\newcommand{\dt}[1]{\left|\begin{matrix}#1\end{matrix}\right|}
\newcommand{\dis}{\displaystyle}

% ---------- Figures ----------
\pgfplotsset{trig format plots=rad}
\pgfplotsset{courbe/.style={axis lines=middle,axis equal image,
  axis line style={Gris,-{Stealth[length=5pt]}},
  xlabel={$x$},ylabel={$y$},ticklabel style={font=\scriptsize\sffamily,color=Gris},
  every axis x label/.style={at={(ticklabel* cs:1)},anchor=west,font=\small},
  every axis y label/.style={at={(ticklabel* cs:1)},anchor=south,font=\small},
  samples=300,smooth,width=0.6\linewidth}}
\tikzset{arcetude/.style={Bleu,line width=1.25pt},
         reste/.style={BleuClair,line width=1pt},
         tgte/.style={Ocre,line width=0.9pt},
         asy/.style={Gris,dashed,line width=0.6pt},
         pt/.style={circle,fill=Ink,inner sep=1.2pt}}

% ---------- Navigation ----------
\makeatletter
\newcommand*\l@rapentry{\@dottedtocline{1}{0em}{0em}}
\newcommand*\l@metentry{\@dottedtocline{1}{0em}{0em}}
\newcommand{\listemethodes}{\@starttoc{met}}
\newcommand{\listerappels}{\@starttoc{rap}}
\makeatother
% onglets de chapitre sur le bord droit
\usepackage{eso-pic}
\newif\ifonglets \ongletsfalse
\newcommand{\dessineronglet}{%
  \edef\ongcol{\ongletcouleur}%
  \edef\ongnum{\grandnum}%
  \pgfmathsetmacro{\ongy}{3.2+2.3*(\ongnum-1)}%
  \begin{tikzpicture}[remember picture,overlay]
  \fill[\ongcol] ([xshift=-0.75cm,yshift=-\ongy cm]current page.north east)
     rectangle ([yshift=-\ongy cm-1.9cm]current page.north east);
  \node[text=white,font=\sffamily\bfseries\large] at ([xshift=-0.37cm,yshift=-\ongy cm-0.95cm]current page.north east)
     {\ongnum};
  \end{tikzpicture}}
\AddToShipoutPictureBG{\ifonglets\ifnum\value{section}>0 \dessineronglet\fi\fi}
% blocs de navigation et d'enseignement
\definirbloc{commencer}{Gris}{Si vous commencez ici}
\definirbloc{horsprogramme}{Brique}{Hors programme}
\newcommand{\HP}{\,{\sffamily\bfseries\scriptsize\textcolor{Brique}{[HP]}}}
\newcommand{\source}[1]{{\footnotesize\sffamily\color{Gris}\,(#1)}}
\newcommand{\coupdepouce}[1]{\par\vspace{-0.2em}\noindent{\footnotesize\sffamily\color{Gris}%
  \textbf{Coup de pouce} (à donner si l'étudiant bloque) : #1}\par}
\newcommand{\pointarret}[1]{\par\vspace{1em}\noindent{\color{Filet}\rule[0.5ex]{\linewidth}{0.4pt}}\par\nobreak
  \noindent{\footnotesize\sffamily\color{Ocre}\bfseries POINT D'ARRÊT POSSIBLE}\quad{\footnotesize\sffamily\color{Gris}#1}\par\vspace{0.6em}}
\newcommand{\demarche}[1]{\noindent{\sffamily\bfseries\small\color{BleuMoyen}Démarche.}\ \textit{#1}\par\smallskip}

% ---------- Corrigés : idées, astuces, pièges, à retenir ----------
\newcommand{\choix}[1]{\par\noindent{\sffamily\bfseries\small\color{BleuMoyen}Choix de la méthode.}\ #1\par\smallskip}
\newcommand{\gain}[1]{\par\noindent{\sffamily\bfseries\small\color{Ocre}Astuce.}\ #1\par}
\newcommand{\piege}[1]{\par\noindent{\sffamily\bfseries\small\color{Brique}Piège.}\ #1\par}
\newtcolorbox{retenirbox}{enhanced,breakable,colback=Fond,colframe=Fond,arc=0pt,boxrule=0pt,
  borderline west={2pt}{0pt}{Ocre},left=8pt,right=6pt,top=4pt,bottom=4pt,before skip=0.7em,after skip=0.5em}
\newcommand{\retenir}[1]{\begin{retenirbox}{\sffamily\bfseries\footnotesize\color{Ocre}À RETENIR}\quad #1\end{retenirbox}}
\newcommand{\TD}[1]{{\color{Ocre}\bfseries TD\,#1}\ ·\ }
\newcommand{\q}[1]{\par\smallskip\noindent{\sffamily\bfseries\color{Bleu}#1}\ }
\providecommand{\iddots}{\mathinner{\mkern1mu\raise1pt\hbox{.}\mkern2mu\raise4pt\hbox{.}\mkern2mu\raise7pt\hbox{.}\mkern1mu}}


% =====================================================================
%  Examen blanc — Algèbre linéaire et déterminants (Le Collé, PT)
%  Une source, deux PDF : \def\CORRIGE{} avant \input donne le corrigé.
% =====================================================================
\providecommand{\Ker}{\operatorname{Ker}}\renewcommand{\Im}{\operatorname{Im}}
\fancyfoot[L]{}\fancyhead[L]{\footnotesize\sffamily\color{Gris}Examen blanc · Algèbre linéaire et déterminants\ifdefined\CORRIGE\ · corrigé\fi}
\newcommand{\pts}[1]{\hfill{\footnotesize\sffamily\color{Gris}(#1)}}
\newcommand{\Ex}[3]{\par\needspace{6\baselineskip}\vspace{1.3em}\noindent
  {\sffamily\bfseries\large\color{Bleu}Exercice #1 \,·\, #2}\hfill{\sffamily\small\color{Ocre}#3}\par
  \vspace{2pt}{\color{Filet}\rule{\linewidth}{0.4pt}}\par\smallskip}
\newcommand{\Pb}[2]{\par\needspace{6\baselineskip}\vspace{1.3em}\noindent
  {\sffamily\bfseries\large\color{Bleu}Problème \,·\, #1}\hfill{\sffamily\small\color{Ocre}#2}\par
  \vspace{2pt}{\color{Filet}\rule{\linewidth}{0.4pt}}\par\smallskip}
\newcommand{\Partie}[1]{\par\needspace{4\baselineskip}\medskip\noindent{\sffamily\bfseries\color{BleuMoyen}#1}\par\nobreak\smallskip}
\newcommand{\Q}[1]{\par\needspace{3\baselineskip}\medskip\noindent{\sffamily\bfseries\color{Bleu}#1}\enspace}
\newcommand{\autre}[1]{\par\noindent{\sffamily\bfseries\small\color{BleuMoyen}Autre méthode.}\ #1\par}
\newcommand{\lienDL}[1]{\par\noindent{\sffamily\bfseries\small\color{Ocre}Lien avec le DL.}\ #1\par}
\setlist[enumerate,1]{label=\arabic*.,leftmargin=1.6em,itemsep=4pt,topsep=4pt}
\setlist[enumerate,2]{label=(\alph*),leftmargin=1.6em,itemsep=2pt,topsep=2pt}

\begin{document}
\pagestyle{fancy}
\thispagestyle{plain}

% ------------------------------------------------------------------ en-tête
\begin{tikzpicture}[remember picture,overlay]
\fill[Bleu] (current page.north west) rectangle ([yshift=-4.2cm]current page.north east);
\node[anchor=north west,text=white,font=\sffamily\footnotesize,inner sep=0] at ([xshift=2.2cm,yshift=-1.1cm]current page.north west)
  {\textls[120]{MATHÉMATIQUES · CPGE PT · EXAMEN BLANC}};
\node[anchor=north west,text=white,font=\sffamily\bfseries\LARGE,inner sep=0] at ([xshift=2.2cm,yshift=-1.75cm]current page.north west)
  {Algèbre linéaire et déterminants};
\node[anchor=north west,text=white,font=\sffamily\normalsize,inner sep=0] at ([xshift=2.2cm,yshift=-2.75cm]current page.north west)
  {\ifdefined\CORRIGE Corrigé détaillé et diagnostic\else Durée : 3 heures \quad·\quad Calculatrice interdite\fi};
\node[anchor=north east,text=white,font=\sffamily\scshape\large,inner sep=0] at ([xshift=-2.2cm,yshift=-1.05cm]current page.north east)
  {Le Collé};
\end{tikzpicture}
\vspace*{2.9cm}

\ifdefined\CORRIGE
% =====================================================================
%                              CORRIGÉ
% =====================================================================
\markboth{Corrigé · Examen blanc}{}
\paragraph{Comment utiliser ce corrigé} Ne l'ouvre qu'après avoir composé les 3 heures. Corrige d'abord ta copie avec un
stylo d'une autre couleur, question par question, en comptant les points. Pour chaque question ratée, lis le \emph{choix de la
méthode} : c'est lui qu'il faut retenir, plus que le calcul. La grille de la dernière page dit quoi revoir.

\subsection*{Ce que le sujet évalue}
\begin{center}
\small\renewcommand{\arraystretch}{1.22}
\begin{tabular}{@{}>{\raggedright\arraybackslash}p{7.6cm}>{\raggedright\arraybackslash}p{5.4cm}@{}}
\toprule
\textbf{Compétence} & \textbf{Questions}\\\midrule
Démontrer qu'une famille est libre (combinaison nulle, évaluation, composition) & 1.2 · 2.3(b) · 2.4(a) · P.3\\
Démontrer qu'une famille est génératrice, trouver une base d'un sous-espace & 1.3 · 2.2(a) · 3.2\\
Utiliser la dimension (« $n$ vecteurs en dimension $n$ ») & 1.1 · 2.2(b) · 2.4(a) · 3.3 · P.3\\
Coordonnées dans une base, changement de base & 2.1(b) · 2.3(c) · 2.4(c) · 3.4 · P.4\\
Sous-espaces, somme directe, supplémentaires, projecteur & 1.4 · 3.1 · 3.3 · 3.4\\
Application linéaire : noyau, image, rang, théorème du rang & 1.6 · P.2\\
Calculer un déterminant (opérations, développement, factorisation) & 2.1(a) · 4.1 · 4.2\\
Propriétés du déterminant ($\det(\lambda A)$, $\det A^T$, $\det(AB)$) & 1.5 · 4.3 · P.6\\
Déterminant, inversibilité, liberté, systèmes de Cramer & 2.1 · 4.1 · 4.2 · P.6\\
« Trouver tous les… » : prouver les deux sens & 2.1(a) · P.5\\\bottomrule
\end{tabular}
\end{center}

\lienDL{Les questions 2.3(b), 4.2 et P.3 à P.5 reprennent les difficultés du DL~3 sous une autre forme : liberté par évaluation
(comme III.3), Vandermonde, famille $(x,Ax,A^2x)$ avec $A^3=0$ (comme II.1 et II.4.a). Les formulations sont différentes : il
faut reconnaître la technique.}

\begin{center}\small\renewcommand{\arraystretch}{1.2}
\begin{tabular}{@{}lcccccc@{}}\toprule
& Ex.~1 & Ex.~2 & Ex.~3 & Ex.~4 & Problème & Total\\\midrule
Points & 3 & 5 & 3,5 & 4 & 4,5 & 20\\
Temps conseillé & 15 min & 50 min & 35 min & 35 min & 45 min & 3 h\\\bottomrule
\end{tabular}
\end{center}

% ---------------------------------------------------------------------
\clearpage
\Ex{1}{Vrai ou faux}{3 points}

\Q{1. Vrai.} $\dim\R^3=3$ : une famille libre de $\R^3$ a au plus 3 vecteurs. Une famille de 4 vecteurs est donc liée.

\Q{2. Faux.} Contre-exemple : $u=(1,0,0)$, $v=(0,1,0)$, $w=(1,0,0)$. Les familles $(u,v)$ et $(v,w)$ sont libres, mais $u-w=0$ :
$(u,v,w)$ est liée.
\piege{la liberté ne se vérifie pas « deux à deux ». Trois vecteurs deux à deux non colinéaires peuvent être liés, par exemple
$(1,0)$, $(0,1)$, $(1,1)$ dans $\R^2$.}

\Q{3. Vrai.} $\dim\R^3=3$ et la famille a 3 vecteurs : pour une telle famille, génératrice $\iff$ libre $\iff$ base.

\Q{4. Faux.} Dans $\R^2$, $F=\vect\big((1,0)\big)$ et $G=\vect\big((0,1)\big)$ : $(1,0)$ et $(0,1)$ sont dans $F\cup G$, mais pas
leur somme $(1,1)$. $F\cup G$ n'est pas stable par addition.

\Q{5. Faux.} Multiplier $A$ par $2$ multiplie chacune des 3 colonnes par $2$, et le déterminant est linéaire par rapport à
\emph{chaque} colonne : $\det(2A)=2^3\det A=8\det A$. Par exemple $\det(2I_3)=8$.
\retenir{$\det(\lambda A)=\lambda^n\det A$ pour $A\in\Mn{\K}$.}

\Q{6. Vrai.} Par le théorème du rang, $\dim\Ker f=4-\rg f$. Or $\Im f\subset\R^3$, donc $\rg f\leq3$ et $\dim\Ker f\geq1$ :
$\Ker f\neq\{0\}$, $f$ n'est pas injective.

% ---------------------------------------------------------------------
\Ex{2}{Familles libres, génératrices, bases}{5 points}

\Partie{A. Une famille de $\R^3$ avec paramètre}

\Q{2.1(a)} \choix{3 vecteurs en dimension 3 avec un paramètre : un seul déterminant donne toutes les valeurs de $m$ d'un coup.
Résoudre un système avec $m$ obligerait à discuter à chaque pivot.}
\[\det(u_1,u_2,u_3)=\begin{vmatrix}1&0&2\\2&1&m\\-1&1&1\end{vmatrix}
\underset{C_3\leftarrow C_3-2C_1}{=}\begin{vmatrix}1&0&0\\2&1&m-4\\-1&1&3\end{vmatrix}=1\times\begin{vmatrix}1&m-4\\1&3\end{vmatrix}=7-m.\]
Le déterminant est non nul si et seulement si $m\neq7$. \textbf{$(u_1,u_2,u_3)$ est une base de $\R^3$ si et seulement si $m\neq7$.}
Les deux sens sont prouvés en même temps, car « base $\iff\det\neq0$ ».

\Q{2.1(b)} Pour $m=1$, on cherche $(a,b,c)$ tel que $au_1+bu_2+cu_3=(1,1,1)$ :
\[\begin{cases}a+2c=1\\2a+b+c=1\\-a+b+c=1\end{cases}\]
La différence des deux dernières lignes donne $3a=0$, donc $a=0$, puis $c=\frac12$ et $b=\frac12$. \textbf{Les coordonnées de
$(1,1,1)$ dans $\mathcal B'$ sont $\left(0,\frac12,\frac12\right)$.} Vérification : $\frac12(u_2+u_3)=\frac12(2,2,2)=(1,1,1)$.
\piege{les coordonnées dépendent de la base : $(1,1,1)$ est la colonne dans la base canonique, $(0,\frac12,\frac12)$ dans
$\mathcal B'$. Toujours vérifier en recombinant.}

\Q{2.1(c)} Pour $m=7$, on cherche une relation $au_1+bu_2+u_3=0$ : la première coordonnée donne $a=-2$, la troisième $b=-3$, et
la deuxième est alors vérifiée. Donc $u_3=2u_1+3u_2$. Comme $u_1$ et $u_2$ ne sont pas colinéaires,
$\rg(u_1,u_2,u_3)=\dim\vect(u_1,u_2)=2$.

\Partie{B. Un sous-espace de $\R^4$}

\Q{2.2(a)} \choix{un sous-espace défini par des équations : on résout le système, et on écrit la solution générale comme
combinaison linéaire. Cela donne directement une famille génératrice.}

$(x,y,z,t)\in F\iff x=y-2z$ et $t=-x-y-z=-2y+z$. Donc
\[F=\{\,y(1,1,0,-2)+z(-2,0,1,1),\ y,z\in\R\,\}=\vect(f_1,f_2),\]
avec $f_1=(1,1,0,-2)$ et $f_2=(-2,0,1,1)$.
\emph{Génératrice :} c'est l'égalité ci-dessus. \emph{Libre :} si $yf_1+zf_2=0$, les deuxième et troisième coordonnées donnent
$y=0$ et $z=0$. \textbf{$(f_1,f_2)$ est une base de $F$, et $\dim F=2$.}
\autre{$F$ est défini par 2 équations indépendantes dans $\R^4$, donc $\dim F=4-2=2$ (théorème du rang appliqué à
$(x,y,z,t)\mapsto(x+y+z+t,\,x-y+2z)$, de rang 2).}

\Q{2.2(b)} \choix{on connaît maintenant $\dim F=2$. Pour 2 vecteurs dans un espace de dimension 2, il suffit qu'ils soient
dans $F$ et libres : inutile de prouver qu'ils engendrent $F$.}

$w_1=(3,1,-1,-3)$ : $3+1-1-3=0$ et $3-1-2=0$, donc $w_1\in F$. $w_2=(0,2,1,-3)$ : $0+2+1-3=0$ et $0-2+2=0$, donc $w_2\in F$. Ils
ne sont pas colinéaires (la première coordonnée de $w_2$ est nulle, pas celle de $w_1$, et $w_2\neq0$). \textbf{$(w_1,w_2)$ est une
famille libre de 2 vecteurs de $F$, avec $\dim F=2$ : c'est une base de $F$.}
\piege{oublier de vérifier que $w_1,w_2\in F$. Deux vecteurs libres de $\R^4$ ne forment pas une base de $F$ s'ils n'y sont pas.}

\Partie{C. Une famille de polynômes}

\Q{2.3(a)}
\begin{center}\small\renewcommand{\arraystretch}{1.2}
\begin{tabular}{@{}cccc@{}}\toprule
& $x=0$ & $x=1$ & $x=2$\\\midrule
$P_1=(X-1)(X-2)$ & $2$ & $0$ & $0$\\
$P_2=X(X-2)$ & $0$ & $-1$ & $0$\\
$P_3=X(X-1)$ & $0$ & $0$ & $2$\\\bottomrule
\end{tabular}
\end{center}

\Q{2.3(b)} \choix{les trois polynômes ont le même degré : l'argument « degrés distincts » ne marche pas. Le tableau montre
que chaque $P_i$ ne vit qu'en un point. On évalue la relation en $0$, $1$, $2$ : chaque évaluation isole un coefficient.}

Soient $a,b,c$ tels que $aP_1+bP_2+cP_3=0$. On évalue en $0$ : $2a=0$. En $1$ : $-b=0$. En $2$ : $2c=0$. Donc $a=b=c=0$ : la
famille est libre. Elle a 3 vecteurs et $\dim\R_2[X]=3$ : \textbf{c'est une base de $\R_2[X]$}.
\lienDL{c'est la même idée qu'en III.3 : une relation entre vecteurs devient une information sur des valeurs, qu'on lit point
par point.}

\Q{2.3(c)} D'après (b), $P=aP_1+bP_2+cP_3$ pour des réels $a,b,c$. On évalue en $0$, $1$, $2$ :
$P(0)=2a$, $P(1)=-b$, $P(2)=2c$. Donc
\[P=\frac{P(0)}2\,P_1-P(1)\,P_2+\frac{P(2)}2\,P_3.\]

\Q{2.3(d)} D'après (c), si $P(0)=1$, $P(1)=0$, $P(2)=3$, nécessairement $P=\frac12P_1+\frac32P_3=2X^2-3X+1$. Réciproquement,
ce polynôme vérifie $P(0)=1$, $P(1)=0$, $P(2)=8-6+1=3$. \textbf{L'unique solution est $P=2X^2-3X+1$.}

\Partie{D. Un raisonnement sans coordonnées}

\Q{2.4(a)} \choix{on ne connaît pas les vecteurs, seulement que $(u,v,w)$ est une base. On écrit une combinaison nulle et on la
réécrit en fonction de $u,v,w$ pour utiliser leur liberté.}

Soient $\alpha,\beta,\gamma$ tels que $\alpha(u+v)+\beta(v+w)+\gamma(w+u)=0$, soit $(\alpha+\gamma)u+(\alpha+\beta)v+(\beta+\gamma)w=0$.
Comme $(u,v,w)$ est libre : $\alpha+\gamma=\alpha+\beta=\beta+\gamma=0$. En additionnant, $\alpha+\beta+\gamma=0$, puis
$\beta=\gamma=\alpha=0$. La famille est libre, a 3 vecteurs, et $\dim E=3$ : \textbf{c'est une base de $E$.}
\autre{la matrice de la famille dans la base $(u,v,w)$ est $\begin{pmatrix}1&0&1\\1&1&0\\0&1&1\end{pmatrix}$, de déterminant
$2\neq0$.}

\Q{2.4(b)} $(u-v)+(v-w)+(w-u)=0$ : c'est une relation à coefficients non tous nuls. \textbf{La famille est liée.}

\Q{2.4(c)} $(u+v)-(v+w)+(w+u)=2u$, donc $u=\frac12(u+v)-\frac12(v+w)+\frac12(w+u)$ : \textbf{coordonnées
$\left(\frac12,-\frac12,\frac12\right)$}.

% ---------------------------------------------------------------------
\Ex{3}{Sous-espaces supplémentaires et projecteur}{3,5 points}

\Q{3.1} L'application $\varphi:P\mapsto\big(P(1),P'(1)\big)$, de $E$ dans $\R^2$, est linéaire (l'évaluation et la dérivation le
sont). $F=\Ker\varphi$ est donc un sous-espace vectoriel de $E$.
\autre{vérifier directement : le polynôme nul est dans $F$, et si $P,Q\in F$, alors $(\lambda P+\mu Q)(1)=0$ et
$(\lambda P+\mu Q)'(1)=0$.}

\Q{3.2(a)} Les polynômes $1$, $X-1$, $(X-1)^2$, $(X-1)^3$ sont non nuls et de degrés $0,1,2,3$ distincts : la famille est libre.
Elle a 4 vecteurs et $\dim E=4$ : \textbf{$\mathcal B_1$ est une base de $E$}.

\Q{3.2(b)} \choix{la formule de Taylor donne directement les coordonnées dans $\mathcal B_1$, et les conditions qui
définissent $F$ portent exactement sur les deux premières.}

Par la formule de Taylor, les coordonnées de $P$ dans $\mathcal B_1$ sont $\left(P(1),\,P'(1),\,\frac{P''(1)}2,\,\frac{P'''(1)}6\right)$.
Donc $P\in F$ si et seulement si ses deux premières coordonnées sont nulles, c'est-à-dire si et seulement si
$P\in\vect\big((X-1)^2,(X-1)^3\big)$. La famille $\big((X-1)^2,(X-1)^3\big)$ engendre $F$, et elle est libre car extraite de la
base $\mathcal B_1$. \textbf{C'est une base de $F$, et $\dim F=2$.}

\Q{3.3} \choix{deux méthodes possibles. La plus rapide : intersection nulle et somme des dimensions.}

Soit $P\in F\cap G$. Comme $P\in G=\R_1[X]$, $P=a+b(X-1)$, donc $P(1)=a$ et $P'(1)=b$. Comme $P\in F$, $a=b=0$ : $P=0$. Ainsi
$F\cap G=\{0\}$, et $\dim F+\dim G=2+2=4=\dim E$. \textbf{Donc $E=F\oplus G$.}
\autre{$G=\vect\big(1,X-1\big)$ et $F=\vect\big((X-1)^2,(X-1)^3\big)$ : en réunissant ces deux bases, on obtient la base
$\mathcal B_1$ de $E$, donc $E=F\oplus G$.}

\Q{3.4(a)} D'après la formule de Taylor,
\[P=\underbrace{P(1)+P'(1)(X-1)}_{\in\,G}+\underbrace{\frac{P''(1)}2(X-1)^2+\frac{P'''(1)}6(X-1)^3}_{\in\,F}.\]
La décomposition de $P$ selon $E=F\oplus G$ est unique ; la composante dans $G$ est $q(P)$. \textbf{Donc
$q(P)=P(1)+P'(1)(X-1)$.} Pour $P=X^3$ : $q(X^3)=1+3(X-1)=3X-2$. Contrôle : $X^3-(3X-2)=(X-1)^2(X+2)$ est bien dans $F$.

\Q{3.4(b)} $q$ laisse fixes les vecteurs de $G$ et annule ceux de $F$ :
\[\Mat_{\mathcal B_1}(q)=\begin{pmatrix}1&0&0&0\\0&1&0&0\\0&0&0&0\\0&0&0&0\end{pmatrix}.\]
Dans $\mathcal B_0$ : $q(1)=1$, $q(X)=X$, $q(X^2)=1+2(X-1)=2X-1$, $q(X^3)=3X-2$, donc
\[\Mat_{\mathcal B_0}(q)=\begin{pmatrix}1&0&-1&-2\\0&1&2&3\\0&0&0&0\\0&0&0&0\end{pmatrix}.\]
On vérifie que cette matrice $Q$ vérifie $Q^2=Q$, comme tout projecteur.
\retenir{dans une base adaptée à $E=F\oplus G$, la matrice d'un projecteur est diagonale avec des $1$ et des $0$. Choisir la
bonne base simplifie tout : c'est l'idée de toute la suite du cours.}

% ---------------------------------------------------------------------
\Ex{4}{Calculs de déterminants}{4 points}

\Q{4.1} \choix{la somme de chaque ligne vaut $a+3b$ : on additionne toutes les colonnes dans la première pour factoriser.}
\begin{align*}D_4&=\begin{vmatrix}a&b&b&b\\b&a&b&b\\b&b&a&b\\b&b&b&a\end{vmatrix}
\underset{C_1\leftarrow C_1+C_2+C_3+C_4}{=}(a+3b)\begin{vmatrix}1&b&b&b\\1&a&b&b\\1&b&a&b\\1&b&b&a\end{vmatrix}\\
&\underset{L_i\leftarrow L_i-L_1}{=}(a+3b)\begin{vmatrix}1&b&b&b\\0&a-b&0&0\\0&0&a-b&0\\0&0&0&a-b\end{vmatrix}.\end{align*}
La dernière matrice est triangulaire : \textbf{$D_4=(a+3b)(a-b)^3$}. La matrice est inversible si et seulement si $D_4\neq0$,
c'est-à-dire \textbf{$a\neq b$ et $a\neq-3b$}.
\piege{développer directement un $4\times4$ : 24 termes et presque sûrement une erreur. Toujours chercher d'abord une
structure (somme constante, colonnes proches, zéros).}

\Q{4.2(a)} On fait $L_2\leftarrow L_2-L_1$ et $L_3\leftarrow L_3-L_1$, puis on factorise :
\[V=\begin{vmatrix}1&a&a^2\\0&b-a&b^2-a^2\\0&c-a&c^2-a^2\end{vmatrix}=(b-a)(c-a)\begin{vmatrix}1&b+a\\1&c+a\end{vmatrix}=(b-a)(c-a)(c-b).\]
\gain{$b^2-a^2=(b-a)(b+a)$ : on fait apparaître les facteurs $(b-a)$ et $(c-a)$ avant de développer.}

\Q{4.2(b)} $P=x_0+x_1X+x_2X^2$ vérifie les trois conditions si et seulement si
\[\begin{pmatrix}1&a&a^2\\1&b&b^2\\1&c&c^2\end{pmatrix}\begin{pmatrix}x_0\\x_1\\x_2\end{pmatrix}=\begin{pmatrix}\alpha\\\beta\\\gamma\end{pmatrix}.\]
Le déterminant de ce système est $V\neq0$ car $a,b,c$ sont distincts : c'est un système de Cramer, il a une unique solution.
\textbf{Il existe un unique $P\in\R_2[X]$ qui convient.}
\lienDL{c'est le même phénomène qu'en 2.3, où $a=0$, $b=1$, $c=2$, et qu'en III.3 du DL : trois valeurs distinctes
« séparent » les polynômes de degré au plus 2. Ici on le prouve par le déterminant, en 2.3 par évaluation.}

\Q{4.3} \choix{on ne peut rien calculer : on utilise les propriétés, $\det A^T=\det A$ et $\det(\lambda A)=\lambda^n\det A$.}

$A^T=-A$, donc $\det A=\det A^T=\det(-A)=(-1)^n\det A=-\det A$ car $n$ est impair. Ainsi $2\det A=0$ : \textbf{$\det A=0$, et
$A$ n'est pas inversible.} Pour $n$ pair, c'est faux : $\begin{pmatrix}0&1\\-1&0\end{pmatrix}$ est antisymétrique, de
déterminant $1$.
\piege{écrire $\det(-A)=-\det A$ sans justifier : c'est vrai seulement parce que $n$ est impair.}

% ---------------------------------------------------------------------
\Pb{une matrice nilpotente}{4,5 points}

\Q{P.1}
\[A^2=\begin{pmatrix}0&0&0\\1&1&1\\-1&-1&-1\end{pmatrix},\qquad A^3=A^2A=0.\]
(Chaque ligne de $A^2$ est proportionnelle à $(1,1,1)$, et $(1,1,1)A=(0,0,0)$.)

\Q{P.2} \emph{Noyau.} $AX=0$ s'écrit $-x-y-z=0$, $x+2y+2z=0$, $-y-z=0$. La dernière équation donne $z=-y$, puis la première
$x=0$ ; la deuxième est alors vérifiée. \textbf{$\Ker A=\vect\big((0,1,-1)\big)$}, de dimension 1.

\emph{Image.} Elle est engendrée par les colonnes $C_1=(-1,1,0)$, $C_2=(-1,2,-1)$, $C_3=C_2$. $C_1$ et $C_2$ ne sont pas
colinéaires : \textbf{$\Im A=\vect(C_1,C_2)$, et $\rg A=2$}. Théorème du rang : $1+2=3$. \checkmark

\emph{Équation.} $C_1$ et $C_2$ vérifient $x+y+z=0$, et ce plan est de dimension 2 : \textbf{$\Im A$ a pour équation $x+y+z=0$}.

\emph{Inclusion.} $(0,1,-1)$ vérifie $0+1-1=0$, donc $\Ker A\subset\Im A$.

\Q{P.3} \choix{l'énoncé interdit le déterminant : on utilise $A^3=0$. Appliquer $A^2$ à une relation fait disparaître tous
les termes sauf un.}

On a $x_0=e_1$, $Ax_0=(-1,1,0)$, $A^2x_0=(0,1,-1)\neq0$. Soient $c_0,c_1,c_2$ tels que $c_0x_0+c_1Ax_0+c_2A^2x_0=0$.
\begin{itemize}[leftmargin=1.2em,itemsep=1pt]
\item On multiplie par $A^2$ : comme $A^3=A^4=0$, il reste $c_0A^2x_0=0$, et $A^2x_0\neq0$, donc $c_0=0$.
\item On multiplie alors par $A$ : il reste $c_1A^2x_0=0$, donc $c_1=0$.
\item Il reste $c_2A^2x_0=0$, donc $c_2=0$.
\end{itemize}
La famille est libre, de 3 vecteurs en dimension 3 : \textbf{$\mathcal B$ est une base de $\R^3$.}
\lienDL{c'est exactement la question II.1 du DL, avec une autre matrice.}

\Q{P.4} $f(x_0)=Ax_0$ est le 2\textsuperscript{e} vecteur de $\mathcal B$, $f(Ax_0)=A^2x_0$ le 3\textsuperscript{e}, et
$f(A^2x_0)=A^3x_0=0$ :
\[N=\Mat_{\mathcal B}(f)=\begin{pmatrix}0&0&0\\1&0&0\\0&1&0\end{pmatrix},\qquad
P=\begin{pmatrix}1&-1&0\\0&1&1\\0&0&-1\end{pmatrix},\qquad A=PNP^{-1}.\]
$P$ a pour colonnes les coordonnées de $x_0$, $Ax_0$, $A^2x_0$ dans la base canonique. La relation vient de la formule de
changement de base $N=P^{-1}AP$.

\Q{P.5} \choix{« tous les $x$ » : deux sens, et la question P.3 fait le premier.}

Soit $x=(x_1,x_2,x_3)$. On a $A^2x=(x_1+x_2+x_3)(0,1,-1)$.
\begin{itemize}[leftmargin=1.2em,itemsep=1pt]
\item Si $x_1+x_2+x_3\neq0$, alors $A^2x\neq0$ et le raisonnement de P.3 s'applique mot pour mot : la famille est une base.
\item Si $x_1+x_2+x_3=0$, alors $A^2x=0$ : la famille contient le vecteur nul, elle est liée.
\end{itemize}
\textbf{$(x,Ax,A^2x)$ est une base si et seulement si $x_1+x_2+x_3\neq0$}, c'est-à-dire $x\notin\Ker A^2$.
\lienDL{c'est la question II.4.a du DL. La condition ne se devine pas : elle vient du calcul de $A^2x$.}

\Q{P.6} \choix{calculer directement $\det(A+\lambda I_3)$ marche, mais la relation $A=PNP^{-1}$ donne le résultat sans calcul :
deux matrices semblables ont le même déterminant.}

$A+\lambda I_3=PNP^{-1}+\lambda PP^{-1}=P(N+\lambda I_3)P^{-1}$, donc
\[\det(A+\lambda I_3)=\det P\,\det(N+\lambda I_3)\,\det P^{-1}=\det(N+\lambda I_3)=\lambda^3,\]
car $N+\lambda I_3$ est triangulaire, de diagonale $(\lambda,\lambda,\lambda)$. Avec $\lambda=1$ : $\det(A+I_3)=1\neq0$,
\textbf{$A+I_3$ est inversible}. Avec $\lambda=0$ : $\det A=0$, \textbf{$A$ n'est pas inversible} (on le savait : $\rg A=2$).
\gain{on peut même donner l'inverse : $(I_3+A)(I_3-A+A^2)=I_3+A^3=I_3$.}

% ---------------------------------------------------------------------
\clearpage
\markboth{Diagnostic}{}
\subsection*{Après la correction : quoi revoir ?}
\begin{center}
\small\renewcommand{\arraystretch}{1.25}
\begin{tabular}{@{}>{\raggedright\arraybackslash}p{5.6cm}>{\raggedright\arraybackslash}p{7.4cm}@{}}
\toprule
\textbf{Si tu as perdu des points en…} & \textbf{Revoir}\\\midrule
1.2, 2.3(b), 2.4(a), P.3 (liberté) & Livret d'algèbre linéaire, méthode R.2 ; livret du DL, outils 1 et 2\\
2.2, 3.2 (génératrice, base d'un sous-espace) & Livret d'algèbre linéaire, rappels « Espaces vectoriels » et « Dimension finie »\\
2.1(b), 2.4(c), 3.4, P.4 (coordonnées, matrices) & Livret d'algèbre linéaire, proposition « changement de base » ; livret du DL, outil 3\\
3.3, 3.4 (supplémentaires, projecteur) & Livret d'algèbre linéaire, §1.2\\
P.2 (noyau, image, rang) & Livret d'algèbre linéaire, méthode R.3 « Calculer un rang, un noyau, une image »\\
4.1, 4.2 (calculs) & Livret des déterminants, §1.2 et §1.5 (somme constante, Vandermonde)\\
1.5, 4.3, P.6 (propriétés) & Livret des déterminants, §1.3\\
2.1(a), P.5 (« tous les… ») & Livret du DL, outil 4\\\bottomrule
\end{tabular}
\end{center}

\begin{essentiel}
\begin{itemize}[leftmargin=1.1em,itemsep=3pt]
\item Libre : combinaison nulle, puis le bon geste pour isoler les coefficients (identifier, évaluer, composer par $A^k$).
\item $n$ vecteurs en dimension $n$ : une seule propriété à prouver. Pour une base d'un sous-espace $F$, vérifier aussi que les
vecteurs sont dans $F$.
\item Un paramètre et $n$ vecteurs en dimension $n$ : le déterminant donne toutes les valeurs d'un coup.
\item Un déterminant se calcule par opérations sur les lignes et colonnes avant de développer.
\item $\det(\lambda A)=\lambda^n\det A$, $\det A^T=\det A$, deux matrices semblables ont le même déterminant.
\end{itemize}
\end{essentiel}

\else
% =====================================================================
%                               SUJET
% =====================================================================
\markboth{Examen blanc · Algèbre linéaire et déterminants}{}
\begin{center}
\begin{tcolorbox}[enhanced,colback=Fond,colframe=Fond,arc=0pt,boxrule=0pt,left=10pt,right=10pt,top=8pt,bottom=8pt,width=\linewidth]
\small
\textbf{\sffamily\color{Bleu}Consignes.} Le sujet comporte quatre exercices et un problème, indépendants : on peut les traiter
dans l'ordre de son choix. Toute réponse doit être \emph{justifiée} : citer le résultat du cours utilisé et vérifier ses
hypothèses. Les résultats sont à encadrer. Une réponse juste sans justification ne rapporte pas de points ; une démarche correcte
avec une erreur de calcul en rapporte. La clarté et la précision de la rédaction sont prises en compte.

\medskip
\begin{tabular}{@{}lcccccc@{}}
\textbf{\sffamily\color{Bleu}Barème} & Ex.~1 : 3 & Ex.~2 : 5 & Ex.~3 : 3,5 & Ex.~4 : 4 & Problème : 4,5 & \textbf{Total : 20}\\
\end{tabular}
\end{tcolorbox}
\end{center}

\Ex{1}{Vrai ou faux}{3 points}
Pour chaque affirmation, dire si elle est vraie ou fausse, et le justifier : une preuve courte si elle est vraie, un
contre-exemple si elle est fausse. \pts{0,5 point par affirmation}
\begin{enumerate}
\item Toute famille de quatre vecteurs de $\R^3$ est liée.
\item Si $(u,v)$ et $(v,w)$ sont deux familles libres d'un espace vectoriel $E$, alors $(u,v,w)$ est libre.
\item Une famille de trois vecteurs qui engendre $\R^3$ est une base de $\R^3$.
\item La réunion de deux sous-espaces vectoriels de $\R^2$ est un sous-espace vectoriel de $\R^2$.
\item Pour toute matrice $A\in\mathcal M_3(\R)$, $\det(2A)=2\det A$.
\item Une application linéaire de $\R^4$ dans $\R^3$ n'est jamais injective.
\end{enumerate}

\Ex{2}{Familles libres, génératrices, bases}{5 points}

\Partie{A. Une famille de $\R^3$ avec paramètre}
Pour $m\in\R$, on pose $u_1=(1,2,-1)$, $u_2=(0,1,1)$ et $u_3=(2,m,1)$.
\begin{enumerate}
\item
\begin{enumerate}
\item Déterminer les valeurs de $m$ pour lesquelles $(u_1,u_2,u_3)$ est une base de $\R^3$. \pts{0,75 pt}
\item On prend $m=1$ et on note $\mathcal B'=(u_1,u_2,u_3)$. Déterminer les coordonnées du vecteur $(1,1,1)$ dans $\mathcal B'$.
\pts{0,5 pt}
\item Pour la valeur de $m$ exclue au (a), écrire $u_3$ comme combinaison linéaire de $u_1$ et $u_2$, et donner le rang de
$(u_1,u_2,u_3)$. \pts{0,25 pt}
\end{enumerate}
\end{enumerate}

\Partie{B. Un sous-espace de $\R^4$}
Soit $F=\{(x,y,z,t)\in\R^4,\ x+y+z+t=0\ \text{et}\ x-y+2z=0\}$.
\begin{enumerate}[resume]
\item
\begin{enumerate}
\item Déterminer une famille génératrice de $F$, puis montrer que c'est une base de $F$. En déduire $\dim F$. \pts{0,75 pt}
\item Montrer que $w_1=(3,1,-1,-3)$ et $w_2=(0,2,1,-3)$ forment une base de $F$. \pts{0,5 pt}
\end{enumerate}
\end{enumerate}

\Partie{C. Une famille de polynômes}
Dans $\R_2[X]$, on pose $P_1=(X-1)(X-2)$, $P_2=X(X-2)$ et $P_3=X(X-1)$.
\begin{enumerate}[resume]
\item
\begin{enumerate}
\item Calculer $P_i(0)$, $P_i(1)$ et $P_i(2)$ pour $i=1,2,3$. \pts{0,25 pt}
\item Montrer que $(P_1,P_2,P_3)$ est une base de $\R_2[X]$. \pts{0,5 pt}
\item Soit $P\in\R_2[X]$. Exprimer les coordonnées de $P$ dans cette base à l'aide de $P(0)$, $P(1)$ et $P(2)$. \pts{0,25 pt}
\item En déduire l'unique polynôme $P\in\R_2[X]$ tel que $P(0)=1$, $P(1)=0$ et $P(2)=3$. \pts{0,25 pt}
\end{enumerate}
\end{enumerate}

\Partie{D. Un raisonnement sans coordonnées}
Soit $E$ un espace vectoriel de dimension 3 et $(u,v,w)$ une base de $E$.
\begin{enumerate}[resume]
\item
\begin{enumerate}
\item Montrer que $(u+v,\ v+w,\ w+u)$ est une base de $E$. \pts{0,5 pt}
\item La famille $(u-v,\ v-w,\ w-u)$ est-elle une base de $E$ ? \pts{0,25 pt}
\item Donner les coordonnées de $u$ dans la base $(u+v,\ v+w,\ w+u)$. \pts{0,25 pt}
\end{enumerate}
\end{enumerate}

\Ex{3}{Sous-espaces supplémentaires et projecteur}{3,5 points}
Soit $E=\R_3[X]$, de base canonique $\mathcal B_0=(1,X,X^2,X^3)$. On pose
\[F=\{P\in E,\ P(1)=0\ \text{et}\ P'(1)=0\}\qquad\text{et}\qquad G=\R_1[X].\]
On rappelle la formule de Taylor : pour tout $P\in E$,
$P=P(1)+P'(1)(X-1)+\frac{P''(1)}{2}(X-1)^2+\frac{P'''(1)}{6}(X-1)^3$.
\begin{enumerate}
\item Montrer que $F$ est un sous-espace vectoriel de $E$. \pts{0,5 pt}
\item
\begin{enumerate}
\item Montrer que $\mathcal B_1=\big(1,\ X-1,\ (X-1)^2,\ (X-1)^3\big)$ est une base de $E$. \pts{0,5 pt}
\item Donner les coordonnées d'un polynôme $P$ dans $\mathcal B_1$. En déduire une base de $F$ et sa dimension. \pts{0,75 pt}
\end{enumerate}
\item Montrer que $E=F\oplus G$. \pts{0,75 pt}
\item On note $q$ le projecteur sur $G$ parallèlement à $F$.
\begin{enumerate}
\item Montrer que $q(P)=P(1)+P'(1)(X-1)$ pour tout $P\in E$. Calculer $q(X^3)$. \pts{0,5 pt}
\item Écrire la matrice de $q$ dans la base $\mathcal B_1$, puis dans la base $\mathcal B_0$. \pts{0,5 pt}
\end{enumerate}
\end{enumerate}

\Ex{4}{Calculs de déterminants}{4 points}
\begin{enumerate}
\item Soient $a,b\in\R$. Calculer, sous forme factorisée, \pts{1,5 pt}
\[D_4=\begin{vmatrix}a&b&b&b\\b&a&b&b\\b&b&a&b\\b&b&b&a\end{vmatrix}.\]
Pour quelles valeurs de $a$ et $b$ la matrice correspondante est-elle inversible ?
\item Soient $a,b,c$ trois réels.
\begin{enumerate}
\item Calculer, sous forme factorisée, $V=\begin{vmatrix}1&a&a^2\\1&b&b^2\\1&c&c^2\end{vmatrix}$. \pts{1 pt}
\item On suppose $a,b,c$ deux à deux distincts. Soient $\alpha,\beta,\gamma\in\R$. Montrer qu'il existe un unique polynôme
$P\in\R_2[X]$ tel que $P(a)=\alpha$, $P(b)=\beta$ et $P(c)=\gamma$. \pts{0,5 pt}
\end{enumerate}
\item Soit $n$ un entier impair et $A\in\mathcal M_n(\R)$ une matrice antisymétrique ($A^T=-A$). Montrer que $A$ n'est pas
inversible. Le résultat reste-t-il vrai si $n$ est pair ? \pts{1 pt}
\end{enumerate}

\Pb{une matrice nilpotente}{4,5 points}
On considère
\[A=\begin{pmatrix}-1&-1&-1\\1&2&2\\0&-1&-1\end{pmatrix}\]
et $f$ l'endomorphisme de $\R^3$ canoniquement associé à $A$. On note $(e_1,e_2,e_3)$ la base canonique de $\R^3$.
\begin{enumerate}
\item Calculer $A^2$ et $A^3$. \pts{0,5 pt}
\item Déterminer une base de $\Ker A$ et une base de $\Im A$. Vérifier le théorème du rang, donner une équation de $\Im A$, et
montrer que $\Ker A\subset\Im A$. \pts{1 pt}
\item On pose $x_0=e_1$. Montrer, \emph{sans calculer de déterminant}, que $\mathcal B=(x_0,Ax_0,A^2x_0)$ est une base de
$\R^3$. \pts{1 pt}
\item Écrire la matrice $N$ de $f$ dans la base $\mathcal B$, la matrice de passage $P$ de la base canonique à $\mathcal B$, et
la relation entre $A$, $N$ et $P$. \pts{0,5 pt}
\item Déterminer tous les vecteurs $x=(x_1,x_2,x_3)$ de $\R^3$ tels que $(x,Ax,A^2x)$ soit une base de $\R^3$. \pts{1 pt}
\item Pour $\lambda\in\R$, calculer $\det(A+\lambda I_3)$ en utilisant la question 4. En déduire que $A+I_3$ est inversible et
que $A$ ne l'est pas. \pts{0,5 pt}
\end{enumerate}

\vfill
\begin{center}\sffamily\small\color{Gris}— Fin du sujet —\end{center}
\fi
\end{document}
